\section{User Input}
\label{SEC:User_Input}

As described in \sref{SEC:Files:Namelist_Input},
GFR uses a Fortran namelist file for specifying user inputs,
and can output a default input file for you to modify.
%
The name of the input file to read should be added to the end of the MPI
command line used to execute the code.
%
This input file contains a number of namelist input blocks,
which are described in the sections that follow.
%
\begin{multicols}{3}
\raggedcolumns % Permit last column to be short.
\begin{itemize}
  \item \nameref{SEC:nml_version}
  \item \nameref{SEC:nml_grid}
  \item \nameref{SEC:nml_reference}
  \item \nameref{SEC:nml_physics}
  \item \nameref{SEC:nml_initialization}
  \item \nameref{SEC:nml_timestep}
  \item \nameref{SEC:nml_bc}
  \item \nameref{SEC:nml_numerics}
  \item \nameref{SEC:nml_limiter}
  \item \nameref{SEC:nml_time_ave}
  \item \nameref{SEC:nml_output}
  \item \nameref{SEC:nml_output_variables}
  \item \nameref{SEC:nml_sem}
  \item \nameref{SEC:nml_chemistry}
  \item \nameref{SEC:nml_testcase_channel}
  \item \nameref{SEC:nml_testcase_pulse}
  \item \nameref{SEC:nml_testcase_sod}
  \item \nameref{SEC:nml_testcase_tcf}
  \item \nameref{SEC:nml_testcase_vortex}
  \item \nameref{SEC:nml_interpolation}
  \item \nameref{SEC:nml_utility}
  \item \nameref{SEC:nml_debug}
\end{itemize}
\end{multicols}

The order of the namelist blocks in the input file does not matter,
nor does the order of the variables within a namelist block.
%
The user may include every namelist block and every input variable
in the input file.
%
However, the only \emph{required} entries are those inputs that the user wants
to have non-default values.
%
It is highliy recommended to include the \nameref{SEC:nml_version}
namelist block so that future versions of the code can properly
interpret and convert your input file.
%
%
Note that a ``!'' character in a Fortran namelist file begins a comment that
extends to the end of the line.
%
This can be useful for documenting your input file.

%%%%%%%%%%%%%%%%%%%%%%%%%%%%%%%%%%%%%%%%%%%%%%%%%%%%%%%%%%%%%%%%%%%%%%%%%%%%%%%%

In each of the namelist input blocks described in the following sections,
we first identify all the variables in the namelist block and then provide
a detailed explanation of each input.
%
Comments in the listings are used to document the expected units of the input
variables (a good practice), and the values shown represent the default
settings.

%%%%%%%%%%%%%%%%%%%%%%%%%%%%%%%%%%%%%%%%%%%%%%%%%%%%%%%%%%%%%%%%%%%%%%%%%%%%%%%%

\AddToHook{cmd/subsection/before}{\clearpage}
\subsection{nml\_version}
\label{SEC:nml_version}
\inp*[nmlVersion]{}[nml\_version][]

This \hyperref[SEC:User_Input]{namelist block}
is used to inform GFR which version of the namelist to expect when reading
through the input file.
%
This namelist block is provided as a means to allow the use of older input files
without having to manually convert them to the current format.

\begin{nmlblocklst}[
,label={LST:nml_version}
,caption={Namelist nml\_version.}
]
&nml_version
  Namelist_Version = 2
/
\end{nmlblocklst}

%==========================================================================

\begin{nmldescitem}
Namelist_Version = 2
\end{nmldescitem}
\inp*[nmlVersion::NamelistVersion]{}[Namelist\_Version][]

\begin{nmldescbody}
The format of the namelist input file.
Valid options are: 1, 2.
%
Input variables may be added, removed, or renamed between versions.
However, the code will attempt to convert an older input file to the
current format.
%
This report documents the format of namelist version 2. Users should add the
\inptext{nml\_version} block shown in \lref{LST:nml_version} to the top of
their input files so that future versions of the code will recognize it.
\end{nmldescbody}

%==========================================================================


\subsection{nml\_grid}
\label{SEC:nml_grid}
\inp*[nmlGrid]{}[nml\_grid][]

This \hyperref[SEC:User_Input]{namelist block}
is used to specify parameters related to the grid input file.

\begin{nmlblocklst}[
,label={LST:nml_grid}
,caption={Namelist nml\_grid.}
]
&nml_grid
  Grid_Scaling_Factor      =   1.0
  Grid_Format              =  -1
  Grid_File                =  ""
  Plot3d_Cutfile           =  ""
  Plot3d_BCSfile           =  ""
  Plot3d_Stride_I          =  -1
  Plot3d_Stride_J          =  -1
  Plot3d_Stride_K          =  -1
  Plot3d_Stride_All        =  -1
  Plot3d_Agglomerate_Order =  -1
  Correct_CGNS_Grid_Points =   F
  Internal_Grid_DoF        =   0
  Internal_Grid_Order      =  -1
  Metis_Option             =   3
  Geom_Wts(1)%Tria         =   1
  Geom_Wts(1)%Quad         =   1
  Geom_Wts(1)%Tetr         =   1
  Geom_Wts(1)%Pyra         =   1
  Geom_Wts(1)%Pris         =   1
  Geom_Wts(1)%Hexa         =   1
/
\end{nmlblocklst}
%%Geom_Wts(1)%Edge         =   1
%%Geom_Wts(1)%Face_Tria    =   1
%%Geom_Wts(1)%Face_Quad    =   1

%==========================================================================

\begin{nmldescitem}
Grid_Scaling_Factor = 1.0
\end{nmldescitem}
\inp*[nmlGrid::GridScalingFactor]{}[Grid\_Scaling\_Factor][]

\begin{nmldescbody}
Multiplication factor ($> 0.0$) to scale the input grid file coordinates to
the desired length units in meters.
%
For example, use a value of 0.0254 to convert a grid containing coordinates
in inches to meters.
%
See the discussion of
  \key[key:GridUnits]{grid units}[grid units]
and
  \key[key:GridScaledUnits]{grid scaled units}[grid scaled units]
in \sref{SEC:Grid_Scaling}.
\end{nmldescbody}

%==========================================================================

\begin{nmldescitem}
Grid_Format = -1
\end{nmldescitem}
\inp*[nmlGrid::GridFormat]{}[Grid\_Format][]

\begin{nmldescbody}
Specifies the format of the input grid file.

\begin{tabular}{ @{} r @{ = } p{0.85\linewidth} @{} }
  -1 & Internally generated grid
       (only supported with certain 
       \inp[nmlInitialization::Initialization]{Initialization} options) \\
   1 & Gambit Neutral (very limited support, high-order curved elements not
       supported) \\
   2 & Gmsh (limited support) \\
   3 & Plot3D \\
   4 & CGNS (recommended)
\end{tabular}

\end{nmldescbody}

%==========================================================================

%##%\FINALFIX{\newpage}
\begin{nmldescitem}
Grid_File = ""
\end{nmldescitem}
\inp*[nmlGrid::GridFile]{}[Grid\_File][]

\begin{nmldescbody}
File path to the input grid file relative to the current directory.
\end{nmldescbody}

%==========================================================================

\begin{nmldescitem}
Plot3d_Cutfile = ""
\end{nmldescitem}
\inp*[nmlGrid::Plot3dCutfile]{}[Plot3d\_Cutfile][]

\begin{nmldescbody}
File path to the input file that specifies the cut conditions for a
multiblock Plot3D grid file.
\end{nmldescbody}

%==========================================================================

\begin{nmldescitem}
Plot3d_BCSfile = ""
\end{nmldescitem}
\inp*[nmlGrid::Plot3dBCSfile]{}[Plot3d\_BCSfile][]

\begin{nmldescbody}
File path to the input file that specifies the boundary conditions for a
Plot3D grid file.
\end{nmldescbody}

%==========================================================================

\begin{nmldescitem}
Plot3d_Stride_All = -1
Plot3d_Stride_I   = -1 
Plot3d_Stride_J   = -1
Plot3d_Stride_K   = -1
\end{nmldescitem}
\inp*[nmlGrid::Plot3dStrideAll]{}[Plot3d\_Stride\_All][]%
\inp*[nmlGrid::Plot3dStrideI]{}[Plot3d\_Stride\_I][]%
\inp*[nmlGrid::Plot3dStrideJ]{}[Plot3d\_Stride\_J][]%
\inp*[nmlGrid::Plot3dStrideK]{}[Plot3d\_Stride\_K][]%
 
\begin{nmldescbody}
If $> 0$, applies a coarsening stride to the specified direction
of all Plot3D blocks.
 
\inptext{Plot3d\_Stride\_All} is applied first to all directions.
%
Then the individual i, j, k options can be used to override the
stride setting in that specific direction.
\end{nmldescbody}

%==========================================================================

\begin{nmldescitem}
Plot3d_Agglomerate_Order = -1
\end{nmldescitem}
\inp*[nmlGrid::Plot3dAgglomerateOrder]{}[Plot3d\_Agglomerate\_Order][]

\begin{nmldescbody}
If $> 0$, agglomerates linear Plot3D cells into high-order grid cells of
this polynomial order.
%
For this to be successfull, the number of {\em cells} in each direction
for every block must be divisible by this value.
\end{nmldescbody}

%==========================================================================

\begin{nmldescitem}
Correct_CGNS_Grid_Points = F
\end{nmldescitem}
\inp*[nmlGrid::CorrectCGNSGridPoints]{}[Correct\_CGNS\_Grid\_Points][]

\begin{nmldescbody}
Experimental code that tries to correct the high-order points given in a
CGNS grid file created by Pointwise. The high-order points should be equally
spaced, but some software such as Pointwise does not require this.
\end{nmldescbody}

%==========================================================================

\begin{nmldescitem}
Internal_Grid_DoF = 0
\end{nmldescitem}
\inp*[nmlGrid::InternalGridDoF]{}[Internal\_Grid\_DoF][]
\key*{}[degree-of-freedom (DoF)!internally generated grid]
\key*{}[internally generated grid!degree-of-freedom (DoF)]

\begin{nmldescbody}
This specifies the grid resolution to use for an
\key[key:InternallyGeneratedGrid]{internally generated grid}[=].
\\
\begin{tabular}{ @{\quad} r @{ \quad } p{0.85\linewidth} @{} }
  $<0$ & User specified number of cells. \\
  $=0$ & 512 degrees of freedom. \\
  $>0$ & User specified degrees of freedom.
\end{tabular}\\
%
If the number of degrees of freedom are specified, then the number of cells
needed to achieve the specified degrees of freedom in each direction will be
computed from the solution order as:
\begin{equation*}
    \left[ \text{\# of cells} \right]
  = \mathrm{nint}
    ( \inptext{Internal\_Grid\_DoF}
%   ( \text{\inp[nmlGrid::InternalGridDoF]{Internal\_Grid\_DoF}}
    / ( \text{\inp[nmlNumerics::SolutionOrder]{Solution\_Order}} + 1 )
    ).
\end{equation*}

\NOTE
An internally generated grid will only be created for specific
\inp[nmlInitialization::Initialization]{Initialization} conditions,
as indicated on \pref{\pageref{nmlInitialization::Initialization}}.
%
For the following cases, one should instead use the grid inputs in the
associated namelist block.

\begin{tabular}{S ll}
  {\inp[nmlInitialization::Initialization]{Initialization}}
& Description
& Namelist Block \\
  \hline
   2 &  Channel flow
     &  \inp[nmlTestcaseChannel]{nml\_testcase\_channel} \\
  21 &  Taylor-Couette flow
     &  \inp[nmlTestcaseTCF]{nml\_testcase\_tcf} \\
  22 &  Freestream Taylor-Couette flow
     &  \inp[nmlTestcaseTCF]{nml\_testcase\_tcf}
\end{tabular}
\end{nmldescbody}

%==========================================================================

\begin{nmldescitem}
Internal_Grid_Order = -1
\end{nmldescitem}
\inp*[nmlGrid::InternalGridOrder]{}[Internal\_Grid\_Order][]

\begin{nmldescbody}
Polynomial order used by the internal grid generator to create a
high-order grid.
%
This only applies to internally generated grids with curved geometries
(\thatis Taylor-Couette flow, channel, and annular pipe channel test cases).
%
All other internally generated grids have linear geometries.
\end{nmldescbody}

%==========================================================================

\FINALFIX{\newpage}
\begin{nmldescitem}
Metis_Option = 3
\end{nmldescitem}
\inp*[nmlGrid::MetisOption]{}[Metis\_Option][]

\begin{nmldescbody}
Specifies the \key{METIS}[=] algorithm to use for
partitioning the global grid. If the specified option fails, it will
loop through all of the other options to try to find a successful partitioning.

%% Note: longtable permits page breaks.
%\begin{tabular}{ @{} r @{ = } p{0.85\linewidth} @{} }
\begin{longtable}{ @{} r @{ = } p{0.85\linewidth} @{} }
  1 & Create non-weighted partitions based on a multilevel recursive bisection
      while minimizing the edge cut.
      \\
  2 & Create non-weighted partitions based on a multilevel $k$-way algorithm
      while minimizing the edge cut.
      \\
  3 & Create non-weighted partitions based on a multilevel $k$-way algorithm
      while minimizing the total communication volume.
      \\
  4 & Create weighted partitions based on a multilevel $k$-way algorithm
      while minimizing the total communication volume.
      This is equivalent to option 3 but also
      includes cell weights (based on geometry type).
      \\
  5 & Create weighted partitions based on a multilevel $k$-way algorithm
      while minimizing the edge cut.
      This is equivalent to option 2 but also
      includes cell weights (based on geometry type).
      \\
  6 & Create weighted partitions based on a multilevel recursive bisection
      while minimizing the edge cut.
      This is equivalent to option 1 but also
      includes cell weights (based on geometry type).
\end{longtable}
%\end{tabular}

Note that the last three options include geometry-dependent
computational cost into load balancing the partitions.
%
Weighting for each geometry type and order is specified through the input
\inp[nmlGrid::GeomWts]{Geom\_Wts}[=][].

\end{nmldescbody}

%==========================================================================

\begin{nmldescitem}
Geom_Wts(1)%Tria = 1
Geom_Wts(1)%Quad = 1
Geom_Wts(1)%Tetr = 1
Geom_Wts(1)%Pyra = 1
Geom_Wts(1)%Pris = 1
Geom_Wts(1)%Hexa = 1
\end{nmldescitem}
%%If including these, then align = signs above.
%Geom_Wts(1)%Edge      = 1
%Geom_Wts(1)%Face_Tria = 1
%Geom_Wts(1)%Face_Quad = 1
\inp*[nmlGrid::GeomWts]{}[Geom\_Wts][]
%\inp*[nmlGrid::GeomWts/Edge]{}[Geom\_Wts(:)\%Edge][]
\inp*[nmlGrid::GeomWts/Tria]{}[Geom\_Wts(:)\%Tria][]
\inp*[nmlGrid::GeomWts/Quad]{}[Geom\_Wts(:)\%Quad][]
\inp*[nmlGrid::GeomWts/Tetr]{}[Geom\_Wts(:)\%Tetr][]
\inp*[nmlGrid::GeomWts/Pyra]{}[Geom\_Wts(:)\%Pyra][]
\inp*[nmlGrid::GeomWts/Pris]{}[Geom\_Wts(:)\%Pris][]
\inp*[nmlGrid::GeomWts/Hexa]{}[Geom\_Wts(:)\%Hexa][]
%\inp*[nmlGrid::GeomWts/FaceTria]{}[Geom\_Wts(:)\%Face\_Tria][]
%\inp*[nmlGrid::GeomWts/FaceQuad]{}[Geom\_Wts(:)\%Face\_Quad][]

\begin{nmldescbody}
These are the geometry weights used by the \key{METIS}[=]
partitioning algorithm to split the mesh between processors.
%
Valid values are integers in the range [0:100], with larger values
indicating those geometries must perform more computational work
relative to geometries with lower values.
%
Note that the \inpstyle{Tria} and \inpstyle{Quad} inputs are for
element volumes for 2D simulations, and thus are not used for 3D simulations.
%
%Note that the \inpstyle{Tria} and \inpstyle{Quad} inputs are for
%element volumes in two dimensions, whereas the \inpstyle{Face\_Tria}
%and \inpstyle{Face\_Quad} inputs are for element surfaces in
%three dimensions.
%%
%Currently, the \inpstyle{Face\_Tria} and \inpstyle{Face\_Quad} weights are not
%used.
%
See the \inp[nmlGrid::MetisOption]{Metis\_Option} input
for additional details.

Preliminary testing of the geometry weights indicates that improvements
in performance may be highly dependent upon the number and distribution
of each cell type throughout the domain and how many partitions are
requested.
%
We therefore offer no specific weighting recommendations.
 
\IDEA{
  Add table of weights for either Rome or Milan nodes since AMD nodes are used
  at both NAS and a lot of DoD machines, just to serve as a rough starting
  point. Some things to probably point out:
  \begin{itemize}
    \item
      Weights are only necessary for unstructured grids.
    \item
      Not using weights can severely degrade performance for large grids using
      large numbers of MPI processes (10k+).
    \item
      Unless the user has compiled a custom METIS library using 64-bit
      integers, the [maximum weight value]*[number of grid cells] needs to be
      smaller than $2^{31}$ or the partitioning will fail. Use one of the
      example Rome\slash Milan weights from the table to illustrate this point
      with an example of too many cells for these weight values, and how to
      scale the weights to get them within the 32-bit integer limit.
    \item
      Possibly provide an example of how unpredictable the weights can be to
      make the user aware of this issue.
    \item
      Include some remarks about using weights with the different partitioning
      methods and how unpredictable the results can be. Maybe include an example
      of this with the example of how unpredictable the weights can be.
  \end{itemize}
}
\end{nmldescbody}

%==========================================================================


%##%\FINALFIX{\newpage}
\subsection{nml\_reference}
\label{SEC:nml_reference}
\inp*[nmlReference]{}[nml\_reference][]

This \hyperref[SEC:User_Input]{namelist block}
is used to specify the reference conditions for the simulation.
%
These conditions are used within the flow solver to nondimensionalize
the internal variables, and in most cases to initialize the
volume solution.
%
They are also written to the screen during code execution.

\begin{nmlblocklst}[
,label={LST:nml_reference}
,caption={Namelist nml\_reference.}
]
&nml_reference
  Mach         = 0.4
  Reynolds     = -1.0      ! rho*U*L/mu, L=1.0[m]
  Gamma        = 1.4
  Gas_Constant = 287.05    ! [J/(kg*K)] = [m^2/(s^2*K)]
  Density      = -1.0      ! [kg/m^3]
  Temperature  = 300.0     ! [K]
  Pressure     = 100000.0  ! [Pa]
  AoA_Alpha    = 0.0       ! [deg]
  AoA_Beta     = 0.0       ! [deg]
  Prandtl      = 0.72
  Schmidt      = 0.72
  Suth_Muref   = 1.716e-5  ! [Pa*s] = [N*s/m^2]
  Suth_Tref    = 273.0     ! [K]
  Suth_Sref    = 110.4     ! [K]
/
\end{nmlblocklst}

\NOTE
Certain initialization options replace the reference conditions
by preset values to ensure repeatability of these benchmark cases.
%
See the \inp[nmlInitialization::Initialization]{Initialization}
input for details.

%==========================================================================

\begin{nmldescitem}
Mach = 0.4
\end{nmldescitem}
\inp*[nmlReference::Mach]{}[Mach][]

\begin{nmldescbody}
Reference Mach number.
Must be non-negative.
\end{nmldescbody}

%==========================================================================

\begin{nmldescitem}
Reynolds = -1.0  ! rho*U*L/mu, L=1.0[m]
\end{nmldescitem}
\inp*[nmlReference::Reynolds]{}[Reynolds][]

\begin{nmldescbody}
Reference Reynolds number, $Re = \rho U L / \mu$, based on
reference conditions.
%
Here,
$\rho$ is the reference density;
$U$ is the velocity magnitude computed from
  the reference Mach number and
  speed of sound at the reference temperature;
$L$ is the reference length of 1.0 meter; and 
$\mu$ is the viscosity computed from
  Sutherland's viscosity law using \eref{EQN:Viscosity_Sutherland}
  at the reference temperature.

How the Reynolds number is used depends on the input value:
\\
\begin{tabular}{ @{\qquad} r @{ \quad } p{0.85\linewidth} @{} }
%\begin{longtable}{ @{\qquad} r @{ \quad } p{0.85\linewidth} @{} }
  $\leq0$
       & The Reynolds number will be computed from the
         other reference flow properties.
         %
         This yields \scinum{8.7319}{6} for the default values shown above,
         but will change if you modify any of the other reference conditions.
         \\
  $>0$ & The Reynolds number will be used with the
         other reference conditions you specify
         to recalibrate the $\mu_0$ constant
         (\inp[nmlReference::SuthMuref]{Suth\_Muref})
         in \key[key:SutherlandsViscosityLaw]{Sutherland's viscosity law}[=],
         see \eref{EQN:Viscosity_Sutherland}.
         %
         This can be used to prescribe an artificially lower Reynolds number
         that permits better resolution of turbulent scales with
         large-eddy simulation.
         %
         It cannot be used to compute other unspecified conditions such as
         density, pressure, or temperature.
%\end{longtable}
\end{tabular}
\end{nmldescbody}

%==========================================================================

\FINALFIX{\newpage}
\begin{nmldescitem}
Gamma = 1.4
\end{nmldescitem}
\inp*[nmlReference::Gamma]{}[Gamma][]

\begin{nmldescbody}
Reference $\gamma$ (ratio of specific heats).
%
If the value specified is negative, it will be reset to the default
of 1.4 (the typical value for air).
\end{nmldescbody}

%==========================================================================

\begin{nmldescitem}
Gas_Constant = 287.05  ! [J/(kg*K)] = [m^2/(s^2*K)]
\end{nmldescitem}
\inp*[nmlReference::GasConstant]{}[Gas\_Constant][]

\begin{nmldescbody}
Reference specific gas constant (units of \unit{\J\per\kg\per\K}).
%
If the value specified is negative, it will be reset to the default
of 287.05 (the typical value for air).
 
\noindent\NOTE
In the code, the default value actually has more precision since it is computed
as $\texttt{Gas\_Constant} = R_\textrm{univ} / M_\textrm{air}$,
where $R_\textrm{univ}$ is the universal gas constant that is hard-coded as a
Fortran constant with the value
$8314.5\ \unit{\J\per\K\per\kmol}$,
and $M_\textrm{air}$ is the
molar mass of air which is also a hard-coded Fortran constant with the value
$28.9651159\ \unit{\kg\per\kmol}$.
\end{nmldescbody}

%==========================================================================

\begin{nmldescitem}
Density     = -1.0      ! [kg/m^3]
Temperature = 300.0     ! [K]
Pressure    = 100000.0  ! [Pa]
\end{nmldescitem}
\inp*[nmlReference::Density]{}[Density][]
\inp*[nmlReference::Temperature]{}[Temperature][]
\inp*[nmlReference::Pressure]{}[Pressure][]

\begin{nmldescbody}
Reference static
density $\rho$ (\unitDens{}),
temperature $T$ (\unitTemp{}), and 
pressure $p$ (\unitPres{}).

\NOTE
If all three inputs are positive, a consistency check will be performed.
%
If one of these three inputs is negative, it will be computed
from the other two using the equation of state, \eref{EQN:EqOfState}.
%
So, for the default values shown above,
density will evaluate to 1.160833 \unitDens{}.
\end{nmldescbody}

%==========================================================================

\begin{nmldescitem}
AoA_Alpha = 0.0  ! [deg]
AoA_Beta  = 0.0  ! [deg]
\end{nmldescitem}
\inp*[nmlReference::AoAAlpha]{}[AoA\_Alpha][]
\inp*[nmlReference::AoABeta]{}[AoA\_Beta][]

\begin{nmldescbody}
Reference flow angles (in degrees):
$\alpha$ is the
\key[key:AngleOfAttack]{angle of attack}[flow angles!angle of attack]
measured in the $x$-$y$ plane, and
$\beta$  is the
\key[key:AngleOfSideslip]{angle of sideslip}[flow angles!angle of sideslip]
measured in the $x$-$z$ plane.
%
These flow angles are defined by
\eref{EQN:FlowAngleVelX,EQN:FlowAngleVelY,EQN:FlowAngleVelZ} and
\fref{FIG:FlowAngle} on \pref{\pageref{FIG:FlowAngle}}.
\end{nmldescbody}

%==========================================================================

\begin{nmldescitem}
Prandtl = 0.72
\end{nmldescitem}
\inp*[nmlReference::Prandtl]{}[Prandtl][]

\begin{nmldescbody}
Prandtl number of the flow.
\end{nmldescbody}

%==========================================================================

\begin{nmldescitem}
Schmidt = 0.72
\end{nmldescitem}
\inp*[nmlReference::Schmidt]{}[Schmidt][]

\begin{nmldescbody}
Schmidt number of the flow.
\end{nmldescbody}

%==========================================================================

\begin{nmldescitem}
Suth_Muref = 1.716e-5  ! [Pa*s] = [N*s/m^2]
Suth_Tref  = 273.0     ! [K]
Suth_Sref  = 110.4     ! [K]
\end{nmldescitem}
\inp*[nmlReference::SuthMuref]{}[Suth\_Muref][]
\inp*[nmlReference::SuthTref]{}[Suth\_Tref][]
\inp*[nmlReference::SuthSref]{}[Suth\_Sref][]

\begin{nmldescbody}
These inputs correspond to the constants
  $\mu_0$,
  $T_0$, and
  $S_0$
in \key[key:SutherlandsViscosityLaw]{Sutherland's viscosity law}[=],
\eref{EQN:Viscosity_Sutherland}.

\noindent
\NOTE
If the reference Reynolds number specified is positive,
it will be used to recalibrate the $\mu_0$ constant.
See \inp[nmlReference::Reynolds]{Reynolds} for details.
\end{nmldescbody}

%==========================================================================


%##%\FINALFIX{\newpage}
\subsection{nml\_physics}
\label{SEC:nml_physics}

This \hyperref[SEC:User_Input]{namelist block}
is used to specify the physical modeling options
for the simulation.

\begin{nmlblocklst}[
,label={LST:nml_physics}
,caption={Namelist nml\_physics.}
]
&nml_physics
  Governing_Equations               = 2
  Moving_Ref_Frame%Rotation(1:3)    = [0.0,0.0,0.0]  ! [RPM]
  Moving_Ref_Frame%Translation(1:3) = [0.0,0.0,0.0]  ! [m/s]
  Moving_Ref_Frame%Spinup%Iter      = 0
  Moving_Ref_Frame%Spinup%Tstar     = 0.0
/
\end{nmlblocklst}
\inp*[nmlPhysics]{}[nml\_physics]


%==========================================================================

\begin{nmldescitem}
Governing_Equations = 2
\end{nmldescitem}
\inp*[nmlPhysics::GoverningEquations]{}[Governing\_Equations][]

\begin{nmldescbody}
Specifies the governing equations to use.

\begin{tabular}{ @{} r @{ = } p{0.85\linewidth} @{} }
  1 & Euler equations \\
  2 & Navier-Stokes equations
\end{tabular}
\end{nmldescbody}

%==========================================================================

\begin{nmldescitem}
Moving_Ref_Frame%Rotation(1:3) = [0.0,0.0,0.0]  ! [RPM]
\end{nmldescitem}
\inp*[nmlPhysics::MovingRefFrame/Rotation]{}[Moving\_Ref\_Frame\%Rotation][]

\begin{nmldescbody}
This specifies the rotation rate (in RPMs) of the moving reference frame
about each of the three Cartesian axes.
\end{nmldescbody}

%==========================================================================

\begin{nmldescitem}
Moving_Ref_Frame%Translation(1:3) = [0.0,0.0,0.0]  ! [m/s]
\end{nmldescitem}
\inp*[nmlPhysics::MovingRefFrame/Translation]{}[Moving\_Ref\_Frame\%Translation][]

\begin{nmldescbody}
This specifies the translation speed (in m\slash s) of the moving reference
frame in each of the three Cartesian directions.
\end{nmldescbody}

%==========================================================================

\begin{nmldescitem}
Moving_Ref_Frame%Spinup%Iter = 0
\end{nmldescitem}
\inp*[nmlPhysics::MovingRefFrame/Spinup/Iter]{}[Moving\_Ref\_Frame\%Spinup\%Iter][]

\begin{nmldescbody}
Number of time steps to spin up the rotation of the moving reference frame
from 0 to max RPMs.
%
The spin-up period used will be the slower between this and
\inp[nmlPhysics::MovingRefFrame/Spinup/Tstar]{Moving\_Ref\_Frame\%Spinup\%Tstar}.
\end{nmldescbody}

%==========================================================================

\begin{nmldescitem}
Moving_Ref_Frame%Spinup%Tstar = 0.0
\end{nmldescitem}
\inp*[nmlPhysics::MovingRefFrame/Spinup/Tstar]{}[Moving\_Ref\_Frame\%Spinup\%Tstar][]

\begin{nmldescbody}
If $>0$, this value specifies the non-dimensional time
(\key[key:tstar]{\tstar}[t *@\tstar]) interval, \eref{EQN:tstar},
to spin up the rotation of the moving reference frame from 0 to max RPMs.
%
The spin-up period used will be the slower between this and
\inp[nmlPhysics::MovingRefFrame/Spinup/Iter]{Moving\_Ref\_Frame\%Spinup\%Iter}.
\end{nmldescbody}

%==========================================================================


%##%\FINALFIX{\newpage}
\subsection{nml\_initialization}
\label{SEC:nml_initialization}
\inp*[nmlInitialization]{}[nml\_initialization][]

This \hyperref[SEC:User_Input]{namelist block}
is used to specify solution initialization.

\begin{nmlblocklst}[
,label={LST:nml_initialization}
,caption={Namelist nml\_initialization.}
]
&nml_initialization
  Initialization                = 1
  Jet_Init_Func                 = 2
  Load_Restart_File             = F
  Restart_File                  = "out.rst"
  Initial_Solution_File         = ""
  Initial_Solution_Format       = -99        ! -99 = Unknown
/
\end{nmlblocklst}
% CFL_Reduction_For_Nan_Restart = 2.0

The \key*[key:InitializationSequence]{initialization sequence}[Initialization Sequence]
is defined with the following precedence:
\begin{itemize}
%
\item
If the file specified by \inp[nmlInitialization::InitialSolutionFile]{Initial\_Solution\_File}
exists and \inp[nmlInitialization::InitialSolutionFormat]{Initial\_Solution\_Format} is set to
a valid value, then the solution data read from the specified file will be
used as the initialization for a new simulation.
%
Note that \inptext{Initial\_Solution\_File} and \inptext{Initial\_Solution\_Format} can be
set to read an externally generated solution or a GFR restart file.
%
The iteration and solution time will be set to zero.
%
\item
If the file specified by \inp[nmlInitialization::RestartFile]{Restart\_File}
exists and
\inp[nmlInitialization::LoadRestartFile]{Load\_Restart\_File}=\texttt{T},
then the solution will be initialized using the data read from
the specified restart file.
The iteration and solution time will continue from that in the restart file.
%
\item
Otherwise, the initial solution will be defined based on
the \inp[nmlInitialization::Initialization]{Initialization} option.
The iteration and solution time will be set to zero.
\end{itemize}

\noindent\NOTE
To continue a simulation from a restart file while preserving the iteration
and solution time, one should set
  \inptext{Initial\_Solution\_File}=\texttt{""},
  \inptext{Load\_Restart\_File}=\texttt{T}, and
  \inptext{Restart\_File} to the name of the input restart file.
Many users forget to make these changes, and their simulations start over
from the previously specified initialization rather than continuing the
simulation.


\devNotes{
%
The files specified by \inptext{Initial\_Solution\_File} and \inptext{Restart\_File}
need to be checked for validity as part of the all the error checks after the
namelist is read.
%
Currently, I think it goes all the way to trying to read the specified restart
file before finding out if it does not exist, and at that stage it is not
immediately clear if it is \inptext{Initial\_Solution\_File} or \inptext{Restart\_File}
that is responsible.
%
Earlier error checking can allow for a more graceful abort with a clearer
message like
%
\begin{InputExample}
  \item[\quoterm{Tried to read \inptext{Initial\_Solution\_File}, but file does not
  exist!}]
\end{InputExample}
%
immediately after reading the namelist.
%
Probably need to also add error checking on if both methods are \quoterm{active}
and whether that is intentional or not (or do we just use the above
precedence?).
%
} % end devNotes
%
\vspace{1.0em} % Add vertical space before next item.

%==========================================================================

%%% The following notes should help when updating this manual in the future.
%%%
%%% For cases that support an internally generated grid (IG),
%%% see subroutine create_internal_grid in internal_grid_mod.f90.
%%% itestcase:
%%%   |Transport_Problems|   = |[Vortex_Problems,Advec_Problems]|
%%%                          = |[Shu_Vortex,Vortex_Transport,Shu_MultiVortex,
%%%                              Density_Transport,Entropy_Transport]|
%%%                          = |[5,6,55,7,8]|
%%%    Taylor_Green_Vortex   =   9
%%%   |Taylor_Couette_Flow|  = |21|
%%%   |Freestream_TCF|       = |22|
%%%   |Pipe_Channel_Flow|    = |25|
%%%   |CH4_Flame_Speed_Test| = |19|
%%%   |Sod_Shock_Tube|       = |23|
%%%   |Density_Pulse_Test|   = |24|
%%%   |Shu_Osher_Test|       = |26|
%%%   |Advection_JiangShu|   = |27|
%%%
%%% For cases that modify the reference conditions (RC),
%%% see subroutine initialize_reference in initialization_mod.f90.
%%% itestcase:
%%%   Advec_Problems = [Density_Transport, Entropy_Transport] = [7,8]
%%%  |Shu_Vortex|         =  5,-5
%%%  |Shu_MultiVortex|    =  55,-55
%%%  -Vortex_Transport    = -6
%%%   Taylor_Green_Vortex =  9
%%%   Taylor_Couette_Flow = 21
%%%   Sod_Shock_Tube      = 23
%%%   Shu_Osher_Test      = 26
%%%   Advection_JiangShu  = 27
%%%
%%% For cases that provide an exact solution (ES),
%%% see sub output_variable in write_cgns_mod.f90.
%%%   Sod_Shock_Tube         = 23, ExactDensity, ExactPressure
%%%   Advection_JiangShu     = 27, ExactDensity, ExactPressure
%%%
%%% FIXME:
%%% For cases that modiy the post-procesing (PP),
%%% see sub init_generic_mod in generic_mod.f90 that writes screen info.
%%% itestcase:
%%%   Inviscid_Gaussian_Bump =  3
%%%   Laminar_BndLyr         =  4
%%%   THX_Phase_3            = 20
%%% see sub get_output_variable in postproc_mod.f90.
%%% itestcase:
%%%   Inviscid_Gaussian_Bump =  3, Entropy calc

\begin{nmldescitem}
Initialization = 1
\end{nmldescitem}
\inp*[nmlInitialization::Initialization]{}[Initialization][]

\begin{nmldescbody}
Specifies the test-case\slash problem to solve, primarily affecting the
initialization and post-processing of the simulation.

\begingroup % limit scope of local commands
\setlength{\tabcolsep}{0.75\tabcolsep} % shrink column gap
%%% If descriptions are too long, need p{} to line wrap.
%\begin{longtable}{ @{} r c c c p{0.75\linewidth} @{} }
\begin{longtable}{ @{} r c c c l @{} }
 \toprule
 Value & IG & RC & PP & Description \\
 \midrule
 \endhead
    1 & . & . & . & Generic flow, initialize to reference conditions.\\
    2 & Y & . & . & Channel flow.\\
    3 & . & . & Y & Inviscid Gaussian bump.\\
    4 & . & . & Y & Laminar subsonic boundary layer.\\
    5 & Y & Y & . & Shu isentropic Euler vortex, propagating diagonally.\\
                    \nopagebreak
   -5 & Y & Y & . & Shu isentropic Euler vortex, stationary.\\
   55 & Y & Y & . & Shu isentropic Euler vortex, propagating diagonally,
                    4 total vortices.\\
                    \nopagebreak
  -55 & Y & Y & . & Shu isentropic Euler vortex, stationary,
                    4 total vortices.\\
    6 & Y & . & . & Isentropic Euler vortex propagating in
                    \inp[nmlReference::AoAAlpha]{AoA\_Alpha}
                    direction.\\
                    \nopagebreak
   -6 & Y & Y & . & Isentropic Euler vortex propagating only in the
                    $y$-direction.\\
    7 & Y & Y & . & Linear advection equation advecting a density profile.\\
    8 & Y & Y & . & Linear advection equation advecting an entropy profile.\\
    9 & Y & Y & . & Taylor-Green vortex problem.\\
   10 & . & . & . & Subsonic jet cases using the SMC000 nozzle.\\
                    \nopagebreak
  -10 & . & . & . & Subsonic jet cases using the ARN2 nozzle.\\
   11 & . & . & . & Supersonic jet cases.\\
   12 & . & . & . & Infinite cylinder.\\
   13 & . & . & . & Freestream preservation.\\
   14 & . & . & . & Reacting backstep.\\
   15 & . & . & . & 2D diffusion flame fueled with hydrogen.\\
   16 & . & . & . & 2D diffusion flame fueled with methane.\\
   17 & . & . & . & Nitrogen\slash Oxygen diffusion.\\
   18 & . & . & . & Compression corner.\\
   19 & Y & . & . & Methane flame speed test.\\
   20 & . & . & Y & THX Phase 3.\\
   21 & Y & Y & . & Taylor-Couette flow.\\
   22 & Y & . & . & Freestream Taylor-Couette flow.\\
   23 & Y & Y & . & 1D Sod shocktube problem.\\
   24 & Y & . & . & 1D Density pulse test.\\
   25 & Y & . & . & Annular pipe channel flow.\\
   26 & Y & Y & . & 1D Shu-Osher problem.\\
   27 & Y & Y & . & 1D Jiang-Shu density pulse test.\\
   \bottomrule
%  \multicolumn{2}{l}{%
%    \footnotesize{%
%      $^\dagger$
%      Supports 
%      \inp[nmlGrid::GridFormat]{Grid\_Format} = -1 (internally generated grid),
%      but external grids can be used as well.
%    }
%  }
\end{longtable}
\endgroup

\noindent
In the list above, the following designations are made:
\begin{enumerate}[align=parleft]
\item[IG]
This input option supports an
\key*[key:InternallyGeneratedGrid]{internally generated grid}[=]
by setting \inp[nmlGrid::GridFormat]{Grid\_Format} = -1 (internal grid),
but external grids can be used as well.
Grid resolution can be controlled as described under the
\inp[nmlGrid::InternalGridDoF]{Internal\_Grid\_DoF} input.
\item[RC]
This input option overrides any user-specified reference conditions
in \inp[nmlReference]{nml\_reference} with preset values inside the code,
which will be written to the screen during code execution.
\item[PP]
This input option only affects certain post-processed quantities,
such as the screen output,
and the flow will be initialized the same as option 1 (generic flow).
\end{enumerate}

\NOTE
This initialization is superseded when either loading a restart file by
specifying \inp[nmlInitialization::LoadRestartFile]{Load\_Restart\_File}=\texttt{T}
or using the input \inp[nmlInitialization::InitialSolutionFile]{Initial\_Solution\_File}
to initializing a simulation from a valid solution file.
%
Instead, the flowfield will be read from the corresponding file.
%
See the description of the
\key[key:InitializationSequence]{initialization sequence}
on \pageref{key:InitializationSequence}.

\end{nmldescbody}

%==========================================================================

\begin{nmldescitem}
Jet_Init_Func = 2
\end{nmldescitem}
\inp*[nmlInitialization::JetInitFunc]{}[Jet\_Init\_Func][]

\begin{nmldescbody}
Sometimes when initializing a jet simulation from freestream conditions,
early transients within the nozzle can be numerically unstable,
and a better initial solution is desirable.

If \inp[nmlInitialization::Initialization]{Initialization}
is $\pm 10$ or $11$, \inptext{Jet\_Init\_Func} can be used to
specify different error function distributions to initialize the
solution for the ARN2, SMC000, and the supersonic jet geometries.
%
The jet geometry (\forexample ARN2, SMC000\cite{Bridges_2011}) used for the
initialization function depends on the value of \inptext{Initialization}.

\begin{tabular}{ @{} r @{ = } p{0.85\linewidth} @{} }
  1 & Constant radial blending. \\
  2 & Blending in both x and r with a constant radial spreading. \\
  3 & Blending in both x and r with a downstream radial spreading.
\end{tabular}
\\
For any other value, the volume will be initialized to the
freestream reference conditions.

\IDEA{More detail could be provided on the profile functions.}
\end{nmldescbody}

%==========================================================================

\begin{nmldescitem}
Load_Restart_File = F
\end{nmldescitem}
\inp*[nmlInitialization::LoadRestartFile]{}[Load\_Restart\_File][]

\begin{nmldescbody}
Specifies whether or not to read a restart file.

\begin{tabular}{ @{} r @{ = } p{0.85\linewidth} @{} }
  T & Read and initialize from the restart file given by 
      \inp[nmlInitialization::RestartFile]{Restart\_File}.\\
  F & Do not read a restart file, but instead initialize according to
      \inp[nmlInitialization::Initialization]{Initialization}.
\end{tabular}

\NOTE
This will be superseded by
\inp[nmlInitialization::InitialSolutionFile]{Initial\_Solution\_File} if the
  specified value is an existing file and
\inp[nmlInitialization::InitialSolutionFormat]{Initial\_Solution\_Format} is set to a valid
value.
%
The only difference is that this option keeps the iteration number and
simulation time from the restart file, whereas \inptext{Initial\_Solution\_File}
resets both to 0.
%
See the description of the
\key[key:InitializationSequence]{initialization sequence}
on \pageref{key:InitializationSequence}.
\end{nmldescbody}

%==========================================================================

\begin{nmldescitem}
Restart_File = "out.rst"
\end{nmldescitem}
\inp*[nmlInitialization::RestartFile]{}[Restart\_File][]

\begin{nmldescbody}
File path and name to the solution restart file that will be \emph{read} if
\inp[nmlInitialization::LoadRestartFile]{Load\_Restart\_File}
is true.
%
%This takes precedence over
This supersedes
\inp[nmlInitialization::Initialization]{Initialization},
but will be superseded by
\inp[nmlInitialization::InitialSolutionFile]{Initial\_Solution\_File}.
%
See the description of the
\key[key:InitializationSequence]{initialization sequence}
on \pageref{key:InitializationSequence}.
%
The solution will continue from the iteration count and simulation time
of the restart file.

\NOTE
This has no effect on the restart files that are \emph{written},
typically to:
\\ % Prevent filename from entering right margin.
  \texttt{Restart\_files/out.<iteration>.rst}.
\end{nmldescbody}

%==========================================================================

%\begin{nmldescitem}
%CFL_Reduction_For_Nan_Restart = 2.0
%\end{nmldescitem}
%\inp*[nmlInitialization::CFLReductionForNanRestart]{}[CFL\_Reduction\_For\_Nan\_Restart][]

%\begin{nmldescbody}
%\NOTE
%This feature is not working.
%
%This feature is intended to provide for an automatic recovery
%for failed solutions by rolling back to a previous state
%and reducing the timestep.
%%
%If an NaN is encountered in a solution, the code will restart from
%the most recent restart file.
%%
%The CFL from when the NaN occurred is divided by
%\inptext{CFL\_Reduction\_For\_Nan\_Restart} and used to continue
%the simulation.
%
%This feature is most useful on computing systems where manually
%modifying and resubmitting a job incurs long delays due to the
%queuing system.
%%
%Unfortunately, in many instances this technique is insufficient
%to overcome the numerical exception, and manual adjustments
%to the limiting \andor filtering are necessary.
%
%%\FIXME{
%%\begin{enumerate}
%%\item
%%REMOVE this input from the code and documentation. Does not work well.
%%\item
%%Current code bypasses automatic restart if
%%  $\inptext{CFL\_Reduction\_For\_Nan\_Restart} > 1$.
%%\item
%%There is no other input that enables\slash disables this automatic recovery.
%%  It is\slash was intended to always be on.
%%\item
%%What happens if the input file uses CFL ramping?  Does ramping continue?
%%\end{enumerate}
%%}
%\end{nmldescbody}

%==========================================================================

\begin{nmldescitem}
Initial_Solution_File = ""
\end{nmldescitem}
\inp*[nmlInitialization::InitialSolutionFile]{}[Initial\_Solution\_File][]

\begin{nmldescbody}
Path to a file that can be \emph{read} for initializing the solution.
%
This can be an externally generated solution or a GFR restart file.
%
This is essentially the same as loading a restart file,
except the \emph{\textbf{iteration count and simulation time are reset to 0}}.
%
See \key[key:InitializationSequence]{initialization sequence}
on \pref{\pageref{key:InitializationSequence}} for important details.
%
When specifying \inptext{Initial\_Solution\_File}, one must also specify a valid value
for \inp[nmlInitialization::InitialSolutionFormat]{Initial\_Solution\_Format} as its default
value will prevent this file from being read.

\NOTE
This has no effect on the solution files that are \emph{written}.

\NOTE
After initialization is complete, set \inptext{Initial\_Solution\_File}=\texttt{""}
so that subsequent simulations will continue from the previous restart file
rather than starting over again.
See \pref{\pageref{key:InitializationSequence}}.


\end{nmldescbody}

%==========================================================================

\begin{nmldescitem}
Initial_Solution_Format = -99  ! -99 = Unknown
\end{nmldescitem}
\inp*[nmlInitialization::InitialSolutionFormat]{}[Initial\_Solution\_Format][]

\begin{nmldescbody}
Format of the file specified by input
\inp[nmlInitialization::InitialSolutionFile]{Initial\_Solution\_File}.

\begin{tabular}{ @{} r @{ = } p{0.85\linewidth} @{} }
  0 & Restart file format.\\
  3 & Plot3D file format (experimental).
\end{tabular}
\\
\NOTE
Any other value will prevent the file specified by
\inp[nmlInitialization::InitialSolutionFile]{Initial\_Solution\_File}
from being read, and another means will be used to initialize the solution.

\end{nmldescbody}

%==========================================================================


%#%%\FINALFIX{\newpage}%
\subsection{nml\_timestep}
\label{SEC:nml_timestep}
\inp*[nmlTimestep]{}[nml\_timestep][]

This \hyperref[SEC:User_Input]{namelist block}
is used to specify time stepping features.

\begin{nmlblocklst}[
,label={LST:nml_timestep}
,caption={Namelist nml\_timestep.}
]
&nml_timestep
  Runge_Kutta_Scheme        =    113
  Num_Timesteps_Min         =    100
  Num_Timesteps             =  25000
  Final_Time                =  -10.0  ! non-dim time, t*
  Timestep_Type             =      1
  Timestep                  = 1.0e-6  ! t[s] * ( a[m/s] / L=1.0[m])
  CFL                       =    1.0
  CFL_Ramping_Beg           =    0.1
  CFL_Ramping_End           =    1.0
  CFL_Ramping_Iterations    =      1
  CFL_Ramping_Start_Iter    =      0
  Convergence_Order_Abs_Res =  -16.0
  Convergence_Order_Max_Res =  -16.0
/
\end{nmlblocklst}

%==========================================================================

\begin{nmldescitem}
Runge_Kutta_Scheme = 113
\end{nmldescitem}
\inp*[nmlTimestep::RungeKuttaScheme]{}[Runge\_Kutta\_Scheme][]

\begin{nmldescbody}
Specifies which Runge-Kutta (RK) method to use.
%
See \sref{SEC:Time:Methods} for additional details.
%
Speedup is relative to the 3-stage/\nth{3}-order SSP ERK method.

\begin{longtable}{ r @{ = } r @{/} p{3.55in} c }
  \toprule
  \multicolumn{1}{l@{ \phantom{=} }}{Option} &
  \multicolumn{2}{@{}l}{Description} &
  Speedup \\
  \midrule%\\[-0.75em]
  \endhead
%\multicolumn{4}{l}{}\\[-0.75em]
         \multicolumn{1}{r@{}}{}
       & \multicolumn{2}{@{}l}{\textbf{Recommended options.}} & \\
   113 &  11-stage & \nth{3}-order VOERK\cite{AlJahdali_2022} & 1.33x \\
   154 &  15-stage & \nth{4}-order VOERK\cite{AlJahdali_2022} & 1.33x \\
     4 &   5-stage & \nth{4}-order SSP ERK\cite{Spiteri_2002} & 1.22x \\
   165 &  16-stage & \nth{5}-order VOERK\cite{AlJahdali_2022} & 1.18x \\
   943 &   9-stage & \nth{4}(\nth{3})-order
                     EPERK FSAL\cite{Ranocha_2022}            & 1.22x \\
   532 &   5-stage & \nth{3}(\nth{2})-order
                     EPERK FSAL\cite{Ranocha_2022}            & 1.09x
\\[0.5em]
\pagebreak[3]
%
         \multicolumn{1}{r@{}}{}
       & \multicolumn{2}{@{}l}{\textbf{Other options that work well.}}   & \\
    -2 &   2-stage & \nth{2}-order Classic RK                       & n/a   \\
    -3 &   3-stage & \nth{3}-order Classic RK                       & n/a   \\
    -4 &   4-stage & \nth{4}-order Classic RK                       & n/a   \\
     2 &   2-stage & \nth{2}-order SSP ERK\cite{Gottlieb_1998}      & n/a   \\
     3 &   3-stage & \nth{3}-order SSP ERK\cite{Gottlieb_1998}      & 1.00x \\
     5 &   5-stage & \nth{4}-order Carpenter-Kennedy low storage
                                   RK\cite{Carpenter_1994}          & 0.99x \\
    53 &   5-stage & \nth{3}-order VOERK\cite{AlJahdali_2022}       & 1.08x \\
    64 &   6-stage & \nth{4}-order VOERK\cite{AlJahdali_2022}       & 1.02x \\
    85 &   8-stage & \nth{5}-order VOERK\cite{AlJahdali_2022}       & 1.05x \\
  -943 &   9-stage & \nth{4}(\nth{3})-order
                     EPERK non-FSAL\cite{Ranocha_2022}              & n/a   \\
  -532 &   5-stage & \nth{3}(\nth{2})-order
                     EPERK non-FSAL\cite{Ranocha_2022}              & n/a
\\[0.5em]
\pagebreak[3]
%
         \multicolumn{1}{r@{}}{}
       & \multicolumn{2}{@{}l}{\textbf{Options that do not perform as well.}}&\\
  $\pm$1 &   1-stage & \nth{1}-order Forward Euler method       & n/a   \\
     -53 &   5-stage & \nth{3}-order AOERK\cite{AlJahdali_2022} & n/a \\
    -113 &  11-stage & \nth{3}-order AOERK\cite{AlJahdali_2022} & n/a \\
     -64 &   6-stage & \nth{4}-order AOERK\cite{AlJahdali_2022} & n/a \\
    -154 &  15-stage & \nth{4}-order AOERK\cite{AlJahdali_2022} & n/a \\
     -85 &   8-stage & \nth{5}-order AOERK\cite{AlJahdali_2022} & n/a \\
    -165 &  16-stage & \nth{5}-order AOERK\cite{AlJahdali_2022} & 0.98x \\
      54 &   6-stage & \nth{5}(\nth{4})-order
                       Dormand-Prince EPERK\cite{Dormand_1980}  & 0.64x \\
     853 &  12-stage & \nth{8}(\nth{5})-order
                       Dormand-Prince EPERK\cite{Hairer_1993}   & 0.61x \\
    1054 &  10-stage & \nth{5}(\nth{4})-order
                       EPERK FSAL\cite{Ranocha_2022}            & 0.85x \\
   -1054 &  10-stage & \nth{5}(\nth{4})-order
                     EPERK non-FSAL\cite{Ranocha_2022}        & n/a \\
  \bottomrule
\end{longtable}

\noindent
Acronyms used in the description above:
\\
\begin{tabular}{ @{\qquad} l @{ : } p{0.75\linewidth} @{} }
  SSP   & Strong stability preserving \\
  VOERK & Vortex-optimized explicit Runge-Kutta (ERK) \\
  AOERK & Advection-optimized ERK \\
  FSAL  & First-same-as-last \\
  EPERK & Embedded-pair ERK with adaptive time stepping
\end{tabular}
\\

\NOTE
EPERK methods do not work with limiting or filtering.

\NOTE
The value in parenthesis for EPERK methods is the order of the embedded method.
%
The difference between the actual method and embedded method represents the
time-error from the time step.
%
This is used by the adaptive time stepping proportional–integral–derivative
(PID) controller to compute the size of the next time step.
\end{nmldescbody}

%==========================================================================

\begin{nmldescitem}
Num_Timesteps_Min = 100
\end{nmldescitem}
\inp*[nmlTimestep::NumTimestepsMin]{}[Num\_Timesteps\_Min][]

\begin{nmldescbody}
The minimum number of time steps to complete before checking whether any of
the convergence criteria have been met.

\NOTE When initializing a simulation to a uniform flow, the very
small changes that occur during the first few time steps can trigger a false
positive convergence check if this is set to a small number,
\forexample $\leq5$.
\end{nmldescbody}

%==========================================================================

\begin{nmldescitem}
Num_Timesteps = 25000
\end{nmldescitem}
\inp*[nmlTimestep::NumTimesteps]{}[Num\_Timesteps][]

\begin{nmldescbody}
Number of time steps for the simulation to run, unless
\inp[nmlTimestep::FinalTime]{Final\_Time}
is reached before completing all time steps.
\end{nmldescbody}

%==========================================================================

\begin{nmldescitem}
Final_Time = -10.0  ! non-dim time, t*
\end{nmldescitem}
\inp*[nmlTimestep::FinalTime]{}[Final\_Time][]

\begin{nmldescbody}
Stops the simulation when the time within the simulation reaches this
final time value.
%
This value is specified in nondimensional units
(\key[key:tstar]{\tstar}[t *@\tstar]), \eref{EQN:tstar}, based on the input
provided for
\inp[nmlOutput::TimeScalingFactor]{Time\_Scaling\_Factor}.
%
For the isentropic vortex problem
(\inp[nmlInitialization::Initialization]{Initialization}=5 or 55), a
negative value specifies the total number of periods to run the simulation.
%
For all other cases, a negative value disables the use of
\inptext{Final\_Time} to stop a simulation.

\NOTE
The nondimensional simulation time
(\key[key:tstar]{\tstar}[t *@\tstar]), \eref{EQN:tstar},
is written to standard output with the convergence information
and to the residual file.
\end{nmldescbody}

%==========================================================================

\begin{nmldescitem}
Timestep_Type = 1
\end{nmldescitem}
\inp*[nmlTimestep::TimestepType]{}[Timestep\_Type][]

\nopagebreak
\begin{nmldescbody}
This determines the method used to compute the time step size and the type of
time stepping to use for the simulation (unless using one of the EPERK methods
for adaptive time stepping, see note below).
%
\\
\begin{tabular}{ @{\qquad} r @{ = } p{0.80\linewidth} @{} }
  -2 & Local time step from the inviscid eigenvalue. (Experimental) \\
  -1 & Local time step from the spectral radius. \\
   0 & Global time step from user specified time step. \\
   1 & Global time step from the spectral radius. \\
   2 & Global time step from the inviscid eigenvalue.  (Experimental)
\end{tabular}
 
\begin{InputList}[font=\mdseries]
  \item[Local time stepping:]
    At the beginning of each solver iteration, each grid cell computes its own
    unique time step size that is used to update the solution within the cell at
    the end of each Runge-Kutta stage.
    %
    Since all cells are updated with different time step sizes, this can only be
    used for steady-state simulations and cannot be used for time-accurate
    simulations (\forexample LES) which require all cells to be synchronized in
    time.
  \item[Global time stepping:]
    At the beginning of each solver iteration, the time step for each cell is
    computed in the same way as with local time stepping.
    %
    However, the minimum time step size across all grid cells is taken to be the
    global time step size that all cells use to update their solutions in time.
    %
    Global time stepping is required for time-accurate simulation like LES, but
    is generally inefficient for steady-state problems.
  \item[Inviscid eigenvalue:]
    The maximum inviscid eigenvalue is used to compute a time step size at
    each solution point within a cell, and the minimum of these values is
    taken to be the time step for the cell.
  \item[Spectral radius:]
    The time step within each cell is taken as the minimum value of the
    spectral radius evaluated at every face\slash flux point of the cell.
  \item[User specified time step:]
    A constant global time step is obtained from the user input for
    \inp[nmlTimestep::Timestep]{Timestep}.
    %
    \NOTE The input \inptext{Timestep} is specified as a nondimensional value.
    It is the physical time step nondimensionalized by the reference speed of
    sound and a length of 1 meter.
%
%   A constant global time step is computed from the inputs
%   $\inp[nmlTimestep::CFL]{CFL} \times
%    \inp[nmlTimestep::Timestep]{Timestep}$.
%   \newline
%   %
%   For this option, \inptext{CFL}=1.0 unless CFL ramping is being used.
%   See \inp[nmlTimestep::Timestep]{Timestep}.
\end{InputList}

\vspace{1.0ex}\noindent\NOTE
For each of the methods above,
once the time step is determined, the effective $\Delta t$ will be computed as
\begin{equation}
  \Delta t = CFL \times \text{time step}
\end{equation}
where $CFL$ is obtained from
  \inp[nmlTimestep::CFL]{CFL}
or from
  \inp[nmlTimestep::CFLRampingBeg]{CFL\_Ramping\_Beg}/%
  \inp[nmlTimestep::CFLRampingEnd]{End}[CFL\_Ramping\_End]
if CFL ramping is being used.
%
If \inptext{Timestep\_Type}=0 (user specified time step), then
\inp[nmlTimestep::CFL]{CFL} will be set to 1.0,
however ramping may still be used.
See \inp[nmlTimestep::Timestep]{Timestep}.
\\ % Adds an extra blank line.

\noindent\NOTE
If using one of the EPERK methods for adaptive time stepping, the CFL column
written to \file{stdout} for each iteration will contain an approximate CFL
value, $\text{CFL}_\textsc{eperk}$, for the adaptive time step size relative
to the time step size that would have been computed for one of the CFL-based
ERK methods, \thatis
%
\[
  \text{CFL}_\textsc{eperk} \approx
                   \frac{\Delta t_\textsc{eperk}}{\Delta t_\textsc{erk}}
\]
%
where $\Delta t_\textsc{eperk}$ is the time step size computed by the PID
controller for the EPERK method, and $\Delta t_\textsc{erk}$ is the time step
size that would have been computed using the method specified by
\inptext{Timestep\_Type} for a CFL-based ERK method with $\text{CFL}=1.0$.

\devNotes{
For \inptext{Timestep\_Type=0} (user specified time step), have the
code force CFL=1.0 rather than just ignore it.
}
\end{nmldescbody}

%==========================================================================

\begin{nmldescitem}
Timestep = 1.0e-6  ! t[s] * (a[m/s] / L=1.0[m])
\end{nmldescitem}
\inp*[nmlTimestep::Timestep]{}[Timestep][]

\begin{nmldescbody}
Specifies the size of the time step that is used when
\inp[nmlTimestep::TimestepType]{Timestep\_Type}=0
(constant global time step).
%\thatis $\inptext{Timestep\_Type}=0$.
%
The value specified by this input is nondimensionalized by
$L/a_{\infty}$, the reference length (1~m) and 
speed of sound.
Thus, the dimensional time step size will be
\inp[nmlTimestep::Timestep]{Timestep}$ * L / a_\infty$ seconds.
 
When a simulation is running, the nondimensional time step size is
the second column written to stdout, so the value of this input is
identical to that output under the conditions
\inp[nmlTimestep::TimestepType]{Timestep\_Type}=0 and
\inp[nmlTimestep::CFL]{CFL}=1.0.
%
%\CHANGECODE
When using \inp[nmlTimestep::TimestepType]{Timestep\_Type}=0
(constant global time step), the value used for
\inptext{CFL} is $1.0$ unless CFL ramping is being used.
%
If CFL ramping is being used, the ramping parameters will be used
to evaluate the value of \inptext{CFL} to be used for each time step,
but \inptext{CFL} will revert to a value of $1.0$ once the
ramping completes.

\NOTE
This nondimensional time step is written to standard output
with the convergence information.
%
The \inptext{Timestep} is nondimensionalized by
the reference length (1~m) and speed of sound, whereas
the solution time that is written to standard output and
\inp[nmlTimestep::FinalTime]{Final\_Time} are both nondimensionalized by
\inp[nmlOutput::TimeScalingFactor]{Time\_Scaling\_Factor}
via \eref{EQN:tstar}.

\end{nmldescbody}

%==========================================================================

\begin{nmldescitem}
CFL = 1.0
\end{nmldescitem}
\inp*[nmlTimestep::CFL]{}[CFL][]

\begin{nmldescbody}
Analogous to the CFL condition, this is used to scale the time step size to
get a maximum stable time step. The maximum value this can be while also
keeping a simulation stable is very dependent on the type of simulation,
flow conditions, and input settings.  Experimentation is generally
required to find a good value.  A lower value is more likely to keep a
simulation stable, but will require more time steps to reach either
convergence or a given simulation time. Typical values are
$0.1\leq\inptext{CFL}\leq2$, but past experience has found simulations that
were stable with values up to $10.0$ for certain Runge-Kutta schemes.

%\NOTE
If CFL ramping is actively being used, that computed value
will be applied.
Once CFL ramping has completed, the value of this input will be used.

\devNotes{
CHANGECODE:
This value will automatically be set to $1.0$ if using
\inp[nmlTimestep::TimestepType]{Timestep\_Type}=0
(constant global time step), unless CFL ramping is active.
}

\end{nmldescbody}

%==========================================================================

\begin{nmldescitem}
CFL_Ramping_Beg = 0.1
CFL_Ramping_End = 1.0
\end{nmldescitem}
\inp*[nmlTimestep::CFLRampingBeg]{}[CFL\_Ramping\_Beg][]
\inp*[nmlTimestep::CFLRampingEnd]{}[CFL\_Ramping\_End][]

\begin{nmldescbody}
Beginning and ending CFL values to use with the linear CFL ramping that is
active when
$\inp[nmlTimestep::CFLRampingIterations]{CFL\_Ramping\_Iterations}\geq2$.

See also \sref{SEC:Files:Interrupt:Adjust_CFL} for adjusting the CFL
while the solver is running.
\end{nmldescbody}

%==========================================================================

\begin{nmldescitem}
CFL_Ramping_Iterations = 1
\end{nmldescitem}
\inp*[nmlTimestep::CFLRampingIterations]{}[CFL\_Ramping\_Iterations][]

\begin{nmldescbody}
The number of time steps used to linearly ramp the CFL value,
starting at \inptext{CFL\_Ramping\_Beg} and
ending   at \inptext{CFL\_Ramping\_End}.
%
$\inptext{CFL\_Ramping\_Iterations}\geq2$ activates CFL ramping.

\NOTE
A new restart file will be written just before the
ramping schedule begins in case it is too aggressive and causes the
simulation to fail.

See also \sref{SEC:Files:Interrupt:Adjust_CFL} for adjusting the CFL
while the solver is running.
\end{nmldescbody}

%==========================================================================

\begin{nmldescitem}
CFL_Ramping_Start_Iter = 0
\end{nmldescitem}
\inp*[nmlTimestep::CFLRampingStartIter]{}[CFL\_Ramping\_Start\_Iter][]

\begin{nmldescbody}
Specifies the time step iteration at which to begin the CFL ramping schedule
specified by
\inp[nmlTimestep::CFLRampingBeg]{CFL\_Ramping\_Beg},
\inp[nmlTimestep::CFLRampingEnd]{CFL\_Ramping\_End}, and
\inp[nmlTimestep::CFLRampingIterations]{CFL\_Ramping\_Iterations}.
%
This input is primarily used with the \inptext{ADJUST\_CFL} control file
(\sref{SEC:Files:Interrupt:Adjust_CFL})
that can be used to modify the CFL ramping schedule of a running simulation.

\NOTE
A new restart file will be written just before the
ramping schedule begins in case it is too aggressive and causes the
simulation to fail.


See also \sref{SEC:Files:Interrupt:Adjust_CFL} for adjusting the CFL
while the solver is running.
\end{nmldescbody}

%==========================================================================

\begin{nmldescitem}
Convergence_Order_Abs_Res = -16.0
\end{nmldescitem}
\inp*[nmlTimestep::ConvergenceOrderAbsRes]{}[Convergence\_Order\_Abs\_Res][]

\begin{nmldescbody}
This is the convergence goal for the orders of magnitude of the absolute
residual.
\end{nmldescbody}

%==========================================================================

\begin{nmldescitem}
Convergence_Order_Max_Res = -16.0
\end{nmldescitem}
\inp*[nmlTimestep::ConvergenceOrderMaxRes]{}[Convergence\_Order\_Max\_Res][]

\begin{nmldescbody}
This is the convergence goal for the orders of magnitude reduction of the
residual relative to the maximum value of the residual during the
simulation.
\end{nmldescbody}

%==========================================================================


\subsection{nml\_bc}
\label{SEC:nml_bc}
\inp*[nmlBC]{}[nml\_bc][]

This \hyperref[SEC:User_Input]{namelist block}
is used to specify boundary conditions.

\begin{nmlblocklst}[
,escapeinside={!@}{!@}
,label={LST:nml_bc}
,caption={Namelist nml\_bc.}
]
&nml_bc
  Walls_Are_Exact                          = T
  Sub_Inflow_Method                        = 4
  Use_BC_Conflicts                         = T
  BC_Input(1)%Name                         = ""
  BC_Input(1)%Set_BC_Value                 = -9999
  BC_Input(1)%Temperature                  = -1.0               ! [K]
  BC_Input(1)%Pressure                     = -1.0               ! [Pa]
  BC_Input(1)%Density                      = -1.0               ! [kg/m^3]
  BC_Input(1)%Mach                         = -1.0
  BC_Input(1)%VelocityX                    = 0.0                ! [m/s]
  BC_Input(1)%VelocityY                    = 0.0                ! [m/s]
  BC_Input(1)%VelocityZ                    = 0.0                ! [m/s]
  BC_Input(1)%Total_Temperature            = -1.0               ! [K]
  BC_Input(1)%Total_Pressure               = -1.0               ! [Pa]
  BC_Input(1)%Total_Density                = -1.0               ! [kg/m^3]
  BC_Input(1)%AoA_Alpha                    = 0.0                ! [deg]
  BC_Input(1)%AoA_Beta                     = 0.0                ! [deg]
  BC_Input(1)%Wall_Temp                    = -1.0               ! [K]
  BC_Input(1)%Rotation_Rate(1:3)           = [0.0,0.0,0.0]      ! [rpm]
  BC_Input(1)%Is_Specified_In_MRF          = F
  BC_Input(1)%Mass_Frac(1:7)               = [7*0.0]
  BC_Input(1)%Species_Name(1:7)            = [7*""]
  BC_Input(1)%Relax%Value                  = 1.0
  BC_Input(1)%Relax%Time                   = 0.0                ! non-dim
  BC_Input(1)%Relax%Iter                   = 0
  BC_Input(1)%SEM%Profile_File             = ""
  BC_Input(1)%SEM%Is_Cylindrical           = F
  BC_Input(1)%SEM%Length_Scale             = 0.0                ! [m]
  BC_Input(1)%SEM%Convective_Velocity      = 0.0                ! [m/s]
  BC_Input(1)%SEM%Profile_Height           = 0.0                ! [m]
  BC_Input(1)%SEM%Sub_In_BC_Grp            = 0
  BC_Input(1)%SEM%Sub_In_BC_Idx            = 0
  BC_Input(1)%Periodic%Matching_BC_Name    = ""
  BC_Input(1)%Periodic%Matching_BC_Idx     = 0
  BC_Input(1)%Periodic%RotationCenter(1:3) = [0.0,0.0,0.0]      ! grid units
  BC_Input(1)%Periodic%RotationAngle(1:3)  = [0.0,0.0,0.0]      ! [deg]
  BC_Input(1)%Periodic%RotationAxis(1:4)   = [0.0,0.0,0.0,0.0]
  BC_Input(1)%Periodic%Translation(1:3)    = [0.0,0.0,0.0]      ! grid units
/
\end{nmlblocklst}
\inpBC*[nmlBC::BCInput]{}

%% These nml_bc entries were used to test \texttt and hyperlinks
%% within the namelist block. To make the standard lstlisting match
%% the \texttt font, must use ',basewidth={0.50em,0.45em}'.
%% An extra space or two before the '=' is gobbled by Latex escape \@,
%% which might affect alignment of '='.
%%
%  !@\texttt{BC\_Input(1)\%Periodic\%Translation(1:3)}!@      = [0.0,0.0,0.0]
%  !@\BCinpBC[nmlBC::BCInput/Periodic/Translation][1]{Periodic\%Translation(1:3)}!@    = [0.0,0.0,0.0]

%==========================================================================

\begin{nmldescitem}
Walls_Are_Exact = T
\end{nmldescitem}
\inp*[nmlBC::WallsAreExact]{}[Walls\_Are\_Exact][]

\begin{nmldescbody}
Specifies whether the wall boundary conditions are set
directly using the exact flow conditions or
indirectly using the ghost state that provides
the exact flow conditions when averaged with the interior state.
\\
\begin{tabular}{ @{\qquad} r @{ = } p{0.75\linewidth} @{} }
  T & The wall boundary solution is the exact wall boundary condition,
      \forexample
      $\left[ u_{\textrm{wall}},
              v_{\textrm{wall}},
              w_{\textrm{wall}}
       \right] = 0$ for a no-slip wall. \\
  F & The average of the wall boundary solution
      and the interior solution on the boundary face
      recovers the exact wall boundary condition,
      \forexample
      $\left[ u_{\textrm{wall}},
              v_{\textrm{wall}},
              w_{\textrm{wall}}
       \right] =
      -\left[ u_{\textrm{wall}},
              v_{\textrm{wall}},
              w_{\textrm{wall}}
       \right]$ for a no-slip wall.
\end{tabular}

\NOTE
Hartmann and Leicht\cite{Hartmann_2013} report (Remark 6.13 on page 83) that
using exact wall boundary conditions seems to be more accurate on coarse grids
and for lower values of
\linebreak[4]% Prevent next word from entering right margin.
\inp[nmlNumerics::SolutionOrder]{Solution\_Order},
whereas using the averaging method is more stable.
%
%Details can be found in \linecite{Hartmann_2013}.
%{\small \url{http://elib.dlr.de/90967/1/hartmann_leicht_VKI_LS_2014-3.pdf}}

%\FIXME{
%\\
%This description could be worded better.\\
%Is this setting the external BC state or the actual face point values?\\
%For each case above, at what location is
%      $\left[ u_{\textrm{wall}},
%              v_{\textrm{wall}},
%              w_{\textrm{wall}}
%       \right]$
%specified? Is it a ghost cell?\\
%}
\end{nmldescbody}

%==========================================================================

\begin{nmldescitem}
Sub_Inflow_Method = 4
\end{nmldescitem}
\inp*[nmlBC::SubInflowMethod]{}[Sub\_Inflow\_Method][]

\begin{nmldescbody}
Specifies which subsonic inflow boundary condition algorithm to use.
\\
\begin{tabular}{ @{\qquad} r @{ = } p{0.85\linewidth} @{} }
  1 & Hold the total pressure and total temperature constant at the inflow,
      use the outgoing characteristic from the interior, and
      perform a Newton iteration to get
      the static temperature, speed of sound, and normal velocity magnitude
      for the exterior state. \\
  2 & Use the BC static density and velocity, and the interior static
      pressure. \\
  3 & Use the method specified in
      \linecite{Carlson_2011}. % NASA TM-2011-217181.
      The interior static pressure, total enthalpy, and
      outward propagating Riemann invariant are extrapolated
      from the interior to the boundary.
      %
      These values with the BC total pressure and total temperature
      are used to compute the BC speed of sound
      by solving a quadratic equation based on adiabatic and
      isentropic assumptions\slash relations.\\
  4 & Use the BC total pressure, total temperature, angle of attack,
      and the interior state to get the interior Mach number at the boundary.
\end{tabular}
\\
See \sref{SEC:BC:Inflow:Subsonic} for a detailed description of how each of
these methods is implemented.

\NOTE
It is recommended to use options 4, 3, or 1 (in that order).
\end{nmldescbody}

%==========================================================================

\begin{nmldescitem}
Use_BC_Conflicts = T
\end{nmldescitem}
\inp*[nmlBC::UseBCConflicts]{}[Use\_BC\_Conflicts][]

\begin{nmldescbody}
Logical flag for determining whether or not to overwrite boundary solutions
in adjacent cells if co-located boundary points have different boundary
conditions.
%
The boundary preference hierarchy can be found in \file{kind\_types\_mod.f90} in
the array \file{bc\_hierarchy}.

\NOTE
This is only relevant if using Legendre-Gauss-Lobatto points
for solution points
($\inp[nmlNumerics::LocSolutionPts]{Loc\_Solution\_Pts}=2$).
\end{nmldescbody}

%==========================================================================

\begin{nmldescitem}
BC_Input(1)%Name = ""
\end{nmldescitem}
\inpBC*[nmlBC::BCInput/Name]{}[Name][]

\begin{nmldescbody}
Character string that matches the name of a boundary group in the grid file.
%
See \sref{SEC:Files:Grid_Input:CGNS:PWBC} for using Pointwise
to set boundary conditions for CGNS files.
%
This input is the means of connecting boundary faces in the grid file with a
boundary condition provided in the input file to specify the boundary type and
flow conditions to use.

\NOTE
The character string is treated as case insensitive.
\end{nmldescbody}

%==========================================================================

\FINALFIX{\newpage}
\begin{nmldescitem}
BC_Input(1)%Set_BC_Value = -9999
\end{nmldescitem}
\inpBC*[nmlBC::BCInput/SetBCValue]{}[Set\_BC\_Value][]

\begin{nmldescbody}
This overrides the boundary condition type specified in the grid file,
using this boundary condition type for the boundary group given by
\BCinptext[1]{Name}.
%
Valid values for boundary condition types are:\\ %[0.25em]

\begingroup
\NewDocumentCommand{\bcgroup}{m}{\textbf{#1}}
%
%% Use threeparttable to add table footnote.
\begin{threeparttable}
\begin{tabular}{ @{} r @{ \quad } p{0.845\linewidth} @{} }
%% Use longtable to permit page breaks.
%\begin{longtable}{ @{} r @{ \quad } p{0.845\linewidth} @{} }
        & \bcgroup{Inflow boundary conditions} \\*
      1 & Generic inflow, subsonic\slash supersonic  \\*
      2 & Subsonic inflow \\*
      3 & Supersonic inflow \\*
      5 & SEM (unsteady turbulent) inflow%
          \tnote{\ensuremath{\ast}} %
          \\[0.25em]
        & \bcgroup{Outflow boundary conditions} \\*
     11 & Generic outflow, subsonic\slash supersonic \\*
     12 & Subsonic outflow \\*
     13 & Supersonic outflow \\[0.25em]
        & \bcgroup{Farfield boundary conditions} \\*
     21 & Characteristic inflow\slash outflow based on Riemann invariants,
          subsonic\slash supersonic%
          \tnote{\ensuremath{\ast}} \\*
     22 & Generic inflow\slash outflow, subsonic\slash supersonic \\*
     23 & Fixed to freestream\slash reference flow conditions \\*
     24 & Fixed flow conditions (equivalent to supersonic inflow) \\[0.25em]
        & \bcgroup{Wall boundary conditions} \\*
     50 & Slip wall \\*
     52 & Adiabatic no-slip wall \\*
     53 & Isothermal no-slip wall \\[0.25em]
        & \bcgroup{Numerical boundary conditions} \\*
     80 & Symmetry boundary condition (equivalent to slip wall) \\*
     81 & Periodic (translated \andor rotated) coupled boundaries \\[0.25em]
        & \bcgroup{Other boundary conditions} \\*
  -9999 & Do not override the boundary condition type
          from the grid file. \\*
  \bottomrule
%\end{longtable}
\end{tabular}
  \begin{tablenotes}
    \item[\ensuremath{\ast}] Not compatible with multi-component gas mixtures.
  \end{tablenotes}
\end{threeparttable}
\endgroup
\end{nmldescbody}

%==========================================================================

\begin{nmldescitem}
BC_Input(1)%Temperature = -1.0  ! [K]
\end{nmldescitem}
\inpBC*[nmlBC::BCInput/Temperature]{}[Temperature][]

\begin{nmldescbody}
Static temperature (units of \unitTemp{}).
\end{nmldescbody}

%==========================================================================

\begin{nmldescitem}
BC_Input(1)%Pressure = -1.0  ! [Pa]
\end{nmldescitem}
\inpBC*[nmlBC::BCInput/Pressure]{}[Pressure][]

\begin{nmldescbody}
Static pressure (units of \unitPres{}).
\end{nmldescbody}

%==========================================================================

\begin{nmldescitem}
BC_Input(1)%Density = -1.0  ! [kg/m^3]
\end{nmldescitem}
\inpBC*[nmlBC::BCInput/Density]{}[Density][]

\begin{nmldescbody}
Static density (units of \unitDens{}).
\end{nmldescbody}

%==========================================================================

\begin{nmldescitem}
BC_Input(1)%Mach = -1.0
\end{nmldescitem}
\inpBC*[nmlBC::BCInput/Mach]{}[Mach][]

\begin{nmldescbody}
Mach number (dimensionless).
\end{nmldescbody}

%==========================================================================

\begin{nmldescitem}
BC_Input(1)%VelocityX = 0.0  ! [m/s]
\end{nmldescitem}
\inpBC*[nmlBC::BCInput/VelocityX]{}[VelocityX][]

\begin{nmldescbody}
Velocity in the $x$ coordinate direction (units of \unitVel{}).
\end{nmldescbody}

%==========================================================================

\begin{nmldescitem}
BC_Input(1)%VelocityY = 0.0  ! [m/s]
\end{nmldescitem}
\inpBC*[nmlBC::BCInput/VelocityY]{}[VelocityY][]

\begin{nmldescbody}
Velocity in the $y$ coordinate direction (units of \unitVel{}).
\end{nmldescbody}

%==========================================================================

\begin{nmldescitem}
BC_Input(1)%VelocityZ = 0.0  ! [m/s]
\end{nmldescitem}
\inpBC*[nmlBC::BCInput/VelocityZ]{}[VelocityZ][]

\begin{nmldescbody}
Velocity in the $z$ coordinate direction (units of \unitVel{}).
\end{nmldescbody}

%==========================================================================

\begin{nmldescitem}
BC_Input(1)%Total_Temperature = -1.0  ! [K]
\end{nmldescitem}
\inpBC*[nmlBC::BCInput/TotalTemperature]{}[Total\_Temperature][]

\begin{nmldescbody}
Total temperature (units of \unitTemp{}).
\end{nmldescbody}

%==========================================================================

\begin{nmldescitem}
BC_Input(1)%Total_Pressure = -1.0  ! [Pa]
\end{nmldescitem}
\inpBC*[nmlBC::BCInput/TotalPressure]{}[Total\_Pressure][]

\begin{nmldescbody}
Total pressure (units of \unitPres).
\end{nmldescbody}

%==========================================================================

\begin{nmldescitem}
BC_Input(1)%Total_Density = -1.0  ! [kg/m^3]
\end{nmldescitem}
\inpBC*[nmlBC::BCInput/TotalDensity]{}[Total\_Density][]

\begin{nmldescbody}
Total density (units of \unitDens{}).
\end{nmldescbody}

%==========================================================================

\begin{nmldescitem}
BC_Input(1)%AoA_Alpha = 0.0  ! [deg]
BC_Input(1)%AoA_Beta  = 0.0  ! [deg]
\end{nmldescitem}
\inpBC*[nmlBC::BCInput/AoAAlpha]{}[AoA\_Alpha][]
\inpBC*[nmlBC::BCInput/AoABeta]{}[AoA\_Beta][]

\begin{nmldescbody}
Flow angles (in degrees):
$\alpha$ is the
\key[key:AngleOfAttack]{angle of attack}[flow angles!angle of attack]
measured in the $x$-$y$ plane, and
$\beta$  is the
\key[key:AngleOfSideslip]{angle of sideslip}[flow angles!angle of sideslip]
measured in the $x$-$z$ plane.


%
These flow angles are defined by
\eref{EQN:FlowAngleVelX,EQN:FlowAngleVelY,EQN:FlowAngleVelZ} and
\fref{FIG:FlowAngle} on \pref{\pageref{FIG:FlowAngle}}.
\end{nmldescbody}

%==========================================================================

\begin{nmldescitem}
BC_Input(1)%Wall_Temp = -1.0  ! [K]
\end{nmldescitem}
\inpBC*[nmlBC::BCInput/WallTemp]{}[Wall\_Temp][]

\begin{nmldescbody}
Temperature for isothermal wall (units of \unitTemp{}).
\end{nmldescbody}

%==========================================================================

\begin{nmldescitem}
BC_Input(1)%Rotation_Rate(1:3) = [0.0,0.0,0.0]  ! [rpm] @(x,y,z)-axes
\end{nmldescitem}
\inpBC*[nmlBC::BCInput/RotationRate]{}[Rotation\_Rate][]

\begin{nmldescbody}
This specifies the rotation rate vector of a moving boundary surface.
%
The components are defined with respect to the Cartesian axes,
and the magnitude is the rotation rate (in RPMs).
\end{nmldescbody}

%==========================================================================

\begin{nmldescitem}
BC_Input(1)%Is_Specified_In_MRF = F
\end{nmldescitem}
\inpBC*[nmlBC::BCInput/IsSpecifiedInMRF]{}[Is\_Specified\_In\_MRF][]

\begin{nmldescbody}
Setting this to true indicates that this boundary condition is specified in
the moving reference frame. If the boundary is moving\slash rotating at the
same speed as the moving reference frame, this works as a shortcut to apply
the same motion to this boundary.
\end{nmldescbody}

%==========================================================================

\begin{nmldescitem}
BC_Input(1)%Mass_Frac(1:7) = [7*0.0]
\end{nmldescitem}
\inpBC*[nmlBC::BCInput/MassFrac]{}[Mass\_Frac][]

\begin{nmldescbody}
\raggedright % Avoid next word spilling in to right margin.
\BCinptext{Mass\_Frac(n)} specifies the mass fraction
of the species given in \BCinptext{Species\_Name(n)}.
\end{nmldescbody}

%==========================================================================

\begin{nmldescitem}
BC_Input(1)%Species_Name(1:7) = [7*""]
\end{nmldescitem}
\inpBC*[nmlBC::BCInput/SpeciesName]{}[Species\_Name][]

\begin{nmldescbody}
Specifies the chemical name\slash formula for each of the species.
\end{nmldescbody}

%==========================================================================

\subsubsection{Boundary Condition Relaxation Inputs}
\begin{nmldescitem}
BC_Input(1)%Relax%Value = 1.0
BC_Input(1)%Relax%Time  = 0.0  ! non-dim, t[s] * (a[m/s] / L=1.0[m])
BC_Input(1)%Relax%Iter  = 0
\end{nmldescitem}
\inpBC*[nmlBC::BCInput/Relax/Value]{}[Relax\%Value][]
\inpBC*[nmlBC::BCInput/Relax/Time]{}[Relax\%Time][]
\inpBC*[nmlBC::BCInput/Relax/Iter]{}[Relax\%Iter][]

\begin{nmldescbody}
Typically the volume solution is initialized from the
reference conditions specified in the
\\% Prevent next word from entering right margin.
\inp[nmlReference]{nml\_reference} namelist.
%
If boundary conditions are immediately applied, the sudden change
can introduce numerical instabilities.
%
These relaxation parameters provide a way to gradually apply
each boundary condition.

\BCinptext{Relax\%Value} specifies
the relaxation value (between 0.0 and 1.0) to limit the difference
between the interior and exterior states.
%
The default value of 1.0 does not impose any limitation.

There are two means for specifying when to end relaxation:
one based on solution time, the other based on solution iterations.
%
Relaxation will end when either criteria is met.

If $\BCinptext{Relax\%Time}>0.0$, it specifies
the total accumulated solution time
(nondimensionalized by the reference length of 1.0 meter and speed of sound)
to end relaxation.

If $\BCinptext{Relax\%Iter}>0$, it specifies
the total accumulated time step iteration to end relaxation.
\end{nmldescbody}


%==========================================================================

\subsubsection{Boundary Condition SEM Inputs}
\begin{nmldescitem}
BC_Input(1)%SEM%Profile_File = ""  ! 150 char max
\end{nmldescitem}
\inpBC*[nmlBC::BCInput/SEM/ProfileFile]{}[SEM\%Profile\_File][]

\begin{nmldescbody}
This gives the path to the file containing the flow profile used for the
synthetic eddy method (SEM) boundary condition.
%
See \sref{SEC:SEM:usage:profile} for a detailed description of the
required file format and flow variables needed for the SEM profile.
\end{nmldescbody}

%==========================================================================

\begin{nmldescitem}
BC_Input(1)%SEM%Is_Cylindrical = F
\end{nmldescitem}
\inpBC*[nmlBC::BCInput/SEM/IsCylindrical]{}[SEM\%Is\_Cylindrical][]

\begin{nmldescbody}
This is used to declare if the synthetic eddy method (SEM) box is cylindrical.
%
This solely affects how the eddies are randomly generated to achieve a
truly random distribution in the radial, $r$, direction.
%
For a cylindrical SEM box, randomly generating
  $r\ (r_{\textrm{min}}\le r \le r_{\textrm{max}})$ and
  $\theta\ (\theta_{\textrm{min}}\le \theta \le \theta_{\textrm{max}})$
coordinate values to place a new eddy will result in
more eddies near the symmetry axis and
less eddies further away.
%
In order to generate a truly uniform random distribution,
we need to use the equation
  $r=\sqrt{\mathtt{Random}()\left[ r^2_{\textrm{max}}
                                 - r^2_{\textrm{min}}
                            \right] + r^2_{\textrm{min}}}$
where $\mathtt{Random}()$ is a function that generates a random number
on the unit interval.
%
This applies an area weighting to the radial distribution and
results in a uniform distribution of eddies in the volume.

\end{nmldescbody}

%==========================================================================

\begin{nmldescitem}
BC_Input(1)%SEM%Length_Scale = 0.0  ! [m]
\end{nmldescitem}
\inpBC*[nmlBC::BCInput/SEM/LengthScale]{}[SEM\%Length\_Scale][]

\begin{nmldescbody}
This gives the length scale (in meters) of the ``zone of influence'' of each
synthetic eddy.
%
Recommended values are between $(1/3.5)\times$ and $(1/5)\times$ the
height of the boundary layer profile from the SEM profile file.
%
See \sref{SEC:SEM:usage:recommendations}
for a more detailed description of this input and how it relates to
\inpBC[nmlBC::BCInput/SEM/ConvectiveVelocity]{SEM\%Convective\_Velocity}.

\NOTE
Previous testing seems to suggest that a slower convective velocity
should be used when using a smaller length scale, and a higher convective
velocity should be used when using a larger length scale.
\end{nmldescbody}

%==========================================================================

\FINALFIX{\newpage}
\begin{nmldescitem}
BC_Input(1)%SEM%Convective_Velocity = 0.0  ! [m/s]
\end{nmldescitem}
\inpBC*[nmlBC::BCInput/SEM/ConvectiveVelocity]{}[SEM\%Convective\_Velocity][]

\begin{nmldescbody}
This gives the velocity (in m\slash s) at which the synthetic eddies are
convected through the SEM box.
%
Recommended values are between $0.25\times$ and $1\times$ the
freestream velocity for this BC.
%
See \sref{SEC:SEM:usage:recommendations}
for a more detailed description of this input and how it relates to
\\ % Prevent next word from entering right margin.
\inpBC[nmlBC::BCInput/SEM/LengthScale]{SEM\%Length\_Scale}.

\NOTE
Previous testing seems to suggest that a slower convective velocity
should be used when using a smaller length scale, and a higher convective
velocity should be used when using a larger length scale.
\end{nmldescbody}

%==========================================================================

\begin{nmldescitem}
BC_Input(1)%SEM%Profile_Height = 0.0  ! [m]
\end{nmldescitem}
\inpBC*[nmlBC::BCInput/SEM/ProfileHeight]{}[SEM\%Profile\_Height][]

\begin{nmldescbody}
\NOTE
\emph{This variable is not an input.}
%
It will be calculated based on the profile data read from
\inpBC[nmlBC::BCInput/SEM/ProfileFile]{SEM\%Profile\_File},
and will be reported with the boundary condition information
written to the display output.
%
This output can be used to verify that the profile data was
properly interpreted.

\end{nmldescbody}

%==========================================================================

\begin{nmldescitem}
BC_Input(1)%SEM%Sub_In_BC_Grp = 0
BC_Input(1)%SEM%Sub_In_BC_Idx = 0
\end{nmldescitem}
\inpBC*[nmlBC::BCInput/SEM/SubInBCGrp]{}[SEM\%Sub\_In\_BC\_Grp][]
\inpBC*[nmlBC::BCInput/SEM/SubInBCIdx]{}[SEM\%Sub\_In\_BC\_Idx][]

\begin{nmldescbody}
\NOTE
\emph{Theses variables are not inputs.}
%
They will be determined after setting up the SEM boundary condition
and will be reported with the boundary condition information
written to the display output.
%
This output can be used to verify that the profile data was
properly interpreted.

If the boundary layer profile specified in the
\inpBC[nmlBC::BCInput/SEM/ProfileFile]{SEM\%Profile\_File}
does not fully span the boundary domain, then
a new inflow BC will automatically be created for the boundary points
that are outside the SEM profile.
% 
\inptext{Sub\_In\_BC\_Grp} identifies the \emph{index of the BC group}
to which the new inflow BC created for this SEM belongs.
%
\inptext{Sub\_In\_BC\_Idx} identifies the \emph{index for the new inflow BC}
created for this SEM.

\NOTE
Inflow boundaries are automatically split in this way because
the synthetic eddy method of imposing inflow turublence
actually overprescribes flow conditions on subsonic boundaries.
%
Separating points outside the boundary layer profile into a
standard inflow boundary allows those points to properly communicate
information upstream as needed.
%
Therefore, we recommend that the SEM profile file not extend much
beyond the edge of the boundary layer.
%
See \sref{SEC:SEM:usage:subsonic} for a more detailed explanation.

\end{nmldescbody}

%==========================================================================

\subsubsection{Boundary Condition Periodicity Inputs}
\begin{nmldescitem}
BC_Input(1)%Periodic%Matching_BC_Name = ""
\end{nmldescitem}
\inpBC*[nmlBC::BCInput/Periodic/MatchingBCName]{}[Periodic\%Matching\_BC\_Name][]

\begin{nmldescbody}
This specifies the name of the other boundary condition that is a matching
periodic pair with this boundary condition.
%
If BC's $i$ and $j$ are a periodic pair, then
the following two relations need to be true:
\\[0.5em] % A little extra separation before tabular.
\begin{tabular}{ @{\qquad} l } % @{ \quad } p{0.85\linewidth} @{} }
  \BCinptext[i]{Periodic\%Matching\_BC\_Name} =
  \BCinptext[j]{Name}, and
\\
  \BCinptext[j]{Periodic\%Matching\_BC\_Name} =
  \BCinptext[i]{Name}.
\end{tabular}

\vspace{0.5em} % A little extra separation after tabular.
\NOTE
This and
\inpBC[nmlBC::BCInput/Periodic/MatchingBCIdx]{Periodic\%Matching\_BC\_Idx}
are redundant, but together they serve as a means of error checking the
inputs for periodic pairs.
\end{nmldescbody}

%==========================================================================

\FINALFIX{\newpage}
\begin{nmldescitem}
BC_Input(1)%Periodic%Matching_BC_Idx = 0
\end{nmldescitem}
\inpBC*[nmlBC::BCInput/Periodic/MatchingBCIdx]{}[Periodic\%Matching\_BC\_Idx][]

\begin{nmldescbody}
This specifies the index of the other boundary condition that is a matching
periodic pair with this boundary condition.
%
If BC's $i$ and $j$ are a periodic pair, then
the following two relations need to be true:
\\[0.5em] % A little extra separation before tabular.
\begin{tabular}{ @{\qquad} l } % @{ \quad } p{0.85\linewidth} @{} }
  \BCinptext[i]{Periodic\%Matching\_BC\_Idx} = j, and \\
  \BCinptext[j]{Periodic\%Matching\_BC\_Idx} = i.
\end{tabular}

\vspace{0.5em} % A little extra separation after tabular.
\NOTE
This and
\inpBC[nmlBC::BCInput/Periodic/MatchingBCName]{Periodic\%Matching\_BC\_Name}
are redundant, but together they serve as a means of error checking the
inputs for periodic pairs.
\end{nmldescbody}

%==========================================================================

\begin{nmldescitem}
BC_Input(1)%Periodic%RotationCenter(1:3) = [0.0,0.0,0.0]  ! [grid units]
\end{nmldescitem}
\inpBC*[nmlBC::BCInput/Periodic/RotationCenter]{}[Periodic\%RotationCenter][]

\begin{nmldescbody}
This array gives the coordinate origin
(in \key[key:GridUnits]{grid units}) used for
defining the rotation angles between rotated periodic boundary conditions.

\NOTE
This input should be the same for a matching pair of rotated periodic
boundary conditions.

\end{nmldescbody}

%==========================================================================

\begin{nmldescitem}
BC_Input(1)%Periodic%RotationAngle(1:3) = [0.0,0.0,0.0]  ! [deg]
\end{nmldescitem}
\inpBC*[nmlBC::BCInput/Periodic/RotationAngle]{}[Periodic\%RotationAngle][]

\begin{nmldescbody}
This array gives the rotation angles (in degrees) about each of the
coordinate axes from the current periodic boundary interface to the
connecting periodic boundary interface.
%
If rotating about more than one axis, the rotation is performed
first about the $x$-axis, then the $y$-axis, and finally the $z$-axis.

\NOTE
One may use either
\inptext{Periodic\%RotationAngle} or  % Short Form
%\BCinptext{Periodic\%RotationAngle} or  % Long Form
%\\ % Prevent next word from entering right margin. Needed with Long Form.
\inpBC[nmlBC::BCInput/Periodic/RotationAxis]{Periodic\%RotationAxis}
to define rotated periodic boundary conditions, but not both.
%
If both are specified, \inptext{RotationAxis} will be ignored and
\inptext{RotationAngle} will be used.

\NOTE If BC's $i$ and $j$ are a periodic pair, the following relation needs
to be true:
\\[0.5em] % A little extra separation before tabular.
\begin{tabular}{ @{\qquad} >{\small}l }
   \BCinptext[i]{Periodic\%RotationAngle(:)} = 
  -\BCinptext[j]{Periodic\%RotationAngle(:)}
\end{tabular}
\end{nmldescbody}

%==========================================================================

\begin{nmldescitem}
BC_Input(1)%Periodic%RotationAxis(1:4) = [0.0,0.0,0.0,0.0]
\end{nmldescitem}
\inpBC*[nmlBC::BCInput/Periodic/RotationAxis]{}[Periodic\%RotationAxis][]

\begin{nmldescbody}
The first 3 values of this array define a vector representing the axis of
rotation for this rotated periodic boundary condition, and
the last value of the array defines the rotation angle (in degrees) about
this axis.

\NOTE
The coordinate origin is used as the base of the vector defining the
axis of rotation, NOT
\inpBC[nmlBC::BCInput/Periodic/RotationCenter]{Periodic\%RotationCenter}.

\NOTE
The magnitude of the vector defining the axis of rotation is not used.
The sole purpose of this vector is to define the \emph{direction} of the
rotation axis.

\NOTE
One may use either
\inpBC[nmlBC::BCInput/Periodic/RotationAngle]{Periodic\%RotationAngle} or
%\\ % Prevent next word from entering right margin. Needed for Long Form.
%\BCinptext{Periodic\%RotationAxis} % Long Form
\inptext{Periodic\%RotationAxis} % Short Form
to define rotated periodic boundary conditions, but not both.
%
If both are specified, \inptext{RotationAxis} will be ignored and
\inptext{RotationAngle} will be used.

\NOTE
If BC's $i$ and $j$ are a periodic pair, one of the two following
sets of relations needs to be true:
\\[0.5em] % A little extra separation before tabular.
{\small
\begin{tabular}{ @{\qquad} l @{ = } l @{} l}
     \BCinptext[i]{Periodic\%RotationAxis(1:3)}
  & &\BCinptext[j]{Periodic\%RotationAxis(1:3)}
     \\
     \BCinptext[i]{Periodic\%RotationAxis(4)}
  &-&\BCinptext[j]{Periodic\%RotationAxis(4)}
\end{tabular}
}
\\[0.5em] % A little extra separation after tabular.
or
\\[0.5em] % A little extra separation before tabular.
{\small
\begin{tabular}{ @{\qquad} l @{ = } l @{} l}
     \BCinptext[i]{Periodic\%RotationAxis(1:3)}
  &-&\BCinptext[j]{Periodic\%RotationAxis(1:3)}
     \\
     \BCinptext[i]{Periodic\%RotationAxis(4)}
  & &\BCinptext[j]{Periodic\%RotationAxis(4)}
\end{tabular}
}
\\[0.5em] % A little extra separation after tabular.
In other words, the opposing periodic face must either negate the
axis direction or the rotation angle, but not both.
\end{nmldescbody}

%==========================================================================

\begin{nmldescitem}
BC_Input(1)%Periodic%Translation(1:3) = [0.0,0.0,0.0]  ! [grid units]
\end{nmldescitem}
\inpBC*[nmlBC::BCInput/Periodic/Translation]{}[Periodic\%Translation][]

\begin{nmldescbody}
This array gives the translation vector
(in \key[key:GridUnits]{grid units}) from the current
periodic boundary interface to the connecting periodic boundary interface.

\NOTE
If BC's $i$ and $j$ are a periodic pair, the following relation needs
to be true:
\\[0.5em] % A little extra separation before tabular.
\begin{tabular}{ @{\qquad} l }
     \BCinptext[i]{Periodic\%Translation(:)} = 
    -\BCinptext[j]{Periodic\%Translation(:)}
\end{tabular}
\\[0.5em] % A little extra separation after tabular.
In other words, the opposing periodic face must negate the translation
vector.
\end{nmldescbody}

%==========================================================================


\subsection{nml\_numerics}
\label{SEC:nml_numerics}
\inp*[nmlNumerics]{}[nml\_numerics][]

This \hyperref[SEC:User_Input]{namelist block}
is used to specify parameters for the numerical method.

\begin{nmlblocklst}[
,label={LST:nml_numerics}
,caption={Namelist nml\_numerics.}
]
&nml_numerics
  All_Orders                      = -3
  Solution_Order                  =  3
  Projection_Order                = -3
  Quadrature_Order                = -3
  Error_Order                     = -3
  All_Points                      = -1
  Loc_Solution_Pts                =  1
  Loc_Flux_Pts                    =  1
  Loc_Quadrature_Pts              =  1
  Loc_Error_Pts                   =  1
  Loc_Triangle_Pts                =  1
  Loc_Tetrahedra_Pts              =  1
  Invs_Flux_Method                =  1
  Visc_Flux_Method                =  2
  Visc_Variables_To_Ave           =  1
  Combine_Common_Flux_Projections =  0
  Correction_Function             =  1
/
\end{nmlblocklst}

%==========================================================================

\begin{nmldescitem}
All_Orders = -3
\end{nmldescitem}
\inp*[nmlNumerics::AllOrders]{}[All\_Orders][]

\begin{nmldescbody}
If $>0$, sets
\inp[nmlNumerics::SolutionOrder]{Solution\_Order},
\inp[nmlNumerics::ProjectionOrder]{Projection\_Order}, and
\inp[nmlNumerics::QuadratureOrder]{Quadrature\_Order}
all to this value.

\NOTE
This overrides the values given for these three input parameters.
\end{nmldescbody}

%==========================================================================

\begin{nmldescitem}
Solution_Order = 3
\end{nmldescitem}
\inp*[nmlNumerics::SolutionOrder]{}[Solution\_Order][]

\begin{nmldescbody}
Specifies the degree of the polynomial space, \p{S}, to use for the 
solution. 
% 
The order-of-accuracy is approximately \p[+1]{S}.

\NOTE
If $\inp[nmlNumerics::AllOrders]{All\_Orders} > 0$, then
that value will override this input.
\end{nmldescbody}

%==========================================================================

\begin{nmldescitem}
Projection_Order = -3
\end{nmldescitem}
\inp*[nmlNumerics::ProjectionOrder]{}[Projection\_Order][]

\begin{nmldescbody}
Specifies the degree of the polynomial space, \p{P}, to project the residual
onto when using over-integration by residual projection.
%
The residual will be computed as a polynomial of degree \p{S} and projected down
to a polynomial of degree \p{P}. This can potentially improve stability if
aliasing errors are suspected of destabilizing a simulation. 
%
If $\inptext{Projection\_Order} \leq 0$,
this will be set to the same value as \inptext{Solution\_Order}.

\NOTE
If $\inp[nmlNumerics::AllOrders]{All\_Orders} > 0$, then
that value will override this input.
 
\NOTE
This is only used when
$0<\inptext{Projection\_Order}<\inptext{Solution\_Order}$,
and using this means the order-of-accuracy of the simulation is
now approximately \p[+1]{P} instead of \p[+1]{S}.
\end{nmldescbody}

%==========================================================================

\FINALFIX{\newpage}
\begin{nmldescitem}
Quadrature_Order = -3
\end{nmldescitem}
\inp*[nmlNumerics::QuadratureOrder]{}[Quadrature\_Order][]

\begin{nmldescbody}
Specifies the degree of the quadrature polynomial space, \p{Q}, to use when
using over-integration by over-sampling at the quadrature points.
%
If $\inptext{Quadrature\_Order} \leq 0$,
this will be set to the same value as \inptext{Solution\_Order}.

Certain operations during a given residual evaluation
(\forexample evaluating the common\slash upwind interface fluxes)
produce either polynomials that are not within the solution polynomial space
(\forexample of degree greater than the solution order)
or functions that are not polynomials at all (\forexample pressure).
%
By setting
  $\inptext{Quadrature\_Order} > \inptext{Solution\_Order}$,
these functions will be over-sampled at the quadrature points to produce
approximations of these functions as polynomials of degree \p{Q},
which are in turn projected down to polynomials of degree \p{S}
before adding their respective contribution to the residual.
%
This polynomial representation of degree \p{S} is generally more
accurate than directly evaluating the function at the
\p{S} polynomial points, and
can potentially improve stability if aliasing errors are suspected
of destabilizing a simulation.
%
See \sref{SEC:Aliasing} for an example of how \key{aliasing error}[=] arises
and the process of using over-integration by over-sampling at the
quadrature points.

\NOTE
If $\inp[nmlNumerics::AllOrders]{All\_Orders} > 0$, then
that value will have first priority.
 
\NOTE
If $\inp[nmlNumerics::SolutionOrder]{Solution\_Order} > 
\inptext{Quadrature\_Order}$, then \inptext{Quadrature\_Order}
will be set to \inptext{Solution\_Order}.
 
\NOTE
This quadrature option is only used when
  \inptext{Quadrature\_Order} $>$ \inptext{Solution\_Order}
or
  \inp[nmlNumerics::LocQuadraturePts]{Loc\_Quadrature\_Pts} $\neq$
  \inp[nmlNumerics::LocSolutionPts]{Loc\_Solution\_Pts},
otherwise the quadrature is evaluated directly at the solution points
without the need for projection.
% 
In all cases,
$\inptext{Quadrature\_Order} \geq \inptext{Solution\_Order}$
because using a lower quadrature order provides poor results.
% 
Unlike over-integration by residual projection,
the order-of-accuracy of the simulation remains \p[+1]{S}.

\NOTE
This quadrature option has not been working since the viscous and
unstructured capability was added to the code.
It should be fixed in an upcoming version.

\IDEA{
Perhaps \sref{SEC:Numerics} needs a detailed discussion of the
difference between over-integration by residual projection
and the over-integration by over-sampling at the quadrature points
used here?
}
\end{nmldescbody}

%==========================================================================

\begin{nmldescitem}
Error_Order = -3
\end{nmldescitem}
\inp*[nmlNumerics::ErrorOrder]{}[Error\_Order][]

\begin{nmldescbody}
Specifies the degree of the polynomial space to use when computing
integrated error quantities, \forexample the $L^2$-norm of the residual.
%
If
$\inptext{Error\_Order} < \inptext{Solution\_Order}$,
this will be set to the same value as \inptext{Solution\_Order}.

\NOTE
This input is only used when
  $\inptext{Error\_Order} > \inptext{Solution\_Order}$
or
  $\inp[nmlNumerics::LocErrorPts]{Loc\_Error\_Pts} \neq
   \inp[nmlNumerics::LocSolutionPts]{Loc\_Solution\_Pts}$.

\end{nmldescbody}

%==========================================================================

\begin{nmldescitem}
All_Points = -1
\end{nmldescitem}
\inp*[nmlNumerics::AllPoints]{}[All\_Points][]

\begin{nmldescbody}
Valid values for this input are the same as those specified in
\inp[nmlNumerics::LocSolutionPts]{Loc\_Solution\_Pts}.
%
If set to one of those values, this will set
\inp[nmlNumerics::LocSolutionPts]{Loc\_Solution\_Pts},
\inp[nmlNumerics::LocFluxPts]{Loc\_Flux\_Pts}, and
\inp[nmlNumerics::LocQuadraturePts]{Loc\_Quadrature\_Pts}
all to this value;
otherwise, this input will be ignored and the three mentioned input parameters
are unchanged.
%

\NOTE
This will override the values given for the three mentioned input
parameters set to a valid value.
%
\end{nmldescbody}

%==========================================================================

\begin{nmldescitem}
Loc_Solution_Pts = 1
\end{nmldescitem}
\inp*[nmlNumerics::LocSolutionPts]{}[Loc\_Solution\_Pts][]

\begin{nmldescbody}
Specifies the location of the 1D \key{solution points}[points!solution]
that are used in tensor products to create higher dimensional sets of
solution points (used for edge, quad, pyramid, prism, and hex cell geometries):
\\
\begin{tabular}{ @{\qquad} r @{ = } p{0.85\linewidth} @{} }
  0 & \key{Equidistant points}[points!equidistant] (not recommended) \\
  1 & \key{Legendre-Gauss (Gauss) points}[points!Legendre-Gauss (Gauss)] (recommended) \\ 
  2 & \key{Legendre-Gauss-Lobatto (Lobatto) points}[points!Legendre-Gauss-Lobatto (Lobatto)]
\end{tabular}
\\
See \fref{FIG:1D_Points} in \sref{SEC:Intro:Solution_Points}
to better understand the differences in these point locations.

\NOTE

For structured elements (quadrilaterals and hexahedra),
a tensor product of these 1D solution points is used.

For pyramid elements, a special tensor product of these solution points
is used. The 1D solution point locations are used to space the points
vertically from the pyramid base to the peak. At each of these vertical
locations, a 2D tensor product of solution points is used, with the base
being a tensor product of degree \PO, the next vertical location \PO[-1],
with the peak being \PO[0] or a single point.

For prism elements, a tensor product of these 1D solution points and the
specified triangle points is used. 
\end{nmldescbody}

%==========================================================================

\begin{nmldescitem}
Loc_Flux_Pts = 1
\end{nmldescitem}
\inp*[nmlNumerics::LocFluxPts]{}[Loc\_Flux\_Pts][]

\begin{nmldescbody}
Specifies the location of the flux points for edge and quadrilateral faces.
\\
\begin{tabular}{ @{\qquad} r @{ = } p{0.85\linewidth} @{} }
  0 & \key{Equidistant points}[points!equidistant] (not recommended) \\
  1 & \key{Legendre-Gauss (Gauss) points}[points!Legendre-Gauss (Gauss)] (recommended) \\ 
  2 & \key{Legendre-Gauss-Lobatto (Lobatto) points}[points!Legendre-Gauss-Lobatto (Lobatto)]
\end{tabular}
\\
For triangle faces, \inp[nmlNumerics::LocTrianglePts]{Loc\_Triangle\_Pts}
is used for the flux points.

%It is generally recommended to set the location of the flux points to the same
%as the solution points.
%
%If using Legendre-Gauss-Lobatto solution points, using Legendre-Gauss flux
%points could potentially reduce aliasing instabilities at interfaces.

\NOTE
Using different types of points for the solution and flux points is untested;
therefore, it is recommended to use the input \inptext{All\_Points} to use the
same set of points for solution, flux, and quadrature points.
%
\end{nmldescbody}

%==========================================================================

\begin{nmldescitem}
Loc_Quadrature_Pts = 1
\end{nmldescitem}
\inp*[nmlNumerics::LocQuadraturePts]{}[Loc\_Quadrature\_Pts][]

\begin{nmldescbody}
Specifies the location of the \key{quadrature points}[points!quadrature]
used for over-integration by over-sampling at the quadrature points.
This input is treated identically to
\inp[nmlNumerics::LocSolutionPts]{Loc\_Solution\_Pts} except it uses
\inp[nmlNumerics::QuadratureOrder]{Quadrature\_Order}
to determine the number of quadrature points for each geometry type.
\end{nmldescbody}

%==========================================================================

\begin{nmldescitem}
Loc_Error_Pts = 1
\end{nmldescitem}
\inp*[nmlNumerics::LocErrorPts]{}[Loc\_Error\_Pts][]

\begin{nmldescbody}
Specifies the location of the error points used for computing integrated
error quantities. This input is treated identically to
\inp[nmlNumerics::LocSolutionPts]{Loc\_Solution\_Pts} except it uses
\inp[nmlNumerics::ErrorOrder]{Error\_Order}
to determine the number of error points for each geometry type.
\end{nmldescbody}

%==========================================================================

\begin{nmldescitem}
Loc_Triangle_Pts = 1
\end{nmldescitem}
\inp*[nmlNumerics::LocTrianglePts]{}[Loc\_Triangle\_Pts][]

\begin{nmldescbody}
Specifies the location of the solution points used for triangles.
\\
\begin{tabular}{ @{\qquad} r @{ = } p{0.85\linewidth} @{} }
  -2 & Hesthaven-Warburton $\alpha$-optimized points\cite{Hesthaven_2010} \\
   0 & Equally spaced points (not recommended) \\
   1 & Williams-Shunn points\cite{Williams_2013} (recommended) \\
   2 & Barycentric Lobatto points
\end{tabular}
\end{nmldescbody}

%==========================================================================

\FINALFIX{\newpage}
\begin{nmldescitem}
Loc_Tetrahedra_Pts = 1
\end{nmldescitem}
\inp*[nmlNumerics::LocTetrahedraPts]{}[Loc\_Tetrahedra\_Pts][]

\begin{nmldescbody}
Specifies the location of the solution points used for tetrahedra.
\\
\begin{tabular}{ @{\qquad} r @{ = } p{0.85\linewidth} @{} }
  -2 & Hesthaven-Warburton $\alpha$-optimized points\cite{Hesthaven_2010} \\
   0 & Equally spaced points (not recommended)\\
   1 & Shunn-Ham points\cite{Shunn_2012} (recommended) \\
   2 & Barycentric Lobatto points \\
\end{tabular}
\end{nmldescbody}

%==========================================================================

\begin{nmldescitem}
Invs_Flux_Method = 1
\end{nmldescitem}
\inp*[nmlNumerics::InvsFluxMethod]{}[Invs\_Flux\_Method][]

\begin{nmldescbody}
Specifies the approximate Riemann solver used to compute the common inviscid
fluxes at the interfaces.
\\
\begin{tabular}{ @{\qquad} r @{ = } p{0.85\linewidth} @{} }
    1 & Roe flux\cite{Roe_1981} with entropy fix\cite{Huynh_1995} \\
    2 & HLLC flux\cite{Toro_1994} \\
    3 & LDFSS flux\cite{Edwards_1997} \\
    4 & Rotated-Roe-HLL flux\cite{Nishikawa_2008} (2D and 3D only) \\
    5 & AUSM+ flux\cite{Liou_1996} \\
    6 & Rusanov flux\cite{Rusanov_1961}
\end{tabular}\\
%
\NOTE
Only the Roe flux with entropy fix or the Rusanov flux are supported
when using a rotating reference frame. If any other option is selected
with a rotating reference frame, the inviscid flux method will be reset to
the Roe flux.

\end{nmldescbody}

%==========================================================================

\begin{nmldescitem}
Visc_Flux_Method = 2
\end{nmldescitem}
\inp*[nmlNumerics::ViscFluxMethod]{}[Visc\_Flux\_Method][]

\begin{nmldescbody}
Specifies the method used to compute the common viscous fluxes at the
interfaces.
\\
\begin{tabular}{ @{\qquad} r @{ = } p{0.85\linewidth} @{} }
  1 & Bassi-Rebay 1 (BR1)\cite{Bassi_1997} (high chances of stability issues) \\
  2 & Bassi-Rebay 2 (BR2)\cite{Bassi_2005} \\
  3 & Local DG (LDG)\cite{Cockburn_1998} (smaller CFL limit than BR2)
\end{tabular}
\end{nmldescbody}

%==========================================================================

\begin{nmldescitem}
Visc_Variables_To_Ave = 1
\end{nmldescitem}
\inp*[nmlNumerics::ViscVariablesToAve]{}[Visc\_Variables\_To\_Ave][]

\begin{nmldescbody}
Determines which variables are averaged to the interfaces when computing the
common viscous fluxes.
\\
\begin{tabular}{ @{\qquad} r @{ = } p{0.85\linewidth} @{} }
  1 & Average the interface fluxes. Compute the common viscous flux by
      averaging the left and right viscous interface fluxes which are computed
      using the left and right interface solutions and gradients. \\
  2 & Average the interface solutions and gradients. Compute common solutions
      and gradients by averaging the left and right interface solutions and
      gradients, and then use the common interface values to compute the common
      viscous flux.
\end{tabular}
\end{nmldescbody}

%==========================================================================

%##%\FINALFIX{\newpage}
\begin{nmldescitem}
Combine_Common_Flux_Projections = 0
\end{nmldescitem}
\inp*[nmlNumerics::CombineCommonFluxProjections]{}[Combine\_Common\_Flux\_Projections][]

\begin{nmldescbody}
% Specifies if the common flux projections are combined when using
Specifies if the \key[key:CommonFluxValue]{common flux}[common flux value] projections
are combined when using over-integration by over-sampling at the quadrature points.
\\
\begin{tabular}{ @{\qquad} r @{ = } p{0.85\linewidth} @{} }
  0 & All projections of the over-sampled common fluxes are performed
      separately. \\
  1 & The projections of the over-sampled interior-face common fluxes are
      performed separately from the projection of the over-sampled boundary-face
      common fluxes.\\
  2 & All of the over-sampled common fluxes are computed and combined before
      projecting.
\end{tabular}
\end{nmldescbody}

%==========================================================================

\FINALFIX{\newpage}
\begin{nmldescitem}
Correction_Function = 1
\end{nmldescitem}
\inp*[nmlNumerics::CorrectionFunction]{}[Correction\_Function][]

\begin{nmldescbody}
Specifies the \key[key:CorrectionFunction]{correction function}[=]
(\sref{SEC:Intro:FR:Correction_Functions})
to use when computing the
\key[key:ContinuousFluxFunction]{continuous flux function}[=].
\\
\begin{tabular}{ @{\qquad} r @{ = } p{0.85\linewidth} @{} }
  1 & \textsl{g}\textsubscript{DG} correction function (recovers nodal DG
    method) \\
  2 & \textsl{g}\textsubscript{2} correction function (only works for 1D)\\
  3 & \textsl{g}\textsubscript{Ga} correction function (only works for 1D)
\end{tabular}
\end{nmldescbody}

%==========================================================================


\subsection{nml\_limiter}
\label{SEC:nml_limiter}
\inp*[nmlLimiter]{}[nml\_limiter][]

This \hyperref[SEC:User_Input]{namelist block}
is used to specify limiter, filter, and marker controls.

\begin{nmlblocklst}[
,label={LST:nml_limiter}
,caption={Namelist nml\_limiter.}
]
&nml_limiter
  Global_Filter_Interval      =  0
  Global_Filter_Option        =  0
  Global_Filter_Order         = 16
  Global_Filter_Alpha         = 36.0
  Global_Filter_Etac          =  0.0
  Sponge_Filter_Interval      =  0
  Sponge_Filter_Option        =  0
  Sponge_Filter_Order         =  8
  Sponge_Filter_Alpha         = 36.0
  Sponge_Filter_Etac          =  0.0
  Sponge_Filter_X_Start       = huge(0.0)  ! [input grid units]
  Local_Limiter_Interval      =  0
  Local_Limiter_Option        =  0
  Local_Filter_Order          = 16
  Local_Filter_Alpha          = 36.0
  Local_Filter_Etac           =  0.0
  Marker_Option               =  0
  Marker_Variable             =  0
  Marker_Min                  =  1.0
  Marker_Max                  =  1.0
  Marker_Dpsi                 = -1.0
  Marker_Stored_Interval      =  1
  Thermofix%Method            =  0
  Thermofix%Iterations        = 10
  Thermofix%Check_Cells       =  T
  Thermofix%Check_Faces       =  T
  Thermofix%Check_Nodes       =  T
  Thermofix%Verbose           =  T
  Thermofix%Scale_Factor      = 0.5
  Thermofix%Min_Pressure      = 10000.  ! [Pa]
  Thermofix%Min_Temperature   =   100.  ! [K]
  Thermofix%Min_Density       = 0.001   ! [kg/m^3]
  Positivity%Option           = 0
  Positivity%Density_Limit    = 1.0e-6  ! [kg/m^3]
  Positivity%Pressure_Limit   = 1.0e-6  ! [Pa]
  Entropy_Limiter%Option      = 0
  Entropy_Limiter%Interpolant = 0
  Entropy_Limiter%Tolerance   = 0.001
  Entropy_Limiter%Check_Faces = F
/
\end{nmlblocklst}

%==========================================================================

%\begin{nmldescitem}
%Pyramid_Filter_Option = 1
%\end{nmldescitem}
%\inp*[nmlLimiter::PyramidFilterOption]{}[Pyramid\_Filter\_Option][]
%
%\begin{nmldescbody}
%This input specifies how to filter the modes of pyramid elements.
%%
%All options use a tensor product of a 1D exponential filter with
%a 2D quad exponential filter to generate the pyramid filter matrices for
%all exponential filters: global, limiter, and sponge.
%
%The 1D filter is relative to the pyramid height from base to peak, with the
%base being the constant mode, and the peak being the \p{}-th mode, where
%$\p{}=\inp[nmlNumerics::SolutionOrder]{Solution\_Order}$.
%%
%The same 1D filter is used for all options below.
%
%The 2D filter at the base is equivalent for all options below, and is also
%the same as the 2D filter for an actual quad element of order \p{}.
%%
%Each increase in height of the pyramid subtracts 1 from the order of the
%quad used to create the 2D filter (\thatis the base is a quad of order \p{},
%the next level up is a quad of order \p{}-1, then a quad of order \p{}-2,
%and eventually at the peak a ``quad'' of order \p{}-\p{}=0).
%%
%The options below essentially amount to differences in how the 2D filter
%is defined for all of the pyramid layers above the base.
%
%Let $k$ be the level of the pyramid ranging from 0 to \p{}, where 0 is the
%base and \p{} is the peak.
%%
%Let the cutoff order \pc{} at level $k$, denoted by \pc{\p{}-k}{},
%be defined as
%  $\pc{\p{}-k}=\eta_c * \mathrm{real}(\p{}{-k})$,
%where $\eta_c$ is the cutoff ratio defined in
%\sref{SEC:Filtering}.
%
%The permissible values for \inptext{Pyramid\_Filter\_Option}
%can be summarized as:
%\\
%\begin{tabular}{ @{\qquad} r @{ = } p{0.85\linewidth} @{} }
%    1 & Scale $\eta_c$ so that the cutoff order \pc{} remains constant
%        relative to \p{} for all $k$. Note that $\eta_c$ will take the
%        value 1 if it evaluates to something larger than 1.\\
%    2 & Use the same value of $\eta_c$ for each level of the pyramid.
%        This means the cutoff order will change at each level.
%        For example, if
%        \p{}=3 and $\eta_c$=0.5: at the base level ($k$=0)
%        $\pc{\p{}-0}=0.5*\mathrm{real}(\p{}{-0})=2$ and at the second level
%        ($k$=2) $\pc{\p{}-2}=0.5*\mathrm{real}(\p{}{-2})=1$.
%        This results in a slightly more dissipative filter than the
%        first option.\\
%    3 & Use a full 3D tensor product of the 1D exponential filter.
%        This is equivalent to extracting only the pyramid modes from
%        the hex filter.
%\end{tabular}
%
%\NOTE
%After examining how the modes of other cell geometries are filtered,
%the default value of 1 appears to be the ``most correct'' way
%to filter the modes of pyramid cells.
%\end{nmldescbody}

%==========================================================================

\begin{nmldescitem}
Global_Filter_Interval = 0
\end{nmldescitem}
\inp*[nmlLimiter::GlobalFilterInterval]{}[Global\_Filter\_Interval][]

\begin{nmldescbody}
Specifies the frequency to apply the global filter.
\\
\begin{tabular}{ @{\quad} r @{ \quad } p{0.85\linewidth} @{} }
  $\le 0$ & Apply the global filter after each Runge-Kutta stage. \\
 %$> 0$   & Apply the global filter after the last Runge-Kutta stage
 %          of every
 %          $\inptext{Global\_Filter\_Interval}^{\;\textrm{th}}$
 %          time step. \\
  $> 0  $ & Apply the global filter after the last Runge-Kutta stage if
            \newline
            $\textrm{mod} \left(         {current\_iteration},
                                 \inptext{Global\_Filter\_Interval}
                          \right) = 0$.
            \newline
            The ${current\_iteration}$ is the total aggregate iteration
            number including previous jobs.
 %          \\
 %$< 0  $ & The global filter will NOT be applied even if
 %          $\inptext{Global\_Filter\_Option} \neq 0$
\end{tabular}
\end{nmldescbody}

%==========================================================================

\begin{nmldescitem}
Global_Filter_Option = 0
\end{nmldescitem}
\inp*[nmlLimiter::GlobalFilterOption]{}[Global\_Filter\_Option][]

\begin{nmldescbody}
Specifies whether the global filter is used for stabilization.
%
This filter is denoted as a global filter because it is used on
all cells regardless of the quality of the solution in each cell.
%
The available options are as follows:
\\
\begin{tabular}{ @{\quad} r @{\quad} p{0.80\linewidth} @{} }
  $0=$ & Disable the global filter.\\
  $1=$ & Use the exponential filter, \eref{EQN:ExpFilterMatrix}, 
         where the polynomial order is given by 
         $\PO=\inp[nmlNumerics::SolutionOrder]{Solution\_Order}$
         for 1D, 2D, and 3D.
\end{tabular}

GFR uses a three-pass filtering approach involving global, local, and
sponge filters.
%
Each of them and their sequence of application is described in
\sref{SEC:Filtering}.
%
Generally, we recommend that the global filter not be used, because it can
introduce significant dissipation.
%
When it is absolutely necessary to use the global filter, the filter settings
should be set as weak as possible.
\end{nmldescbody}

%==========================================================================

\begin{nmldescitem}
Global_Filter_Order = 16
\end{nmldescitem}
\inp*[nmlLimiter::GlobalFilterOrder]{}[Global\_Filter\_Order][]

\begin{nmldescbody}
This is an integer value between ($2 \leq s \leq 32$) that specifies the
exponent\slash order of the global exponential filter,
\eref{EQN:ExpFilterMatrix}.
%
As discussed in \sref{SEC:Filtering} and shown in \fref{FIG:Filter_s},
the filter order controls how rapidly the filter decays.
Larger values act to concentrate the filtering toward the highest mode.

Because the global filter is applied to all cells regardless of the
solution quality in each cell, one generally wants to use
as large a value as possible so as not to affect the lower modes.
%
Starting with \inptext{Global\_Filter\_Order} of 32 or 16
and halving if necessary is a good approach.
\end{nmldescbody}

%==========================================================================

\begin{nmldescitem}
Global_Filter_Alpha = 36.0
\end{nmldescitem}
\inp*[nmlLimiter::GlobalFilterAlpha]{}[Global\_Filter\_Alpha][]

\begin{nmldescbody}
This is a real value between ($0. < \alpha \leq 36.$) that specifies the damping
parameter of the global exponential filter, \eref{EQN:ExpFilterMatrix}.
%
As discussed in \sref{SEC:Filtering}, and shown in \fref{FIG:Filter_alpha},
this parameter controls the level of damping applied to the highest mode;
larger values apply more damping.
%
A value of 36.0 effectively filters the highest mode to machine zero.

Because the global filter is applied to all cells regardless of the
solution quality in each cell, one ideally wants to use
as small a value as possible so as not to introduce unwanted dissipation.
%
However, since numerical error often first appears in the highest modes
it may be necessary to aggressively filter those modes before it
contaminates the lower modes.
%
Starting with \inptext{Global\_Filter\_Alpha} of 2 or 4
and doubling if necessary is a good approach.

In the special case of \PO[=1], the \PO[1] mode will be multiplied by
$e^{-\alpha}$.
\end{nmldescbody}

%==========================================================================

\begin{nmldescitem}
Global_Filter_Etac = 0.0
\end{nmldescitem}
\inp*[nmlLimiter::GlobalFilterEtac]{}[Global\_Filter\_Etac][]

\begin{nmldescbody}
This is a real value between ($0. \leq \eta_c \le 1.$) that specifies the modal
cutoff fraction for the global exponential filter, \eref{EQN:ExpFilterMatrix}.
%
As discussed in \sref{SEC:Filtering} and shown in \fref{FIG:Filter_etac},
modes at or below this cutoff fraction will be entirely preserved.
%
A value of $1.0$ applies no filter and wastes computational resources.

For most simulations it is desirable to preserve the linear mode,
and so a good approach is to set this input value to $\eta_c=1/{\PO}$ where
$\PO=\inp[nmlNumerics::SolutionOrder]{Solution\_Order}$.
%
However, if ${\PO} \leq 2$ it may be best to use a value of 0,
and if ${\PO} > 4$ it may be possible to use a larger fraction.

In the special case of \PO[=1], any value $<1.0$ will just multiply the \PO[1]
mode by $e^{-\alpha}$.
\end{nmldescbody}

%==========================================================================

\begin{nmldescitem}
Sponge_Filter_Interval =  0
Sponge_Filter_Option   =  0
Sponge_Filter_Order    =  8
Sponge_Filter_Alpha    = 36.0
Sponge_Filter_Etac     =  0.0
Sponge_Filter_X_Start  = huge(0.0)  ! [input grid units]
\end{nmldescitem}
\inp*[nmlLimiter::SpongeFilterInterval]{}[Sponge\_Filter\_Interval][]
\inp*[nmlLimiter::SpongeFilterOption]{}[Sponge\_Filter\_Option][]
\inp*[nmlLimiter::SpongeFilterOrder]{}[Sponge\_Filter\_Order][]
\inp*[nmlLimiter::SpongeFilterAlpha]{}[Sponge\_Filter\_Alpha][]
\inp*[nmlLimiter::SpongeFilterEtac]{}[Sponge\_Filter\_Etac][]
\inp*[nmlLimiter::SpongeFilterXStart]{}[Sponge\_Filter\_X\_Start][]

\begin{nmldescbody}
GFR uses a three-pass filtering approach involving global, local, and
sponge filters.
%
Each of them and their sequence of application is described in
\sref{SEC:Filtering}.
%
Setting \inptext{Sponge\_Filter\_Option}=1 activates the sponge filter,
which is ideally used as an extra-strong exponential filter to damp
non-physical pressure reflections near the outflow boundary.
%
It is applied to any cells where all the $x$-coordinate values
(in \key[key:GridUnits]{grid units}) are larger than
that specified by \inptext{Sponge\_Filter\_X\_Start}.

The details for these inputs are the same as for the similarly
named \inptext{Global\_Filter} options.
% 
The user must manually specify the starting $x$-location
(in \key[key:GridUnits]{grid units}) of the sponge region via
\\% Prevent next word from entering right margin.
\inptext{Sponge\_Filter\_X\_Start},
since the default value of \inptext{huge(0.0)} will be beyond
any current grid location.

\noindent\NOTE
If $\inp[nmlReference::AoAAlpha]{AoA\_Alpha}=\pm90$\textdegree{} and
$\inp[nmlReference::AoABeta]{AoA\_Beta}=0$\textdegree{}, the $y$-axis will be
assumed as the axial coordinate direction, and this input will be applied to the
$y$-coordinate.
%
If $\inptext{AoA\_Alpha}=0$\textdegree{} and
$\inptext{AoA\_Beta}=\pm90$\textdegree{}, the $z$-axis will be assumed as the
axial coordinate direction, and this input will be applied to the
$z$-coordinate.
\end{nmldescbody}

%==========================================================================

\FINALFIX{\newpage}
\begin{nmldescitem}
Local_Limiter_Interval = 0
\end{nmldescitem}
\inp*[nmlLimiter::LocalLimiterInterval]{}[Local\_Limiter\_Interval][]

\begin{nmldescbody}
Specifies the frequency to apply the local limiter.
\\
\begin{tabular}{ @{\quad} r @{ \quad } p{0.85\linewidth} @{} }
  $\le 0$ & Apply the limiter after each Runge-Kutta stage. \\
  $> 0  $ & Apply the limiter after the last Runge-Kutta stage if
            \newline
            $\textrm{mod} \left(         {current\_iteration},
                                 \inptext{Local\_Limiter\_Interval}
                          \right) = 0$.
            \newline
            For example, if $\inptext{Local\_Limiter\_Interval}=1$, the local
            limiter will be applied at the end of every time step. The
            ${current\_iteration}$ is the total aggregate iteration
            number including previous jobs.
 %          \\
 %$< 0  $ & The limiter will NOT be applied even if \newline
 %          $\inptext{Local\_Limiter\_Option} \neq 0$
\end{tabular}
\end{nmldescbody}

%==========================================================================

\begin{nmldescitem}
Local_Limiter_Option = 0
\end{nmldescitem}
\inp*[nmlLimiter::LocalLimiterOption]{}[Local\_Limiter\_Option][]

\begin{nmldescbody}
GFR uses a three-pass filtering approach involving global, local, and
sponge filters.
%
Each of them and their sequence of application is described in
\sref{SEC:Filtering}.
%
\inptext{Local\_Limiter\_Option} specifies the type of limiter to apply to
locally marked cells.
\\
\begin{tabular}{ @{\quad} r @{ $=$ } p{0.85\linewidth} @{} }
  $   0$ & Do not apply any kind of limiter. \\
  $   1$ & Apply filtering. \\
  $   2$ & Apply scaled filtering.
           The filter value is scaled by the magnitude of the scaled marker. \\
  $   3$ & Apply the multi-dimensional limiting process (MLP). \\
  $9001$ & Apply recursive projection.
           For any troubled cell that is marked, project the solution down one
           order; recheck that cell to see if it is still troubled; keep
           repeating this process until the cell is no longer marked or has been
           reduced to a \PO[1] solution.
           \\
  $9002$ & Apply scaled projection.
           For any troubled cell that is marked,
           project the solution down to a order that is dynamically
           chosen by the value of the local marker.
           Cells with higher marker values will be projected to lower orders,
           whereas cells with lower marker values will be projected down
           only an order or two from the solution order.
\end{tabular}

\noindent
\NOTE
%
The locally marked cells are determined from the inputs for
  \inp[nmlLimiter::MarkerOption]{Marker\_Option} and
  \inp[nmlLimiter::MarkerVariable]{Marker\_Variable}.
%
The local filter settings are defined by input values for
  \inp[nmlLimiter::LocalFilterOrder]{Local\_Filter\_Order},
  \inp[nmlLimiter::LocalFilterAlpha]{Local\_Filter\_Alpha}, and
  \inp[nmlLimiter::LocalFilterEtac]{Local\_Filter\_Etac}.

\end{nmldescbody}

%==========================================================================

\begin{nmldescitem}
Local_Filter_Order  = 16
Local_Filter_Alpha  = 36.0
Local_Filter_Etac   =  0.0
\end{nmldescitem}
\inp*[nmlLimiter::LocalFilterOrder]{}[Local\_Filter\_Order][]
\inp*[nmlLimiter::LocalFilterAlpha]{}[Local\_Filter\_Alpha][]
\inp*[nmlLimiter::LocalFilterEtac]{}[Local\_Filter\_Etac][]

\begin{nmldescbody}
As discussed in \sref{SEC:Filtering}, the local filter is applied
to those cells identified by the troubled cell marker.
%
A good approach is to use the local options to apply a stronger filter
to only those cells that have been marked as a troubled.
%
This will help prevent the oscillatory solutions in these cells
from causing the simulation to diverge.
%
Then, if necessary, a weaker global filter can be applied
to all cells to damp any other disturbances.

The details for these inputs are the same as for the similarly
named \inptext{Global\_Filter} options, but one likely wants
to apply a stronger local filter setting.
%
Use
\inptext{Local\_Filter\_Order}=16 and halve if necessary, and
\inptext{Local\_Filter\_Alpha}=36 and halve if possible.
For \inptext{Local\_Filter\_Etac}, use the same recommendations as
for \inp[nmlLimiter::GlobalFilterEtac]{Global\_Filter\_Etac}.

\IDEA{
Create two paragraphs here: one with recommendations for
$\inptext{Local\_Limiter\_Option}=1$, and one with recommendations for 
$\inptext{Local\_Limiter\_Option}=2$.
}
\end{nmldescbody}

%==========================================================================

\FINALFIX{\newpage}
\begin{nmldescitem}
Marker_Stored_Interval = 1
\end{nmldescitem}
\inp*[nmlLimiter::MarkerStoredInterval]{}[Marker\_Stored\_Interval][]

\begin{nmldescbody}
This determines if the marker values written to the CGNS solution files
are from the final Runge-Kutta stage or from each Runge-Kutta stage.
\\
\begin{tabular}{ @{\qquad} r @{ \quad } p{0.85\linewidth} @{} }
  $<0$ & Store markers at the same frequency of
         \inp[nmlLimiter::LocalLimiterInterval]{Local\_Limiter\_Interval}.
         This can be used to synchronize the marker storage and
         limiter application.
         \\
  $=0$ & Store markers after each Runge-Kutta stage. \\
  $>0$ & Store markers after the last Runge-Kutta stage.
\end{tabular}
\\
\NOTE
Storing markers each stage increases memory and file size
(if \inp[nmlOutputVariables::MarkerVariables]{Marker\_Variables}
are being output).
\end{nmldescbody}

%==========================================================================

%##%\FINALFIX{\newpage}
\begin{nmldescitem}
Marker_Option = 0
\end{nmldescitem}
\inp*[nmlLimiter::MarkerOption]{}[Marker\_Option][]

\begin{nmldescbody}
Specifies which type of local marker to use.
\\
\begin{tabular}{ @{\qquad} r @{ = } p{0.85\linewidth} @{} }
  0 & Use the default marker option for the chosen
      \inp[nmlLimiter::LocalLimiterOption]{Local\_Limiter\_Option}. \\
  1 & Use the resolution marker described in \sref{SEC:ResolutionMarker}.\\
  2 & Use the jump marker described in \sref{SEC:JumpMarker}.\\
  3 & Use the MLP marker described in \sref{SEC:MLPMarker}.\\
  4 & Use the Ducros marker described in \sref{SEC:DucrosMarker}.
\end{tabular}\\
%
\Tref{TAB:Permissible_Marker_Option} lists the
permissible and default values for \inptext{Marker\_Option},
depending on
\\ % Prevent next word from entering right margin.
\inptext{Local\_Limiter\_Option}.

\IDEA{Add Ducros and Ducros$\times$Jump markers.}

\begin{table}[h]
\begin{center}
\caption{Permissible values for \inptext{Marker\_Option}.}
\label{TAB:Permissible_Marker_Option}
\begin{tabular}{rlc}
  \multicolumn{2}{c}{\inptext{Local\_Limiter\_Option}} &
  \inptext{Marker\_Option} \\
  \hline
     0 & None                & $1^*,2,3,4$ \\
     1 & Filtering           & $1^*,2,3,4$ \\
     2 & Scaled Filtering    & $1^*,2,3,4$ \\
     3 & MLP                 & $3^*$     \\
  9001 & Recusive Projection & $1^*,2$   \\
  9002 & Scaled   Projection & $1^*,2$   \\
  \hline
\end{tabular}
\\
${}^*$ indicates default value.
\end{center}
\end{table}

\end{nmldescbody}

%==========================================================================

\begin{nmldescitem}
Marker_Variable = 0
\end{nmldescitem}
\inp*[nmlLimiter::MarkerVariable]{}[Marker\_Variable][]

\begin{nmldescbody}
This determines which variable is used to evaluate the local marker.
\\
\begin{tabular}{ @{\qquad} r @{ = } p{0.85\linewidth} @{} }
  -2 & Use the Ducros variable based upon velocity divergence (only valid for
       $\inptext{Marker\_Option}=4$).\\
  -1 & Use pressure.\\
   0 & Use the default marker variable for the chosen
       \inp[nmlLimiter::MarkerOption]{Marker\_Option}.
       \\
   1 & Use density.\\
   2 & Use $x$-momentum.\\
   3 & Use $y$-momentum.\\
   4 & Use $z$-momentum.\\
   5 & Use total energy.\\
   6 & Use the max value after evaluating the marker
       for each conserved variable.
\end{tabular}\\
%
\Tref{TAB:Permissible_Marker_Variable} lists the
permissible and default values for \inptext{Marker\_Variable},
depending on
\\ % Prevent next word from entering right margin.
\inptext{Marker\_Option}.

\begin{table}[h]
\begin{center}
\caption{Permissible values for \inptext{Marker\_Variable}.}
\label{TAB:Permissible_Marker_Variable}
\begin{tabular}{rlc}
  \multicolumn{2}{c}{\inptext{Marker\_Option}} & \inptext{Marker\_Variable} \\
  \hline
     1 & Resolution & $ 1^*,2,3,4,5,6$ \\
     2 & Jump       & $-1^*,1,2,3,4,5$ \\
     3 & MLP        & $ 1^*,2,3,4,5$ \\
     4 & Durcos     & $-2^*$ \\
  \hline
\end{tabular}
\\
${}^*$ indicates default value.
\end{center}
\end{table}

\NOTE
Using all conserved variables (option 6) will result in marking
the most number of cells.
This will increase stability, but has the tendency to add unnecessary
dissipation where it is not needed.
The default of using only density seems to result in the best mix of
stability with minimal unnecessary dissipation.
\end{nmldescbody}

%==========================================================================

\begin{nmldescitem}
Marker_Min = 1.0
\end{nmldescitem}
\inp*[nmlLimiter::MarkerMin]{}[Marker\_Min][]

\begin{nmldescbody}
This gives the value for $\Psi_{min} = \Psi_0 - \Delta \Psi$,
which is the bottom-end of the $s$-curve used to scale
the marker value from the interval ($-\infty$,0] to the interval [0,1].
%
Any cell with a marker value less than the value of this input will have
a scaled marker value of 0, which means that cell will receive no limiting.
%
See the discussion of the scaled marker in \sref{SEC:ScaledMarker}
for a more information
\end{nmldescbody}

%==========================================================================

\begin{nmldescitem}
Marker_Max = 1.0
\end{nmldescitem}
\inp*[nmlLimiter::MarkerMax]{}[Marker\_Max][]

\begin{nmldescbody}
This gives the value for $\Psi_{max} = \Psi_0 + \Delta \Psi$,
which is the top-end of the $s$-curve used to scale
the marker value from the interval ($-\infty$,0] to the interval [0,1].
%
Any cell with a marker value greater than the value of this input will have
a scaled marker value of 1, which means that cell will receive
the maximum strength of the limiting method being employed.
%
See the discussion of the scaled marker in \sref{SEC:ScaledMarker}
for a more information
\end{nmldescbody}

%==========================================================================

\begin{nmldescitem}
Marker_Dpsi = -1.0
\end{nmldescitem}
\inp*[nmlLimiter::MarkerDpsi]{}[Marker\_Dpsi][]

\begin{nmldescbody}
This gives the value for the constant $\Delta\Psi$ that is used to scale
the marker value from the interval ($-\infty$,0] to the interval [0,1].
%
See the discussion of the scaled marker in \sref{SEC:ScaledMarker}
for a more information
\end{nmldescbody}

%==========================================================================

\FINALFIX{\newpage}
\begin{nmldescitem}
Thermofix%Method          = 0
Thermofix%Iterations      = 10
Thermofix%Check_Cells     = T
Thermofix%Check_Faces     = T
Thermofix%Check_Nodes     = T
Thermofix%Verbose         = T
Thermofix%Scale_Factor    = 0.5
Thermofix%Min_Pressure    = 10000.  ! [Pa]
Thermofix%Min_Temperature = 100.    ! [K]
Thermofix%Min_Density     = 0.001   ! [kg/m^3]
\end{nmldescitem}
\inp*[nmlLimiter::Thermofix/Method]{}[Thermofix\%Method][]
\inp*[nmlLimiter::Thermofix/Iterations]{}[Thermofix\%Iterations][]
\inp*[nmlLimiter::Thermofix/CheckCells]{}[Thermofix\%Check\_Cells][]
\inp*[nmlLimiter::Thermofix/CheckFaces]{}[Thermofix\%Check\_Faces][]
\inp*[nmlLimiter::Thermofix/CheckNodes]{}[Thermofix\%Check\_Nodes][]
\inp*[nmlLimiter::Thermofix/Verbose]{}[Thermofix\%Verbose][]
\inp*[nmlLimiter::Thermofix/ScaleFactor]{}[Thermofix\%Scale\_Factor][]
\inp*[nmlLimiter::Thermofix/MinPressure]{}[Thermofix\%MinPressure][]
\inp*[nmlLimiter::Thermofix/MinTemperature]{}[Thermofix\%Min\_Temperature][]
\inp*[nmlLimiter::Thermofix/MinDensity]{}[Thermofix\%Min\_Density][]

\begin{nmldescbody}
Sometimes during the course of a simulation, the solution will
encounter a negative thermodynamic property
(pressure, temperature, or density),
which can cause a NaN (\forexample when computing speed of sound).
%
The \inptext{Thermofix} inputs are intended to help remedy this
situation.
%
While they provide additional stability when needed, they also introduce
additional dissipation, and should be avoided if possible.

\inptext{Method} specifies the method used to repair
cells with bad thermodynamic properties.
Checks are performed during each Runge-Kutta stage.
\\
\begin{tabular}{ @{\qquad} r @{ = } p{0.85\linewidth} @{} }
  0 & None \\
  1 & Limit the high-order modes using the exponential sponge filter. \\
  2 & Scale the high-order modes by \inptext{Scale\_Factor}.
\end{tabular}

\inptext{Iterations} specifies how many passes to
make before abandoning further repair efforts.

\inptext{Scale\_Factor} specifies the scaling factor to apply
to the high-order modes when using \inptext{Method}=2.

\inptext{Check\_Cells},
\inptext{Check\_Faces}, and
\inptext{Check\_Nodes}
specifies whether to check the thermodynamic properties at points within the
cell volume, the cell faces, or the cell nodes (corners).
%
If any of these checks identifies a troubled point,
the solution within the cell volume is treated and
re-interpolated to the faces and nodes.

\inptext{Verbose} is used to output information about
the bad cells to the display output.

\inptext{Min\_Pressure},
\inptext{Min\_Temperature}, and
\inptext{Min\_Density}
are used to specify threshold values below which action should be taken.
%
These limits are imposed individually, so the three specified values need not
satisfy the equation of state.
%
Using small yet positive values is intended to identify problem
regions before the solution has deteriorated too severely.
%
The actual values for these parameter may need to be adjusted
for each flow problem.
%
These values should be small enough to not be physically encountered within
the flow but large enough that the solution can still be recovered before
fully diverging.

The thermofix limiter is applied at the end of each Runge-Kutta stage.
%
The order in which the thermofix limiter is applied relative to
other limiting\slash filtering methods can be found in the enumerated list
towards the beginning \sref{SEC:Numerics}.

\end{nmldescbody}

%==========================================================================

\FINALFIX{\newpage}
\begin{nmldescitem}
Positivity%Option         = 0
Positivity%Density_Limit  = 1.0e-6 ! [kg/m^3]
Positivity%Pressure_Limit = 1.0e-6 ! [Pa]
\end{nmldescitem}
\inp*[nmlLimiter::Positivity/Option]{}[Positivity\%Option][]
\inp*[nmlLimiter::Positivity/DensityLimit]{}[Positivity\%Density\_Limit][]
\inp*[nmlLimiter::Positivity/PressureLimit]{}[Positivity\%Pressure\_Limit][]

\begin{nmldescbody}
 
%ADD REFERENCE TO ZHANG AND SHU, and JOHNSON AND CHING.

Sometimes during the course of a simulation, the solution will encounter a
negative thermodynamic property (pressure, temperature, or density), which can
cause a NaN (\forexample when computing speed of sound).
%
The \inptext{Positivity} inputs are intended to help remedy this situation using
a positivity-preserving limiter.
%
While this method provides additional stability when needed, it will also
introduce additional dissipation, and should therefore be avoided if possible.

\inptext{Option} simply specifies if the positivity limiter is activated.
\\
\begin{tabular}{ @{\qquad} r @{ = } p{0.85\linewidth} @{} }
  0 & Off (Default) \\
  1 & On
\end{tabular}

\inptext{Density\_Limit} specifies the minimum dimensional density allowable for
any of the solution or flux points of a cell.
 
\inptext{Pressure\_Limit} specifies the minimum dimensional pressure allowable
for any of the solution or flux points of a cell.

When the positivity limiter is activated, if any solution or flux point of a
cell has a density or pressure value below their respective limits, the
positivity limiter will be applied to this cell.

The positivity limiter is applied at the end of each Runge-Kutta stage.
%
The order in which the positivity-preserving limiter is applied relative to
other limiting\slash filtering methods can be found in the enumerated list
towards the beginning \sref{SEC:Numerics}.
%
%The positivity limiter will be applied \emph{after} any of the following (if
%also active): sponge filter or weak global filter (\thatis
%$\inptext{Global\_Filter\_Alpha} \leq 1$).
%%
%The positivity limiter will be applied \emph{before} any of the following (if
%also active): entropy limiter, local limiter (\thatis
%$\inptext{Local\_Limiter\_Option}\neq 0$), strong global filter (\thatis
%$\inptext{Global\_Filter\_Alpha}>1$), or \inptext{Thermofix}.

When the positivity limiter is applied to a cell, the solution in the cell is
modified in the following way:
%
\begin{equation}
  \tilde{u} = \check{u} + \theta_{PL}\left(u-\check{u}\right)
  \label{EQN:PLSolution}
\end{equation}
%
where $u$ is the conserved variable being limited, $\check{u}$ is the average
value of the polynomial $u$ in the cell, $\theta_{PL}$ is the limiter scaling
factor, and $\tilde{u}$ is the resulting limited polynomial.
%
The limiter scaling factor, $\theta_{PL}$, is computed as
%
\begin{equation}
  \theta_{PL} = 0.99 \left(
                           \frac{\check{a}-a_{\textrm{Lim}}}
                                {\check{a}-a_{\textrm{min}}}
                     \right)
  \label{EQN:PLScalingFactor}
\end{equation}
%
where $a$ is either density or pressure depending on which positivity limit was
violated in the cell, $\check{a}$ is the average value of the variable $a$ in
the cell, $a_{\textrm{Lim}}$ is the value of the respective input
\inptext{Density\_Limit} or \inptext{Pressure\_Limit}, and $a_{\textrm{min}}$ is
the minimum value of the respective variable from all solution and flux points
of the cell.

\NOTE
If \inptext{Density\_Limit} is violated in the cell, only the density polynomial
is limited by the positivity limiter, with the scaling factor $\theta_{PL}$
evaluated with $a=\rho$.
%
If \inptext{Pressure\_Limit} is violated in the cell, all conserved variables
are limited by the positivity limiter, with the scaling factor $\theta_{PL}$
evaluated with $a=p$.

\NOTE
The density and pressure limits needed for the positivity limiter are generally
much smaller than what is needed for \inptext{Thermofix}.

\end{nmldescbody}

%==========================================================================

\FINALFIX{\newpage}
\begin{nmldescitem}
Entropy_Limiter%Option      = 0
Entropy_Limiter%Interpolant = 0
Entropy_Limiter%Tolerance   = 0.001
Entropy_Limiter%Check_Faces = F
\end{nmldescitem}
\inp*[nmlLimiter::EntropyLimiter/Option]{}[Entropy\_Limiter\%Option][]
\inp*[nmlLimiter::EntropyLimiter/Interpolant]{}[Entropy\_Limiter\%Interpolant][]
\inp*[nmlLimiter::EntropyLimiter/Tolerance]{}[Entropy\_Limiter\%Tolerance][]
\inp*[nmlLimiter::EntropyLimiter/CheckFaces]{}[Entropy\_Limiter\%Check\_Faces][]

\begin{nmldescbody}
 
%ADD REFERENCE TO ZHANG AND SHU, and JOHNSON AND CHING.

Sometimes during the course of a simulation, the solution will encounter a
negative thermodynamic property (pressure, temperature, or density), which can
cause a NaN (\forexample when computing speed of sound).
%
The \inptext{Entropy\_Limiter} inputs are intended to help remedy this situation
using a limiter to prevent violations of the local entropy bound in a cell.
%
While this method provides additional stability when needed, it will also
introduce additional dissipation, and should therefore be avoided if possible.

\inptext{Option} specifies if the entropy limiter is active, and if so, which
method to use to limit any cells that violate the minimum entropy bounds.
\\
\begin{tabular}{ @{\qquad} r @{ = } p{0.85\linewidth} @{} }
  0 & Off (Default) \\
  1 & Flat filter \\
  2 & Scaled filter
\end{tabular}

\inptext{Interpolant} specifies the interpolant used by the entropy limiter.
%
For each cell that is found to be in violation of the local entropy bounds, the
interpolant is used to compute the filter strength, $\theta_{EL}$, necessary to
prevent the cell from committing the violation.
\\
\begin{tabular}{ @{\qquad} r @{ = } p{0.85\linewidth} @{} }
  0 & method used by Lv-Ihme; equation 44 from \linecite{Lv_2015} (Default) \\
  1 & linear \\
  2 & density weighted \\
  3 & log scaling \\
  4 & Pad\'{e} approximation
\end{tabular}
   
\inptext{Tolerance} specifies how much the local entropy bounds can be violated
before limiting the cell.
%
As this value gets closer to zero, the limiter gets more strict about not
allowing the solution to violate the bounds by adding more dissipation to
prevent the violation from occurring.

\inptext{Check\_Faces} specifies if the entropy values from the faces of each
cell (both internal and from the adjacent neighboring cells) are included when
computing the minimum entropy for said cell.
%
Including the face entropy values in the minimum entropy will help with
maintaining a stable solution, but at the expense of additional dissipation.

%%%%%%%%%%%%%%%%%%%%%%%%%%%%%%%%%%%%%%%%%%%%%%%%%%%%%%%%%%%%%%%%%%%%%%%%%%%%%%%%

When the entropy limiter uses the flat filter to modify the solution in a cell,
the limited solution polynomial is given by
%
\begin{equation}
  \tilde{u} = \check{u} + \theta_{EL}\left(u-\check{u}\right)
  \label{EQN:ELFlatSolution}
\end{equation}
%
where $u$ is the conserved variable being limited, $\check{u}$ is the average
value of the polynomial $u$ in the cell, $\tilde{u}$ is the resulting
limited polynomial, and the filter strength is in the range $\theta_{EL} =
\left[0,1\right]$.
%
When using the scaled filter, the entropy limiter modifies the solution in one
of two ways depending on the value of $\theta_{EL}$:
%
\begin{equation}
  \tilde{u} = 
  \begin{cases}
    u + \theta_{EL}\left({\cal F}\left(u\right) - u\right) & 
       \textrm{if } \theta_{EL} \leq 1 \\
    \check{u} + \left(2-\theta_{EL}\right) \left({\cal F}\left(u\right) -
    \check{u}\right) &
       \textrm{if } \theta_{EL} > 1
  \end{cases}
  \label{EQN:ELScaledSolution}
\end{equation}
%
where ${\cal F}\left(u\right)$ is the result of applying the full local filter
to the unmodified solution $u$, and the filter strength is in the range
$\theta_{EL} = \left[0,2\right]$.
  
%%%%%%%%%%%%%%%%%%%%%%%%%%%%%%%%%%%%%%%%%%%%%%%%%%%%%%%%%%%%%%%%%%%%%%%%%%%%%%%%

The entropy limiter is applied at the end of each Runge-Kutta stage.
%
The order in which the entropy limiter is applied relative to other
limiting\slash filtering methods can be found in the enumerated list towards the
beginning \sref{SEC:Numerics}.
%
%The entropy limiter will be applied \emph{after} any of the following (if also
%active): sponge filter, weak global filter (\thatis
%$\inptext{Global\_Filter\_Alpha} \leq 1$), or positivity-preserving limiter.
%%
%The entropy limiter will be applied \emph{before} any of the following (if also
%active): local limiter (\thatis $\inptext{Local\_Limiter\_Option}\neq 0$),
%strong global filter (\thatis $\inptext{Global\_Filter\_Alpha}>1$), or
%\inptext{Thermofix}.

%%%%%%%%%%%%%%%%%%%%%%%%%%%%%%%%%%%%%%%%%%%%%%%%%%%%%%%%%%%%%%%%%%%%%%%%%%%%%%%%

\NOTE
When using $\inptext{Option}=2$, the local filter will be used to compute a
filtered solution to use as part of the limiting process; therefore, the input
coefficients for the local filter (\inptext{Local\_Filter\_Order},
\inptext{Local\_Filter\_Alpha}, and \inptext{Local\_Filter\_Etac}) will become
an important part of the entropy limiter settings and need to be adjusted
accordingly.

\NOTE
When the entropy limiter is applied to a cell, all conserved variables are
modified together with the same filter strength.

\NOTE
So far, our experience seems to indicate that the \inptext{Interpolant} has
little effect on the stability or accuracy of the solution; however, this could
be something that is problem dependent and we have yet to find a case that
demonstrates the effect of different interpolant methods.
%
Because of this, it is recommended to first look at adjusting other inputs when
encountering stability issues.

\end{nmldescbody}

%==========================================================================


\subsection{nml\_time\_ave}
\label{SEC:nml_time_ave}
\inp*[nmlTimeAve]{}[nml\_time\_ave][]

This \hyperref[SEC:User_Input]{namelist block}
controls time averaging of solution variables.

\begin{nmlblocklst}[
,label={LST:nml_time_ave}
,caption={Namelist nml\_time\_ave.}
]
&nml_time_ave
  Compute_Time_Averaging      =  F
  Num_Saved_Time_Ave_Restarts =  2
  Input_Time_Ave_File         = "out.ave"
  Time_Ave_Vel_Is_Axisymm     =  F
  Start_Averaging_Timestep    = -1
  Start_Averaging_Tstar       = -1.0
/
\end{nmlblocklst}

%==========================================================================

\begin{nmldescitem}
Compute_Time_Averaging = F
\end{nmldescitem}
\inp*[nmlTimeAve::ComputeTimeAveraging]{}[Compute\_Time\_Averaging][]

\begin{nmldescbody}
This flag enables or disables the accumulation of time-averaged
flow variables.
%
If enabled, time-averaged restart and CGNS solution files will be
written whenever the standard restart and CGNS solution files are written.


%%BegCut:
%%Describe the restart files using a single paragraph with no line breaks.
%%This was hard to read, so it was replaced with the paragraph below it.
% 
% The time-averaged restart file will be written to the
%   ``\inp[nmlOutput::OutputDir]{Output\_Dir}\slash{Restart\_files}\slash''
% directory.
% %
% In order to save disk space, the name of the time-averaged restart file
% will alternate between
%   ``out.1.ave'' and
%   ``out.2.ave''.
% %
% After successfully writing a new time-averaged restart file, a
% symbolic link will be created that points to the latest
% time-averaged restart file.
% %
% The name of this link will be:
%   ``out.current.ave''.
%%EndCut:


The time-averaged restart file will be written to the
  ``\inp[nmlOutput::OutputDir]{Output\_Dir}\slash{Restart\_files}\slash''
directory.
%
In order to save disk space, the name of the time-averaged restart file
will cycle between the following:
\newline\indent
``\inp[nmlOutput::OutputDir]{Output\_Dir}\slash{Restart\_files}\slash{out.1.ave}''
\newline\indent
``\inp[nmlOutput::OutputDir]{Output\_Dir}\slash{Restart\_files}\slash{out.2.ave}''
\newline\indent
\ldots
\newline\indent
``\inp[nmlOutput::OutputDir]{Output\_Dir}\slash{Restart\_files}\slash{out.N.ave}''
\newline
where N=\inp[nmlTimeAve::NumSavedTimeAveRestarts]{Num\_Saved\_Time\_Ave\_Restarts}.
%
After a new time-averaged restart file is successfully written,
%After successfully writing a new time-averaged restart file,
a symbolic link will be created that points to the latest time-averaged restart
file. The name of this link will be:
\newline\indent
``\inp[nmlOutput::OutputDir]{Output\_Dir}\slash{Restart\_files}\slash{out.current.ave}'' 
\end{nmldescbody}

%==========================================================================

\begin{nmldescitem}
Num_Saved_Time_Ave_Restarts = 2
\end{nmldescitem}
\inp*[nmlTimeAve::NumSavedTimeAveRestarts]{}[Num\_Saved\_Time\_Ave\_Restarts][]

\begin{nmldescbody}
This is used to change the number of time-averaged restart files
(out.$\ast$.ave) that are saved to the Restart\_files directory.
The default is 2, so output will cycle between out.1.ave and out.2.ave.

Through the setting of
\inp[nmlOutput::RestartInterval]{Restart\_Interval},
\inp[nmlOutput::RestartFrequency]{Restart\_Frequency}, and
this input,
one should ensure there is a viable fallback solution to
restart from should an NaN occur.
\end{nmldescbody}

%==========================================================================

%##%\FINALFIX{\newpage}
\begin{nmldescitem}
Input_Time_Ave_File = "out.ave"
\end{nmldescitem}
\inp*[nmlTimeAve::InputTimeAveFile]{}[Input\_Time\_Ave\_File][]

\begin{nmldescbody}
\IDEA{Change this input to \inptext{Input\_Time\_Ave\_File}.}

File path to the time-averaged solution restart file that will be
read during initialization if the following criteria are met:
%
\begin{itemize}
  \item
    $\inp[nmlInitialization::LoadRestartFile]{Load\_Restart\_File}=\mathtt{T}$
  \item One of the following is met:
    \begin{itemize}
      \item
        $\inp[nmlTimeAve::ComputeTimeAveraging]{Compute\_Time\_Averaging}=\texttt{T}$
      \item
        $\inp[nmlTimeAve::StartAveragingTstar]{Start\_Averaging\_Tstar}>0$ and
        $\inptext{Start\_Averaging\_Tstar} \leq \tstar$,
        where \key[key:tstar]{\tstar}[t *@\tstar] is the
        non-dimensional simulation time, \eref{EQN:tstar}, of the data that was
        read in from \inp[nmlInitialization::RestartFile]{Restart\_File}.
      \item
        $\inp[nmlTimeAve::StartAveragingTimestep]{Start\_Averaging\_Timestep}>0$
        and $\inptext{Start\_Averaging\_Timestep} \leq iteration$,
        where $iteration$ is the iteration number of the data that was
        read in from \inptext{Restart\_File}.
    \end{itemize}
\end{itemize}
%
The code will read the time-averaged flow variables of the
previous simulation from which the new simulation was restarted,
and then add to those statistics during the current run job.

\NOTE
If the above criteria are met, but the file specified by
\inptext{Input\_Time\_Ave\_File} does not exist, a completely new time averaging will
be started.
%
This is the case even if you were intending to restart a previous time
averaging, and the reason the file could not be found was due to a small typo in
the file name.

\NOTE
This file is only for \emph{reading} during initialization and
is \emph{not an output file}.
\end{nmldescbody}

%==========================================================================

\begin{nmldescitem}
Time_Ave_Vel_Is_Axisymm = F
\end{nmldescitem}
\inp*[nmlTimeAve::TimeAveVelIsAxisymm]{}[Time\_Ave\_Vel\_Is\_Axisymm][]

\begin{nmldescbody}
% Specifies whether axisymmetric velocity components should be included
% in the time-averaged flow variables.
% 
% If false, the time-averaged Cartesian velocity components
% ($v_x$, $v_y$, $v_z$) will be saved
% along with the Cartesian stress and heat flux data.
% %
% If true, the grid is assumed to be axisymmetric about the $x$-axis and
% the time-averaged
%   streamwise      ($v_x$),
%   radial          ($v_r$), and
%   circumferential ($v_\theta$)
% velocity components will added.
% %
% This option is useful for recording statistics for axisymmetric flows.

This option is useful for recording statistics for axisymmetric flows.
%
It specifies whether axisymmetric velocity components should be included
in the time-averaged flow variables.
%
Note, regardless of the value for this input variable, the Cartesian momentum
and velocity components ($[\rho v_x,\rho v_y,\rho v_z]$ and $[v_x,v_y,v_z]$)
will always be averaged and saved to the time-averaged restart file.
\\[0.25em]
\begin{tabular}{ @{\qquad} r @{ \quad } p{0.85\linewidth} @{} }
  T & The geometry is assumed to be axisymmetric about the $x$-axis.
      %
      The axisymmetric momentum and velocity components in the
      radial ($\rho v_r, v_r$) and
      circumferential ($\rho v_\theta, v_\theta$)
      directions will also be averaged and saved to the
      time-averaged restart file.
      %
      The stress and heat flux data will be for the axisymmetric terms,
      \emph{not} the Cartesian terms
      (\thatis the axisymmetric velocity components are used in the
      time-averaged terms).
      %
      For example, the averaged term \ReyAve{v_y T} will be replaced with
      \ReyAve{v_r T} and \ReyAve{\rho v_x v_z} will be replaced with
      \ReyAve{\rho v_x v_\theta}.
  \\
  F & % The time-averaged Cartesian velocity components
      % ($v_x$, $v_y$, $v_z$) will be saved.
      %
      The stress and heat flux data will be for the Cartesian terms
      (\thatis the Cartesian velocity components are used in the
      time-averaged terms).
\end{tabular}
\\[0.25em]
\NOTE
See \sref{SEC:Files:Time_Averaging_Output} for additional details.

\end{nmldescbody}

%==========================================================================

\begin{nmldescitem}
Start_Averaging_Timestep = -1
\end{nmldescitem}
\inp*[nmlTimeAve::StartAveragingTimestep]{}[Start\_Averaging\_Timestep][]

\begin{nmldescbody}
Iteration number to begin time averaging.
\end{nmldescbody}

%==========================================================================

\begin{nmldescitem}
Start_Averaging_Tstar = -1.0
\end{nmldescitem}
\inp*[nmlTimeAve::StartAveragingTstar]{}[Start\_Averaging\_Tstar][]

\begin{nmldescbody}
Non-dimensional simulation time (\key[key:tstar]{\tstar}[t *@\tstar]),
\eref{EQN:tstar},
to begin time averaging.
\end{nmldescbody}

%==========================================================================


\subsection{nml\_output}
\label{SEC:nml_output}
\inp*[nmlOutput]{}[nml\_output][]

This \hyperref[SEC:User_Input]{namelist block}
is used to control various output files.

\begin{nmlblocklst}[
,label={LST:nml_output}
,caption={Namelist nml\_output.}
]
&nml_output
  Time_Scaling_Factor                =  1.0
  Display_Interval                   =  1
  Residual_Interval                  =  0
  Restart_Interval                   =  0
  Restart_Frequency                  = -1.0  ! tstar
  Restart_Output_Uses_Host_Roots     =  T
  Restart_Output_Disabled            =  F
  Output_Dir                         = "."
  Output_Interval                    = -5
  Output_Frequency                   = -1.0  ! tstar
  Output_Order                       = -3
  Output_Summary_File                =  F
  Dmat_Grad_Temp_Order               = -1
  Loc_Output_Pts                     =  0
  CGNS_High_Order_Output             =  F
  CGNS_Mixed_Output                  =  T
  CGNS_Output_Disabled               =  F
  Probes_Collection_Interval         =  0
  Probes_Write_Interval              =  100
  Point_Probes(1)%xyz(1:3)           =  [-999.,-999.,-999.] ! grid units
  Rake_Probes(1)%num_probes          =  0
  Rake_Probes(1)%xyz_beg(1:3)        =  [-999.,-999.,-999.] ! grid units
  Rake_Probes(1)%xyz_end(1:3)        =  [-999.,-999.,-999.] ! grid units
  Cell_Probes(1)%xyz(1:3)            =  [-999.,-999.,-999.] ! grid units
  Profile_IO_All                     =  F
  Profile_IO_Restart                 =  F
  Profile_IO_CGNS                    =  F
  Profile_IO_CGNS_Subs               =  F
  Lustre_Stripe_Count                =  64
  Lustre_Stripe_Size                 = "4m"
/
\end{nmlblocklst}

%==========================================================================

\begin{nmldescitem}
Time_Scaling_Factor = 1.0  ! [1/s]
\end{nmldescitem}
\inp*[nmlOutput::TimeScalingFactor]{}[Time\_Scaling\_Factor][]

\begin{nmldescbody}
This scaling factor is used to define a non-dimensional time
(\key*[key:tstar]{\tstar}[t *@\tstar])
\begin{gather}
  \label{EQN:tstar}
  \tstar = ( t - t_0 ) \cdot \inptext{Time\_Scaling\_Factor}
\end{gather}
where $t$ is the dimensional time (in seconds).
%
If time-averaging has been activated, then $t_0$ is the dimensional time
(in seconds) when time-averaging began.
Otherwise, $t_0=0$.
%
This non-dimensional time (\tstar)
is what is written to standard output and the residual file.
It is also used by the following input variables to
determine when certain events should occur:
\begin{itemize}
\item
  \inp[nmlPhysics::MovingRefFrame/Spinup/Tstar]{Moving\_Ref\_Frame\%Spinup\%Tstar} for the moving reference frame startup,
\item
  \inp[nmlTimeAve::StartAveragingTstar]{Start\_Averaging\_Tstar}
  for time averaging startup,
\item
  \inp[nmlOutput::RestartFrequency]{Restart\_Frequency} of restart files,
\item
  \inp[nmlOutput::OutputFrequency]{Output\_Frequency} of CGNS files,
  and
\item
  \inp[nmlTimestep::FinalTime]{Final\_Time} for when to stop a simulation.
\end{itemize}


For example, many simulations are based on a number of flow-through periods.
%
In this case, the \inptext{Time\_Scaling\_Factor} should be set to
the value of $U/L$, where
  $U$ is the convective velocity and
  $L$ is the flow-through length of the domain.
\end{nmldescbody}

%==========================================================================

\begin{nmldescitem}
Display_Interval = 1
\end{nmldescitem}
\inp*[nmlOutput::DisplayInterval]{}[Display\_Interval][]

\begin{nmldescbody}
Number of time steps between writing the time-stepping, residual,
and convergence statistics to standard output.
%
The input is treated as follows:
\\
\begin{tabular}{ @{\qquad} r @{ \quad } p{0.75\linewidth} @{} }
  $< 0$ & Reset to
          $\inp[nmlTimestep::NumTimesteps]{Num\_Timesteps} /
           | \inptext{Display\_Interval} |$.\\
  $= 0$ & Reset to 1.\\
  $> 0$ & Use as-is.
\end{tabular}

\NOTE
When running the Taylor-Green vortex problem or the
Taylor-Couette problem, this input is also used as the interval for
writing integrated kinetic energy dissipation rate (KEDR)
statistics to the file ``TaylorGreen\_KEDR.dat'' or
``TaylorCouette\_KEDR.dat'', respectively.
\end{nmldescbody}

%==========================================================================

\begin{nmldescitem}
Residual_Interval = 0
\end{nmldescitem}
\inp*[nmlOutput::ResidualInterval]{}[Residual\_Interval][]

\begin{nmldescbody}
Number of time steps between writing the residual and convergence
statistics for all conserved and primitive variables to the
residual file ``residual.dat''.
%
The input is treated as follows:
\\
\begin{tabular}{ @{\qquad} r @{ \quad } p{0.75\linewidth} @{} }
  $< 0$ & Reset to
          $\inp[nmlTimestep::NumTimesteps]{Num\_Timesteps} /
           | \inptext{Residual\_Interval} |$.\\
  $= 0$ & Reset to
          $\inp[nmlTimestep::NumTimesteps]{Num\_Timesteps} / 500$.\\
  $> 0$ & Use as-is.
\end{tabular}
\end{nmldescbody}

%==========================================================================

\begin{nmldescitem}
Restart_Interval = 0
\end{nmldescitem}
\inp*[nmlOutput::RestartInterval]{}[Restart\_Interval][]

\begin{nmldescbody}
Specifies the number of time steps between writing restart files.
%
Note that a restart file is always written before the first (non-restarted)
time step and after completing the last time step.
%
With the default value of 0, no intermediate files will be written,
only the restarts before the first time step and after the final time step.
%
After successfully writing a new restart file, a
symbolic link will be created that points to the latest restart
file.
%
The name of this link will be:
\newline\indent
``\inp[nmlOutput::OutputDir]{Output\_Dir}\slash{Restart\_files}\slash{out.current.rst}'' 
\end{nmldescbody}

%==========================================================================

\begin{nmldescitem}
Restart_Frequency = -1.0  ! tstar
\end{nmldescitem}
\inp*[nmlOutput::RestartFrequency]{}[Restart\_Frequency][]

\begin{nmldescbody}
If $>0$, this value specifies the non-dimensional time
(\key[key:tstar]{\tstar}[t *@\tstar]) interval, \eref{EQN:tstar},
to write restart files.
%
This is complimentary to \inptext{Restart\_Interval}
and will write additional restart files.
%
After successfully writing a new restart file, a
symbolic link will be created that points to the latest restart
file.
%
The name of this link will be:
\newline\indent
``\inp[nmlOutput::OutputDir]{Output\_Dir}\slash{Restart\_files}\slash{out.current.rst}'' 
\end{nmldescbody}

%==========================================================================

\begin{nmldescitem}
Restart_Output_Uses_Host_Roots = T
\end{nmldescitem}
\inp*[nmlOutput::RestartOutputUsesHostRoots]{}[Restart\_Output\_Uses\_Host\_Roots][]

\begin{nmldescbody}
Indicates which MPI processes are involved in writing to the
restart file.
%
If true, then only the root processes on each host node will
preform the output.
%
Otherwise, all MPI processes will be involved in the output.

\NOTE
This is only applicable for parallel simulations run across
multiple physical computers, \forexample multiple nodes on
the NASA HECC supercomputers.
%
Only change this if you encounter MPI errors that occur when
trying to write restart files.
%
Changing this to false will most likely result in higher
I\slash O file contention and possibly increased memory usage.
\end{nmldescbody}

%==========================================================================

\begin{nmldescitem}
Restart_Output_Disabled = F
\end{nmldescitem}
\inp*[nmlOutput::RestartOutputDisabled]{}[Restart\_Output\_Disabled][]

\begin{nmldescbody}
If set to true, this option will prevent any instantaneous or time-averaged
restart files from being written.
%
This is intended for simulations that will complete with the allotted
resources and will not need to be restarted later.
\end{nmldescbody}

%==========================================================================

\begin{nmldescitem}
Output_Dir = "."
\end{nmldescitem}
\inp*[nmlOutput::OutputDir]{}[Output\_Dir][]

\begin{nmldescbody}
File path to the directory used for writing output files.

\NOTE
If one is available, it is recommended that the output directory be
located on a Lustre file system in order to get the best I\slash O
performance possible.
\end{nmldescbody}

%==========================================================================

\begin{nmldescitem}
Output_Interval = -5
\end{nmldescitem}
\inp*[nmlOutput::OutputInterval]{}[Output\_Interval][]

\begin{nmldescbody}
This controls the number of time steps between writing CGNS solution files.
The input value is interpreted as follows:
\\
\begin{tabular}{ @{\qquad} r @{ \quad } p{0.85\linewidth} @{} }
  $<0$ & For most situations, every
         $\inp[nmlTimestep::NumTimesteps]{Num\_Timesteps} /
           | \inptext{Output\_Interval} |$
         time steps.\\
  $=0$ & For most situations, every
         $\inp[nmlTimestep::NumTimesteps]{Num\_Timesteps}/50$ time steps.\\
  $>0$ & For all situations, every
         \inptext{Output\_Interval} time steps.
\end{tabular}
\\
The first two options result in a fixed number of output files
each time the code is run, assuming the code does not terminate
early.


\NOTE
For the following special cases, the input value is interpreted differently:
\begin{itemize}
%---------------------------------------------------------------------
\item 
If either
  $\inp[nmlOutput::OutputFrequency]{Output\_Frequency}>0$ or
  $\inp[nmlOutput::CGNSOutputDisabled]{CGNS\_Output\_Disabled}=\texttt{T}$,
the action associated with those inputs will take precedence
and this input will not be used.
%---------------------------------------------------------------------
\item
If running any of the transport problems (isentropic vortex,
density\slash entropy advection), this input value is used to
determine when a CGNS solution file will be written as follows:
\\
\begin{tabular}{ @{\qquad} r @{ \quad } p{0.75\linewidth} @{} }
  $<0$ & At the completion of every $|\inptext{Output\_Interval}|$
         flow-through periods.\\
  $=0$ & At the completion of every flow-through period.
\end{tabular}
%---------------------------------------------------------------------
\item
If running either the Taylor-Green vortex problem or the
Taylor-Couette problem, CGNS solution files will automatically be
written at certain specific values of
non-dimensional time (\key[key:tstar]{\tstar})
that are of interest
to each specific problem. If $\inptext{Output\_Interval}=0$,
only these problem specific CGNS solution files will be written;
otherwise ($\inptext{Output\_interval}\neq 0$) this behaves the
same as ``most situations'' in addition to writing the
problem-specific CGNS solution files.
%---------------------------------------------------------------------
\end{itemize}

\end{nmldescbody}

%==========================================================================

\begin{nmldescitem}
Output_Frequency = -1.0  ! tstar
\end{nmldescitem}
\inp*[nmlOutput::OutputFrequency]{}[Output\_Frequency][]

\begin{nmldescbody}
If $>0$, this value specifies the non-dimensional time
(\key[key:tstar]{\tstar}[t *@\tstar]) interval, \eref{EQN:tstar},
to write CGNS files.
%
\inp[nmlOutput::OutputInterval]{Output\_Interval} will be disabled.
%% Users need not know this detail:
%  by setting it to $\inp[nmlTimestep::NumTimesteps]{Num\_Timesteps}+1$.
%
\end{nmldescbody}

%==========================================================================

\begin{nmldescitem}
Output_Order = -3
\end{nmldescitem}
\inp*[nmlOutput::OutputOrder]{}[Output\_Order][]

\begin{nmldescbody}
Specifies the degree of the polynomial space used to represent the solution
in the CGNS solution files. This allows the ability to over-sample the
solution polynomial within each cell which provides a smoother solution 
when visualizing within your favorite post-processing software.

\NOTE
If using \inp[nmlOutput::CGNSHighOrderOutput]{CGNS\_High\_Order\_Output}, the
output order will be limited to a maximum of 4, since that is the
maximum order supported by the CGNS library.

\NOTE
If
$\inptext{Output\_Order} < \inp[nmlNumerics::SolutionOrder]{Solution\_Order}$,
the solution order will be used instead.
\end{nmldescbody}

%==========================================================================

\begin{nmldescitem}
Output_Summary_File = F
\end{nmldescitem}
\inp*[nmlOutput::OutputSummaryFile]{}[Output\_Summary\_File][]

\begin{nmldescbody}
Setting this to true will create the file ``summary.log'' when the
simulation successfully stops.
%
This file contains summarized information about the simulation like:
\begin{itemize}[noitemsep,nolistsep]
\item
  computational cost
\item
  solution order
\item
  number of cells
\item
  degrees of freedom
\item
  total time or \key[key:tstar]{\tstar}[t *@\tstar]
  elapsed during this simulation
\item
  final time or \key[key:tstar]{\tstar}[t *@\tstar]
  at the end of the simulation\\
  (\thatis combining the elapsed value with that from the restart file)
\item
  the current length of time averaging (if applicable)
\item
  and more
\end{itemize}

When running the channel flow test case using the internal grid generator,
it will also contain detailed grid information such as the number of
grid cells and degrees of freedom in each direction, as well as their
respective spacings.
\end{nmldescbody}

%==========================================================================

\begin{nmldescitem}
Dmat_Grad_Temp_Order = -1
\end{nmldescitem}
\inp*[nmlOutput::DmatGradTempOrder]{}[Dmat\_Grad\_Temp\_Order][]

\begin{nmldescbody}
If outputting the derivative of temperature, this is the quadrature order
used for computing temperature before projecting down to the solution order
and computing the derivative.

\NOTE
If $\inptext{Dmat\_Grad\_Temp\_Order} < 0$,
\inp[nmlNumerics::SolutionOrder]{Solution\_Order}
will be used instead.
\end{nmldescbody}

%==========================================================================

\begin{nmldescitem}
Loc_Output_Pts = 0
\end{nmldescitem}
\inp*[nmlOutput::LocOutputPts]{}[Loc\_Output\_Pts][]

\begin{nmldescbody}
Specifies the location of the \key{output nodal points}[points!output]
 within the grid cells when writing the CGNS solution files.
\\
\begin{tabular}{ @{\qquad} r @{ = } p{0.85\linewidth} @{} }
  0 & \key{Equidistant points}[points!equidistant] nodes \\
  1 & \key{Legendre-Gauss (Gauss) points}[points!Legendre-Gauss (Gauss)] nodes \\
  2 & \key{Legendre-Gauss-Lobatto (Lobatto) points}[points!Legendre-Gauss-Lobatto (Lobatto)] nodes
\end{tabular}

\NOTE
Using Legendre-Gauss nodes for the CGNS solution output files
will result in empty space between all cells because
those points do not extend all the way to the cell edge and
there is no connectivity information between cells.
%
This only affects your ability to visualize the solution
and does not affect the solution accuracy in any way.
\end{nmldescbody}

%==========================================================================

%\begin{nmldescitem}
%Continuous_Output = F
%\end{nmldescitem}
%\inp*[nmlOutput::ContinuousOutput]{}[Continuous\_Output][]
%
%\begin{nmldescbody}
%Specifies whether the output routine should attempt to create a smooth,
%continuous solution across all cells before writing to the
%CGNS solution file.
% 
%This option is intended for use with
%\inp[nmlOutput::LocOutputPts]{Loc\_Output\_Pts}=1 (Legendre-Gauss nodes),
%which have empty space between all cells because
%those points do not extend all the way to the cell edge and
%there is no connectivity information between cells.
%%
%When activated, additional output points will be added within
%each cell to fill this gap.
%%
%This will create CGNS solution files that are significantly larger in file
%size compared to the regular CGNS solution files.
%
%\NOTE
%This only works with a single MPI process.
%%
%So far, tests of this feature have been underwhelming with regards
%to creating a cleaner visualization of the solution for things
%like isosurfaces.
%
%\end{nmldescbody}

%==========================================================================

%\begin{nmldescitem}
%Interpolate_Before_Output_Variable = T
%\end{nmldescitem}
%\inp*[nmlOutput::InterpolateBeforeOutputVariable]{}[Interpolate\_Before\_Output\_Variable][]
%
%\begin{nmldescbody}
%Determines whether interpolation to the output points is done
%before or after the output variables are computed.
%
%\begin{tabular}{ @{} r @{ = } p{0.85\linewidth} @{} }
%  T & Interpolation is done before the output variable is computed.
%      Flow variables are interpolated from the solution points
%      to the output points, then output quantities are computed
%      at the output points.
%      This provides consistency between the various output variables
%      at the output points.\\
%  F & Interpolation is done after the output variable is computed.
%      Output quantities are computed from the flow variables
%      at the solution points, then interpolated to the output points.
%\end{tabular}
%
%In GFR, the conserved variables are represented as polynomials
%within each element.
%%
%However, derived quantities such as velocity, temperature, and
%pressure cannot be exactly represented as a polynomial
%because they involve nonlinear relations or fractions of polynomials.
%%
%If one interpolates the conserved variables to the output points
%and then computes the derived quantities, all the output variables
%at the output points will be self-consistent.
%%
%If one instead computes the derived quantities at the solution points
%and uses those approximate polynomials to interpolate onto the
%output points, the conserved and derived values at the output
%points will not be self-consistent.
%
%\end{nmldescbody}

%==========================================================================

\begin{nmldescitem}
CGNS_High_Order_Output = F
\end{nmldescitem}
\inp*[nmlOutput::CGNSHighOrderOutput]{}[CGNS\_High\_Order\_Output][]

\begin{nmldescbody}
Setting this to true will write the CGNS grid file (solution.grid.cgns)
using high-order cells. If not enabled (the default), the CGNS grid file
breaks each cell into linear sub-cells.

\NOTE
Visualizing high-order CGNS solution files requires specialized software,
such as Tecplot 2023 or newer.
\end{nmldescbody}

%==========================================================================

\begin{nmldescitem}
CGNS_Mixed_Output = T
\end{nmldescitem}
\inp*[nmlOutput::CGNSMixedOutput]{}[CGNS\_Mixed\_Output][]

\begin{nmldescbody}
This determines how cell connectivity is handled in high-order CGNS grid files.
%
If true, all grid cells will be combined into a single zone
regardless of cell geometry or order.
%
If false, each individual combination of cell geometry and order
will be written to its own zone.
%
Some post-processing software may prefer the latter format.
\end{nmldescbody}

%==========================================================================

\begin{nmldescitem}
CGNS_Output_Disabled = F
\end{nmldescitem}
\inp*[nmlOutput::CGNSOutputDisabled]{}[CGNS\_Output\_Disabled][]

\begin{nmldescbody}
Completely disables CGNS solutions from being output at the output interval.

On some systems with a large number of processors,
the output of CGNS files can be slow.
%
It can be advantageous to instead output restart files,
and later use
\inp[nmlUtility::ConvertRestartToCGNS]{Convert\_Restart\_To\_CGNS}
in the \inptext{nml\_utility} namelist block of \sref{SEC:nml_utility}
to have GFR batch convert those restart files to CGNS.
%
This post-processing step can be done using a minimal number
of processors, because most of the work is simple file I\slash O.

\end{nmldescbody}

%==========================================================================

\begin{nmldescitem}
Probes_Collection_Interval = 0
\end{nmldescitem}
\inp*[nmlOutput::ProbesCollectionInterval]{}[Probes\_Collection\_Interval][]

\begin{nmldescbody}
If $>0$, this value specifies the iteration interval used to collect
time-history data for each of the probes. The information collected at this
interval is saved to a buffer and only written out after completing the number
of collections specified by the input
\inp[nmlOutput::ProbesWriteInterval]{Probes\_Write\_Interval}.
%
With the default value of 0, no time-history probe data will be collected or
output.
\end{nmldescbody}

%==========================================================================

\begin{nmldescitem}
Probes_Write_Interval = 100
\end{nmldescitem}
\inp*[nmlOutput::ProbesWriteInterval]{}[Probes\_Write\_Interval][]

\begin{nmldescbody}
This value specifies how often to write out the buffer used to collect
time-history data for the probes.

\NOTE
This interval is relative to the number of collections completed, and is not
relative to the simulation iteration.
\end{nmldescbody}

%==========================================================================

\begin{nmldescitem}
Point_Probes(1)%xyz(1:3) = [-999.,-999.,-999.]  ! input grid units
\end{nmldescitem}
\inp*[nmlOutput::PointProbes]{}[Point\_Probes][]

\begin{nmldescbody}
This specifies the location of a single point \key[key:probes]{probe}[probes] to
collect time-history data.
%
The values specified for this input are assumed to be in the units of the input
mesh file (\thatis \emph{grid units} as discussed in \sref{SEC:Grid_Scaling}).
%
Up to 200 point probes may be used.
%
Probes set to the default location are ignored.
%
For more information on the point probe files, see
\sref{SEC:Files:Probes_Output}.

\NOTE
The solution point closest to these coordinates will be used as the source for
the probe data.

\NOTE
If an individual solution point is identified as the source of a point probe by
more than one point\slash rake\slash cell probe, the duplicate probes will be
removed and only a single probe will be collected for that solution point.
%
This will make it seem like there are missing probe files.
\end{nmldescbody}

%==========================================================================

\begin{nmldescitem}
Rake_Probes(1)%Num_Probes   = 0
Rake_Probes(1)%xyz_beg(1:3) = [-999.,-999.,-999.]  ! input grid units
Rake_Probes(1)%xyz_end(1:3) = [-999.,-999.,-999.]  ! input grid units
\end{nmldescitem}
\inp*[nmlOutput::RakeProbes]{}[Rake\_Probes][]

\begin{nmldescbody}
These inputs provide the details necessary to create a rake of point
\key[key:probes]{probes}[probes].
%
The number of probes to create within rake $i$ is specified using the input
\inptext{Rake\_Probes($i$)\%Num\_Probes}.
%
\inptext{Rake\_Probes($i$)\%xyz\_beg(1:3)} and
\inptext{Rake\_Probes($i$)\%xyz\_end(1:3)}
specify the coordinates for the start and end points of the rake, with both of
these coordinate inputs assumed to be in the units of the mesh file (\thatis
\emph{grid units} as discussed in \sref{SEC:Grid_Scaling}).

In the case that $\inptext{Num\_Probes}=1$, the single probe created by the
rake will be located halfway between the start and end points for the rake.
%
Otherwise, probes will be located at both the start and end points for the
rake, and $\inptext{Num\_Probes}-2$ probes will be equally spaced along the
line between the start and end points.

Up to 10 rake probes may be used, with no limit on the number of probes per rake.
%
Rakes with $\inptext{Num\_Probes}=0$ or $\inptext{xyz\_beg}=\inptext{xyz\_end}$
are ignored, and no time-history probe data will be collected or
output for the rake.
%
For more information on the rake probe files, see
\sref{SEC:Files:Probes_Output}.

\NOTE
For each individual rake probe, the solution point closest to the probe
coordinates will be used as the source for the probe data.

\NOTE
If an individual solution point is identified as the source of a point probe by
more than one point\slash rake\slash cell probe, the duplicate probes will be
removed and only a single probe will be collected for that solution point.
%
This will make it seem like there are missing probe files.

\NOTE
It should be warned that too many probes (1000+) may have a large negative
impact on code and file system performance.
\end{nmldescbody}

%==========================================================================

\begin{nmldescitem}
Cell_Probes(1)%xyz(1:3) = [-999.,-999.,-999.]  ! input grid units
\end{nmldescitem}
\inp*[nmlOutput::CellProbes]{}[Cell\_Probes][]

\begin{nmldescbody}
This input provides the ability to create a set of point probes for every
solution point within a grid cell.
%
For the grid cell that contains the coordinates of this input, individual point
\key[key:probes]{probes}[probes] will be created for every solution point within
that grid cell, resulting in time-history data for the entire cell.
%
The values specified for this input are assumed to be in the units of the input
mesh file (\thatis \emph{grid units} as discussed in \sref{SEC:Grid_Scaling}).
%
Up to 10 cell probes may be used.
%
Probes set to the default location are ignored.
%
For more information on the cell probe files, see
\sref{SEC:Files:Probes_Output}.

\NOTE
If an individual solution point is identified as the source of a point probe by
more than one point\slash rake\slash cell probe, the duplicate probes will be
removed and only a single probe will be collected for that solution point.
%
This will make it seem like there are missing probe files.

\NOTE
It should be warned that too many cell probes may have a large negative impact
on code and file system performance.
%
A single cell probe for a \PO[3] hexahedral cell will result in 64 separate
point probes, and 125 probes for a \PO[4] hexahedral cell.
\end{nmldescbody}

%==========================================================================

\begin{nmldescitem}
Profile_IO_All = F
\end{nmldescitem}
\inp*[nmlOutput::ProfileIOAll]{}[Profile\_IO\_All][]

\begin{nmldescbody}
If true, sets both
  \inp[nmlOutput::ProfileIOCGNS]{Profile\_IO\_CGNS} and
  \inp[nmlOutput::ProfileIORestart]{Profile\_IO\_Restart}
to true.
%
If false, it has no effect on the other input parameters.
\end{nmldescbody}

%==========================================================================

\begin{nmldescitem}
Profile_IO_Restart = F
\end{nmldescitem}
\inp*[nmlOutput::ProfileIORestart]{}[Profile\_IO\_Restart][]

\begin{nmldescbody}
Activates timers to measure how long it takes to write each restart or
time-averaged restart file, and outputs the results to standard output.
\end{nmldescbody}

%==========================================================================

\begin{nmldescitem}
Profile_IO_CGNS = F
\end{nmldescitem}
\inp*[nmlOutput::ProfileIOCGNS]{}[Profile\_IO\_CGNS][]

\begin{nmldescbody}
Activates timers to measure how long it takes to write each CGNS solution file,
and outputs the results to standard output.
\end{nmldescbody}

%==========================================================================

\begin{nmldescitem}
Profile_IO_CGNS_Subs = F
\end{nmldescitem}
\inp*[nmlOutput::ProfileIOCGNSSubs]{}[Profile\_IO\_CGNS\_Subs][]

\begin{nmldescbody}
Activates timers to measure how long it takes to complete each individual call
to a function from the CGNS library, and outputs the results to standard output.
%Outputs timing information for each individual call to a CGNS function.
\end{nmldescbody}

%==========================================================================

\begin{nmldescitem}
Lustre_Stripe_Count = 64
\end{nmldescitem}
\inp*[nmlOutput::LustreStripeCount]{}[Lustre\_Stripe\_Count][]

\begin{nmldescbody}
Specifies the approximate number of stripes to use when writing a CGNS
solution file or a restart file.

\NOTE
This is only used if \inp[nmlOutput::OutputDir]{Output\_Dir}
is on a Lustre file system.
%
The code will take this value and find the closest integer value satisfying:
$\textrm{mod} \left( \textsl{ncpu}
                   , \inptext{Lustre\_Stripe\_Count}
              \right) = 0$
\end{nmldescbody}

%==========================================================================

\begin{nmldescitem}
Lustre_Stripe_Size = "4m"
\end{nmldescitem}
\inp*[nmlOutput::LustreStripeSize]{}[Lustre\_Stripe\_Size][]

\begin{nmldescbody}
Specifies the number of bytes to store on each stripe before moving to the
next stripe.

\NOTE
This is only used if \inp[nmlOutput::OutputDir]{Output\_Dir}
is on a Lustre file system.
%
The value of \linebreak[4]\inptext{Lustre\_Stripe\_Size}
must be a multiple of 65,536 bytes (64 KB), and the suffixes
`k', `m', or `g' can be used to specify units of
KB, MB, or GB, respectively.

\EXAMPLE
$\textrm{``4m''} \Rightarrow 4~\textrm{MB}
                 \Rightarrow 4,194,304~\textrm{bytes}$
\end{nmldescbody}

%==========================================================================


\subsection{nml\_output\_variables}
\label{SEC:nml_output_variables}
\inp*[nmlOutputVariables]{}[nml\_output\_variables][]

This \hyperref[SEC:User_Input]{namelist block}
is used to specify which variables to output
to the CGNS files.

\begin{nmlblocklst}[
,label={LST:nml_output_variables}
,caption={Namelist nml\_output\_variables.}
]
&nml_output_variables
  Conserved_Variables            = F
  Momentum                       = F
  Vorticity                      = F
  Cylindrical_Velocity           = F
  Cylindrical_Coordinates        = F
  Relative_Velocity              = F
  Relative_Conditions            = F
  Thermo_Variables               = F
  Temperature                    = F
  Gradients                      = F
  Density_Gradient               = F
  Density_Gradient_Magnitude     = F
  Pressure_Gradient              = F
  Pressure_Gradient_Magnitude    = F
  Pressure_Dmat_Gradient         = F
  Temperature_Gradient           = F
  Temperature_Gradient_Magnitude = F
  Temperature_Dmat_Gradient      = F
  Schlieren                      = F
  Modified_Schlieren             = F
  Transport_Variables            = F
  Stagnation_Variables           = F
  Stagnation_Pressure            = F
  Stagnation_Temperature         = F
  Marker_Variables               = F
  Qcriterion                     = F
  Grid_Details                   = F
  Local_Timestep                 = F
  Species_Mass_Fraction          = T
  Species_Density                = F
  Species_Mole_Fraction          = F
  Extra_Species_Variables        = F
  Debug_Species_Variables        = F
  Grdnt_Species_Variables        = F
  Local_Filter_Dissipation       = F
  CellAve_Density                = F
  Exact_Density                  = F
  Exact_Pressure                 = F
/
\end{nmlblocklst}

%% For future reference:
%%   The default output variable list is in
%%   subroutine get_cgns_sol_variables of write_cgns_mod.f90.

\FINALFIX{\newpage}
The following variables are automatically output to the CGNS files:
\\
  \ov{Density},
  \ov{VelocityX},
  \ov{VelocityY},
  \ov{VelocityZ},
  \ov{Pressure}, and
  \ov{Mach}.

\vspace{1.0ex}
\noindent\textbf{NOTES:}\\
\begin{itemize}
  \item
    The grid coordinates are output in the \key[key:GridUnits]{grid units}
    of the input grid file.  For example, say your grid file has dimensions of
    inches. You will need to include the input
    $\inp[nmlGrid::GridScalingFactor]{Grid\_Scaling\_Factor}=0.0254$ to convert
    the grid file to meters for it to be used correctly in the solver. However,
    this conversion will be undone when creating the CGNS solution files, which
    will contain coordinate dimensions of inches just like the grid file.
  \item
    All variables are output in SI units (except for the coordinate variables as
    discussed in the previous bullet).
  \item
    Due to the first two bullets, if using post-processing software to compute
    integrated quantities, care needs to be taken to ensure the coordinate
    variables and variable(s) being integrated are in compatible units.
  \item
    In the following descriptions, any vector quantities that are output for 1D
    and 2D cases will exclude vector components from higher spatial dimensions.
    For example, \ov{VelocityY} and \ov{VelocityZ} will not be included for a 1D
    flow, and for a gradient variable from a 2D simulation, the
    \ov{[VariableName]GradientZ} variables would not be included.
  \item
    If multiple inputs add the same variable to the list of output variables,
    only one instance of the variable will be written, \thatis any duplicates in
    the list of output variables will be removed.
  \item
    The names of the variables written to the CGNS solution files attempt to
    follow the variable naming conventions of the CGNS standard, \forexample the
    $x$-coordinate variable will have the name \ov{CoordinateX} instead of just
    \quoterm{\textsf{X}}.
\end{itemize}

\devNotes{
CHANGECODE:
Default variables can not be disabled.
Instead of a hard-coded default list, consider adding individual options
to the namelist above with defaults:
\\\indent\raggedright
  \inptext{Density = T},
  \inptext{Velocity = T},
  \inptext{Pressure = T}, and
  \inptext{Mach = T}
}
%==========================================================================

\begin{nmldescitem}
Conserved_Variables = F
\end{nmldescitem}
\inp*[nmlOutputVariables::ConservedVariables]{}[Conserved\_Variables][]

\begin{nmldescbody}
Setting this to true adds the following variables to the CGNS solution files:
\\
  \ov{Density},
  \ov{MomentumX},
  \ov{MomentumY},
  \ov{MomentumZ}, and
  \ov{EnergyStagnationDensity} (\thatis $\rho e_t$).
\end{nmldescbody}

%==========================================================================

\begin{nmldescitem}
Momentum = F
\end{nmldescitem}
\inp*[nmlOutputVariables::Momentum]{}[Momentum][]

\begin{nmldescbody}
Setting this to true adds the following variables to the CGNS solution files:
\\
  \ov{MomentumX},
  \ov{MomentumY}, and
  \ov{MomentumZ}.
\end{nmldescbody}

%==========================================================================

\begin{nmldescitem}
Vorticity = F
\end{nmldescitem}
\inp*[nmlOutputVariables::Vorticity]{}[Vorticity][]

\begin{nmldescbody}
Setting this to true adds the following variables to the CGNS solution files:
\\\raggedright
  \ov{VorticityX},
  \ov{VorticityY},
  \ov{VorticityZ},
  \ov{VorticityMagnitude},
  \ov{DivergenceOfVelocity},
  \ov{StrainRateMagnitude},
  \ov{Qcriterion}, and
  \ov{Qcriterion\_Positive}.
\end{nmldescbody}

%==========================================================================

\begin{nmldescitem}
Cylindrical_Velocity = F
\end{nmldescitem}
\inp*[nmlOutputVariables::CylindricalVelocity]{}[Cylindrical\_Velocity][]

\begin{nmldescbody}
Setting this to true adds the following variables to the CGNS solution files:
\\
  \ov{VelocityR} and
  \ov{VelocityTheta}.
\\
If either
\inp[nmlOutputVariables::RelativeVelocity]{Relative\_Velocity} or
\inp[nmlOutputVariables::RelativeConditions]{Relative\_Conditions}
is also true, then
the following variables will also be added to the CGNS solution files:
\\
  \ov{RelativeVelocityR} and
  \ov{RelativeVelocityTheta}.
\end{nmldescbody}

%==========================================================================

\begin{nmldescitem}
Cylindrical_Coordinates = F
\end{nmldescitem}
\inp*[nmlOutputVariables::CylindricalCoordinates]{}[Cylindrical\_Coordinates][]

\begin{nmldescbody}
Setting this to true adds the following variables to the CGNS \emph{grid} file:
\\
  \ov{CoordinateAxial},
  \ov{CoordinateR}, and
  \ov{CoordinateTheta}.
\\
This will also automatically set 
\inp[nmlOutputVariables::CylindricalVelocity]{Cylindrical\_Velocity} to true.

\noindent\NOTE
The axial coordinate direction will be set to the $x$-axis, \thatis
$\ov{CoordinateAxial}=\ov{CoordinateX}$, unless:
  \begin{itemize}
    \item 
      $\inp[nmlReference::AoAAlpha]{AoA\_Alpha}=\pm90$\textdegree{} and
      $\inp[nmlReference::AoABeta]{AoA\_Beta}=0$\textdegree{}:
      $\quad\Rightarrow\quad \ov{CoordinateAxial}=\ov{CoordinateY}$
    \item 
      $\inptext{AoA\_Alpha}=0$\textdegree{} and
      $\inptext{AoA\_Beta}=\pm90$\textdegree{}:
      $\quad\Rightarrow\quad \ov{CoordinateAxial}=\ov{CoordinateZ}$
  \end{itemize}
The other two cylindrical coordinate variables are determined by the axial
direction and the right-hand rule.

\noindent\NOTE
These variables will only be written to the CGNS grid file. These variables will
be added to each of the solution files by using CGNS links.
\end{nmldescbody}

%==========================================================================

%\FINALFIX{\newpage}
\begin{nmldescitem}
Relative_Velocity = F
\end{nmldescitem}
\inp*[nmlOutputVariables::RelativeVelocity]{}[Relative\_Velocity][]

\begin{nmldescbody}
Setting this to true adds the following variables to the CGNS solution files:
\\
  \ov{RelativeVelocityX},
  \ov{RelativeVelocityY}, and
  \ov{RelativeVelocityZ}.
\end{nmldescbody}

%==========================================================================

\begin{nmldescitem}
Relative_Conditions = F
\end{nmldescitem}
\inp*[nmlOutputVariables::RelativeConditions]{}[Relative\_Conditions][]

\begin{nmldescbody}
Setting this to true adds the following variables to the CGNS solution files:
\\
  \ov{RelativeVelocityX},
  \ov{RelativeVelocityY},
  \ov{RelativeVelocityZ},
\\
  \ov{RelativeMomentumX},
  \ov{RelativeMomentumY},
  \ov{RelativeMomentumZ}, and
  \ov{RelativeMach}.
\end{nmldescbody}

%==========================================================================

\begin{nmldescitem}
Thermo_Variables = F
\end{nmldescitem}
\inp*[nmlOutputVariables::ThermoVariables]{}[Thermo\_Variables][]

\begin{nmldescbody}
Setting this to true adds the following variables to the CGNS solution files:
\\
  \ov{Pressure},
  \ov{Temperature},
  \ov{Enthalpy},
  \ov{SpecificHeatPressure} (\thatis $c_p$), and
  \ov{Entropy}.
\end{nmldescbody}

%==========================================================================

\begin{nmldescitem}
Temperature = F
\end{nmldescitem}
\inp*[nmlOutputVariables::Temperature]{}[Temperature][]

\begin{nmldescbody}
Setting this to true adds the following variable to the CGNS solution files:
  \ov{Temperature}.
\end{nmldescbody}

%==========================================================================

\begin{nmldescitem}
Gradients = F
\end{nmldescitem}
\inp*[nmlOutputVariables::Gradients]{}[Gradients][]

\begin{nmldescbody}
Setting this to true will also include in the CGNS solution file the
directional gradient components for any of the following variables that are
also set to be included in the CGNS solution file:
  \ov{Density},
  \ov{Velocity},
  \ov{EnergyStagnationDensity},
  \ov{Temperature},
  \ov{Momentum},
  \ov{MassFraction}, and
  \ov{SpeciesDensity}.

The name of the gradient variables will be the name of the
non-gradient variable with
  \ov{GradientX}, \ov{GradientY}, and \ov{GradientZ}
appended on to the variable name,
\forexample
\\
  \ov{MomentumYGradientZ}.
\end{nmldescbody}

%==========================================================================

\begin{nmldescitem}
Density_Gradient = F
\end{nmldescitem}
\inp*[nmlOutputVariables::DensityGradient]{}[Density\_Gradient][]

\begin{nmldescbody}
Setting this to true will add the directional gradient components of density
to the CGNS solution files, with the names of the variables being
\ov{DensityGradientX}, \ov{DensityGradientY}, and \ov{DensityGradientZ}.
\end{nmldescbody}

%==========================================================================

\FINALFIX{\newpage}
\begin{nmldescitem}
Density_Gradient_Magnitude = F
\end{nmldescitem}
\inp*[nmlOutputVariables::DensityGradientMagnitude]{}[Density\_Gradient\_Magnitude][]

\begin{nmldescbody}
Setting this to true will add the magnitude of the density gradient
to the CGNS solution files under the variable name \ov{GradDensityMagnitude}.
\end{nmldescbody}

%==========================================================================

\begin{nmldescitem}
Pressure_Gradient = F
\end{nmldescitem}
\inp*[nmlOutputVariables::PressureGradient]{}[Pressure\_Gradient][]

\begin{nmldescbody}
Setting this to true will add the directional gradient components of pressure
to the CGNS solution files, with the names of the variables being
\ov{PressureGradientX}, \ov{PressureGradientY}, and \ov{PressureGradientZ}.

\noindent\NOTE
This will use the chain rule to compute a pressure gradient that is consistent
with the gradients of the conserved variables.
\end{nmldescbody}

%==========================================================================

\begin{nmldescitem}
Pressure_Gradient_Magnitude = F
\end{nmldescitem}
\inp*[nmlOutputVariables::PressureGradientMagnitude]{}[Pressure\_Gradient\_Magnitude][]

\begin{nmldescbody}
Setting this to true will add the magnitude of the pressure gradient
to the CGNS solution files under the variable name \ov{GradPresMagnitude}.

\noindent\NOTE
This will use the chain rule to compute a pressure gradient that is consistent
with the gradients of the conserved variables.
\end{nmldescbody}

%==========================================================================

\begin{nmldescitem}
Pressure_Dmat_Gradient = F
\end{nmldescitem}
\inp*[nmlOutputVariables::PressureDmatGradient]{}[Pressure\_Dmat\_Gradient][]

\begin{nmldescbody}
Setting this to true will use the derivative matrices to compute the gradient
of the pressure and add the resulting directional gradient components and the
gradient magnitude to the CGNS solution files, with the names of the variables
being \ov{PressureDmatrixX}, \ov{PressureDmatrixY}, \ov{PressureDmatrixZ}, and
\ov{PressureDmatrixMagnitude}.

\noindent\NOTE
The pressure field in a cell is not in the solution polynomial space, so the
resulting pressure gradients will most likely contain some amount of aliasing
errors since derivative matrices are only exact for a function within the
solution polynomial space.
\end{nmldescbody}

%==========================================================================

\begin{nmldescitem}
Temperature_Gradient = F
\end{nmldescitem}
\inp*[nmlOutputVariables::TemperatureGradient]{}[Temperature\_Gradient][]

\begin{nmldescbody}
Setting this to true will add the directional gradient components of temperature
to the CGNS solution files, with the names of the variables being
\ov{TemperatureGradientX}, \ov{TemperatureGradientY}, and \ov{TemperatureGradientZ}.

\noindent\NOTE
This will use the chain rule to compute a temperature gradient that is consistent
with the gradients of the conserved variables.
\end{nmldescbody}

%==========================================================================

\begin{nmldescitem}
Temperature_Gradient_Magnitude = F
\end{nmldescitem}
\inp*[nmlOutputVariables::TemperatureGradientMagnitude]{}[Temperature\_Gradient\_Magnitude][]

\begin{nmldescbody}
Setting this to true will add the magnitude of the temperature gradient
to the CGNS solution files under the variable name \ov{GradTempMagnitude}.

\noindent\NOTE
This will use the chain rule to compute a temperature gradient that is consistent
with the gradients of the conserved variables.
\end{nmldescbody}

%==========================================================================

\FINALFIX{\newpage}
\begin{nmldescitem}
Temperature_Dmat_Gradient = F
\end{nmldescitem}
\inp*[nmlOutputVariables::TemperatureDmatGradient]{}[Temperature\_Dmat\_Gradient][]

\begin{nmldescbody}
Setting this to true will use the derivative matrices to compute the gradient of
the temperature and add the resulting directional gradient components and the
gradient magnitude to the CGNS solution files, with the names of the variables
being \ov{TemperatureDmatrixX}, \ov{TemperatureDmatrixY},
\ov{TemperatureDmatrixZ}, and \ov{TemperatureDmatrixMagnitude}.

\noindent\NOTE
The temperature field in a cell is not in the solution polynomial space, so the
resulting temperature gradients will most likely contain some amount of aliasing
errors since derivative matrices are only exact for a function within the
solution polynomial space.
\end{nmldescbody}

%==========================================================================

\begin{nmldescitem}
Schlieren = F
\end{nmldescitem}
\inp*[nmlOutputVariables::Schlieren]{}[Schlieren][]

\begin{nmldescbody}
Setting this to true will compute a numerical Schlieren variable to include
in the CGNS solution files under the variable name \ov{Schlieren}.

\noindent\NOTE
The numerical Schlieren is simply the magnitude of the density gradient, \thatis
$\nabla\rho$.
%
This is not to be confused with what is sometimes referred to as
``shadowgraph'' which is the Laplacian of density, \thatis $\nabla^2\rho$.
\end{nmldescbody}

%==========================================================================

\begin{nmldescitem}
Modified_Schlieren = F
\end{nmldescitem}
\inp*[nmlOutputVariables::ModifiedSchlieren]{}[Modified\_Schlieren][]

\begin{nmldescbody}
Setting this to true will add the variable \ov{ModifiedSchlieren} to the
CGNS solution files, which is a $\log_{10}$ scaling of the numerical Schlieren
variable, \thatis $\log_{10}\left(1+\nabla\rho\right)$ where the offset by $1$
is simply to avoid $\log_{10}(0)$ when the density gradient is $0$.
\end{nmldescbody}

%==========================================================================

\begin{nmldescitem}
Transport_Variables = F
\end{nmldescitem}
\inp*[nmlOutputVariables::TransportVariables]{}[Transport\_Variables][]

\begin{nmldescbody}
Setting this to true adds the following variables to the CGNS solution files:
  \ov{ViscosityMolecular} and
  \ov{ThermalConductivity}.
\\
If using a multi-component gas mixture and
\inp[nmlOutputVariables::ExtraSpeciesVariables]{Extra\_Species\_Variables}
is also true, then the variable
  \ov{Prandtl}
giving the mixture Prandtl number will also be included in the
CGNS solution file.
\end{nmldescbody}

%==========================================================================

\begin{nmldescitem}
Stagnation_Variables = F
\end{nmldescitem}
\inp*[nmlOutputVariables::StagnationVariables]{}[Stagnation\_Variables][]

\begin{nmldescbody}
Setting this to true adds the following variables to the CGNS solution files:
  \ov{PressureStagnation},
  \ov{TemperatureStagnation}, and
  \ov{EnthalpyStagnation}.
\end{nmldescbody}

%==========================================================================

\begin{nmldescitem}
Stagnation_Pressure = F
\end{nmldescitem}
\inp*[nmlOutputVariables::StagnationPressure]{}[Stagnation\_Pressure][]

\begin{nmldescbody}
Setting this to true adds the following variable to the CGNS solution files:
  \ov{PressureStagnation}.
\end{nmldescbody}

%==========================================================================

\begin{nmldescitem}
Stagnation_Temperature = F
\end{nmldescitem}
\inp*[nmlOutputVariables::StagnationTemperature]{}[Stagnation\_Temperature][]

\begin{nmldescbody}
Setting this to true adds the following variable to the CGNS solution files:
  \ov{TemperatureStagnation}.
\end{nmldescbody}

%==========================================================================

\FINALFIX{\newpage}
\begin{nmldescitem}
Marker_Variables = F
\end{nmldescitem}
\inp*[nmlOutputVariables::MarkerVariables]{}[Marker\_Variables][]

\begin{nmldescbody}
If using \inp[nmlLimiter::MarkerOption]{Marker\_Option},
the following variables will be added to the CGNS solution files:
\\
%\begin{tabular}{ @{\quad} c @{ \quad } p{0.85\linewidth} @{} }
\begin{tabular}{ @{\quad} c l l l @{} }
  \inptext{Marker\_Option} & \quad & \multicolumn{2}{c}{Output Variables} \\
  \cline{1-1}                         \cline{3-4}
  1 && \texttt{Resolution\_Marker}& \texttt{Scaled\_Resolution\_Marker} \\
  2 && \texttt{Jump\_Marker}      & \texttt{Scaled\_Jump\_Marker} \\
  3 && \texttt{MLP\_Marker}       & \texttt{Scaled\_MLP\_Marker}
\end{tabular}
\\
The variable specified by
  \inp[nmlLimiter::MarkerVariable]{Marker\_Variable}
will also be output.
\end{nmldescbody}

%==========================================================================

\begin{nmldescitem}
Qcriterion = F
\end{nmldescitem}
\inp*[nmlOutputVariables::Qcriterion]{}[Qcriterion][]

\begin{nmldescbody}
Setting this to true adds the following variables to the CGNS solution files:
\\
  \ov{Qcriterion} and
  \ov{Qcriterion\_Positive}.

For the variable \ov{Qcriterion}, the computed Q-criterion is
normalized then rescaled such that the average Q-criterion is 0.0,
and the values $\pm 1$ respectively represent $\pm 1/100$ of a
standard deviation from the average.
%
The variable \ov{Qcriterion\_Positive} contains the true values
of the Q-criterion, but only the positive parts, with the
negative parts set to zero.
\end{nmldescbody}

%==========================================================================

\begin{nmldescitem}
Grid_Details = F
\end{nmldescitem}
\inp*[nmlOutputVariables::GridDetails]{}[Grid\_Details][]

\begin{nmldescbody}
Setting this to true adds the following variables to the CGNS solution files:
\\
  \ov{Partition},
  \ov{CellVolume}, and
  \ov{CellLengthScale}.

The variable \ov{Partition} will contain the partition ID for each grid cell.
%
The variable \ov{CellVolume} will contain the area or volume of each grid
cell depending on whether the grid is 2D or 3D, respectively.
%
The variable \ov{CellLengthScale} will contain an approximate length scale for
each grid cell which is computed as either the square-root of the cell area
for a 2D grid or the cube-root of the cell volume for a 3D grid.
\end{nmldescbody}

%==========================================================================

\begin{nmldescitem}
Local_Timestep = F
\end{nmldescitem}
\inp*[nmlOutputVariables::LocalTimestep]{}[Local\_Timestep][]

\begin{nmldescbody}
Setting this to true adds the following variables to the CGNS solution files:
\\
  \ov{LocalTimestep} and
  \ov{LocalTimestep\_Scaled}.

The variable \ov{LocalTimestep} will contain the physical time step size
computed for each grid cell, whereas the variable \ov{LocalTimestep\_Scaled}
will contain that same value normalized by the minimum time step size across
all grid cells.
\end{nmldescbody}

%==========================================================================

\begin{nmldescitem}
Species_Mass_Fraction = T
\end{nmldescitem}
\inp*[nmlOutputVariables::SpeciesMassFraction]{}[Species\_Mass\_Fraction][]

\begin{nmldescbody}
If using a multi-component gas mixture, setting this to true adds the
mass fractions for each of the gas species to the CGNS solution files.
%
The variable name will be \ov{MassFraction} with the species' chemical
formula appended to the end, \forexample \ov{MassFractionH2O}.
\end{nmldescbody}

%==========================================================================

\begin{nmldescitem}
Species_Density = F
\end{nmldescitem}
\inp*[nmlOutputVariables::SpeciesDensity]{}[Species\_Density][]

\begin{nmldescbody}
If using a multi-component gas mixture, setting this to true adds the
individual density for each of the gas species to the CGNS solution files.
%
The variable name will be \ov{SpeciesDensity} with the species' chemical
formula appended to the end, \forexample \ov{SpeciesDensityH2O}.
\end{nmldescbody}

%==========================================================================

\begin{nmldescitem}
Species_Mole_Fraction = F
\end{nmldescitem}
\inp*[nmlOutputVariables::SpeciesMoleFraction]{}[Species\_Mole\_Fraction][]

\begin{nmldescbody}
If using a multi-component gas mixture, setting this to true adds the
mole fractions for each of the gas species to the CGNS solution files.
%
The variable name will be \ov{MoleFraction} with the species' chemical
formula appended to the end, \forexample \ov{MoleFractionH2O}.
\end{nmldescbody}

%==========================================================================

\begin{nmldescitem}
Extra_Species_Variables = F
\end{nmldescitem}
\inp*[nmlOutputVariables::ExtraSpeciesVariables]{}[Extra\_Species\_Variables][]

\begin{nmldescbody}
If using a multi-component gas mixture, setting this to true adds the
following variables to the CGNS solution files:
\\\raggedright
  \ov{SpeciesDensity},
  \ov{Enthalpy},
  \ov{StagnationEnthalpy},
  \ov{SpecificHeatRatio} (\thatis mixture $\gamma$),
  \ov{SpecificHeatPressure} (\thatis mixture $c_p$),
  \ov{IdealGasConstant} (\thatis mixture $R_{\mathrm{gas}}$),
  \ov{MomentumX},
  \ov{MomentumY},
  \ov{MomentumZ},
  \ov{EnergyStagnationDensity},
  \ov{PressureStagnation}, and
  \ov{TemperatureStagnation}.

Each gas component will have its own \ov{SpeciesDensity} variable
with its respective chemical formula appended to the end,
\forexample \ov{SpeciesDensityCO2}.
\end{nmldescbody}

%==========================================================================

\begin{nmldescitem}
Debug_Species_Variables = F
\end{nmldescitem}
\inp*[nmlOutputVariables::DebugSpeciesVariables]{}[Debug\_Species\_Variables][]

\begin{nmldescbody}
If using a multi-component gas mixture, setting this to true will output the
individual contribution of each species to various mixture thermodynamic
properties, \forexample enthalpy, $c_p$, etc.

\NOTE
These variables are intended to aid developers in debugging the
thermodynamic properties of gas mixtures.
\end{nmldescbody}

%==========================================================================

\begin{nmldescitem}
Grdnt_Species_Variables = F
\end{nmldescitem}
\inp*[nmlOutputVariables::GrdntSpeciesVariables]{}[Grdnt\_Species\_Variables][]

\begin{nmldescbody}
If using a multi-component gas mixture, this is essentially the same as
\inp[nmlOutputVariables::Gradients]{Gradients}, but it also outputs
the gradients of the species mass fractions and densities.
\end{nmldescbody}

%==========================================================================

%\AddEntry[nmlOutputVariables::LocalFilterDissipation]{nml\_output\_Variables}{Local\_Filter\_Dissipation}

\begin{nmldescitem}
Local_Filter_Dissipation = F
\end{nmldescitem}
\inp*[nmlOutputVariables::LocalFilterDissipation]{}[Local\_Filter\_Dissipation][]

\begin{nmldescbody}
This output can be used to assess filter dissipation.
%\vspace*{0.5em}%
%\FIXME{
%\begin{enumerate}
%\item
%What do these variables represent? Units?
%\item
%Proof this description and verify the option is coded.
%\end{enumerate}
%}
\end{nmldescbody}

%==========================================================================

\begin{nmldescitem}
CellAve_Density = F
\end{nmldescitem}
\inp*[nmlOutputVariables::CellAveDensity]{}[CellAve\_Density][]

\begin{nmldescbody}
Setting this to true computes the cell average of the instantaneous density
polynomial to include in the CGNS solution files under the variable name
\ov{CellAveDensity}.
\end{nmldescbody}

%==========================================================================

\begin{nmldescitem}
Exact_Density = F
Exact_Pressure = F
\end{nmldescitem}
\inp*[nmlOutputVariables::ExactDensity]{}[Exact\_Density][]
\inp*[nmlOutputVariables::ExactPressure]{}[Exact\_Pressure][]

\begin{nmldescbody}
For cases that have an analytical solution, the exact density and pressure
can be computed at the output points and written to the solution file
for direct comparison with the numerical solution.
%
Cases that support this output include the following:
\\
\begin{tabular}{cl}
  \inp[nmlInitialization::Initialization]{Initialization} & Description \\
  23 & 1D Sod shock tube \\
  26 & 1D Jiang-Shu density pulse test
\end{tabular}
\end{nmldescbody}




\subsection{nml\_sem}
\label{SEC:nml_sem}
\inp*[nmlSEM]{}[nml\_sem][]

This \hyperref[SEC:User_Input]{namelist block}
is used to specify the global controls of the Synthetic Eddy Method.
%
Additional controls for specifying the SEM on individual boundaries
are contained in the \inptext{nml\_bc} namelist block
described in \sref{SEC:nml_bc}.

\begin{nmlblocklst}[
,label={LST:nml_sem}
,caption={Namelist nml\_sem.}
]
&nml_sem
  Use_Improved_SEM_Normalization = T
  Use_Profile_For_Initialization = F
  Use_Random_Seed_From_Restart   = F
  Use_Spherical_Eddy_Volume      = F
  Include_Thermo_Fluctuations    = F
  Box_Num_Eddies_Multiplier      = 1.0
  Use_Restart_Eddies             = T
  Tent_Function_Continuity       = -1
/
\end{nmlblocklst}

%==========================================================================

\begin{nmldescitem}
Use_Improved_SEM_Normalization = T
\end{nmldescitem}
\inp*[nmlSEM::UseImprovedSEMNormalization]{}[Use\_Improved\_SEM\_Normalization][]

\begin{nmldescbody}
This uses an exponential weighted moving average of the eddy concentration as
the normalization factor for the unscaled fluctuating velocity field instead of
the value $\sqrt{\mathscr{V}_{\mathrm{box}} / \left(\mathscr{V}_{\mathrm{eddy}}
N_{\mathrm{eddy}}\right)}$ used by the classic SEM formulation.
%
More details can be found in \linecite{Skillen_2016}.
\end{nmldescbody}

%==========================================================================

\begin{nmldescitem}
Use_Profile_For_Initialization = F
\end{nmldescitem}
\inp*[nmlSEM::UseProfileForInitialization]{}[Use\_Profile\_For\_Initialization][]

\begin{nmldescbody}
Specifies whether the SEM profile will be used to initialize the
\emph{entire volume} of the flow domain based on
the $y$ coordinate value and the profile height variable.
%
This is very experimental and only works for very simple geometries
like a flat plate.

%%Hide this, because few users will need it.
%\NOTE
%This does some more specialized flow initializations if used with the
%THX3 test case.

%\FIXME{
%What are the prerequisites regarding structured or unstructured grids,
% and uniformity in a specific direction?
%}
\end{nmldescbody}

%==========================================================================

\begin{nmldescitem}
Use_Random_Seed_From_Restart = F
\end{nmldescitem}
\inp*[nmlSEM::UseRandomSeedFromRestart]{}[Use\_Random\_Seed\_From\_Restart][]

\begin{nmldescbody}
When using an SEM boundary condition, the SEM data will also be
written to the restart files including the seed used for the random number
generator.

If \inptext{Use\_Random\_Seed\_From\_Restart} is true, the seed for the
random number generator will be read from
the restart file and used to reset the random number generator to exactly
where it left off when the restart file was written, as if the simulation
never stopped.

If it is false, the current date and time will be used to create a new
``random'' seed for the random number generator. This might have a very small
effect on the ``randomness'' of the generated strength and locations of the
synthetic eddies in the SEM BC.
%
However, the total number of time steps needed for a simulation will almost
certainly prevent this from having any kind of measurable effect on the
simulation results.
%
\NOTE Very rarely (but it does happen), the SEM can randomly cause a
simulation to die due to an NaN or other floating point exception.
Setting this input to false is the only way to get past one of these
occasions without modifying the solution with limiting \andor filtering.
\end{nmldescbody}

%==========================================================================

\FINALFIX{\newpage}
\begin{nmldescitem}
Use_Spherical_Eddy_Volume = F
\end{nmldescitem}
\inp*[nmlSEM::UseSphericalEddyVolume]{}[Use\_Spherical\_Eddy\_Volume][]

\begin{nmldescbody}
This input determines if the volume of each synthetic eddy is treated as a
sphere (true) or as a box (false).
%
%It really only has an effect on the number of eddies that will be used
It affects the number of eddies that will be used
for an individual SEM boundary,
%
%[{Number of eddies}] = \inptext{Box\_Num\_Eddies\_Multiplier} $\times$
%                       [{Volume of SEM box}]
%                     / [{Volume of synthetic eddy}].
%
\begin{gather*}
     [\text{Number of eddies}]
  =  [\inptext{Box\_Num\_Eddies\_Multiplier}] \times
     \frac{ [\text{Volume of SEM box}] }
          { [\text{Volume of synthetic eddy}] }
\end{gather*}
%
where the volumes are described in \sref{SEC:SEM:theory:Nomenclature}.

\end{nmldescbody}

%==========================================================================

\begin{nmldescitem}
Include_Thermo_Fluctuations = F
\end{nmldescitem}
\inp*[nmlSEM::IncludeThermoFluctuations]{}[Include\_Thermo\_Fluctuations][]

\begin{nmldescbody}
Specifies whether SEM fluctuations will be applied to the temperature and
density profiles based on the Strong Reynolds Analogy for temperature
fluctuations.
%
If this input is false, no fluctuations are added to the temperature and
density profiles.
%
See \sref{SEC:SEM:theory:ComputeBC} for details.
%
\end{nmldescbody}

%==========================================================================

\begin{nmldescitem}
Box_Num_Eddies_Multiplier = 1.0
\end{nmldescitem}
\inp*[nmlSEM::BoxNumEddiesMultiplier]{}[Box\_Num\_Eddies\_Multiplier][]

\begin{nmldescbody}
This input is used as a multiplier to the total number of synthetic eddies
used for the SEM boundary conditions, \thatis
%
%[{Number of eddies}] = [\inptext{Box\_Num\_Eddies\_Multiplier}] $\times$
%                       [{Volume of SEM box}] /
%                       [{Volume of synthetic eddy}].
%
\begin{gather*}
     [\text{Number of eddies}]
  =  [\inptext{Box\_Num\_Eddies\_Multiplier}] \times
     \frac{ [\text{Volume of SEM box}] }
          { [\text{Volume of synthetic eddy}] }
\end{gather*}
%
where the volumes are described in \sref{SEC:SEM:theory:Nomenclature}.
 
\NOTE 
The number of eddies for an individual SEM box are limited to be between a
minimum of 5,000 and a maximum of 500,000.
%
These limits were arbitrarily picked and hardwired into the SEM code, and
either limit can be modified in the source code if needed.
%
They were implemented to prevent
(1) too few eddies from being present to adequately model the flow, and
(2) too many eddies that add little to the modeling, but increase
    computational cost.
%
\end{nmldescbody}

%==========================================================================

\begin{nmldescitem}
Use_Restart_Eddies = T
\end{nmldescitem}
\inp*[nmlSEM::UseRestartEddies]{}[Use\_Restart\_Eddies][]

\begin{nmldescbody}
This input determines whether or not the SEM information contained in a restart
file is used or discarded after reading in the restart file.
%
It is recommended to use the SEM information in the restart file.
%
However, there are certain situations where starting with a fresh SEM is needed
or helpful which requires not using the SEM information in the restart file.
%
Probably the most common scenario is if the SEM convective velocity or length
scale is changed between simulations.
%
A change in either of these two values will result in a change to the size
of the SEM box, which will subsequently cause an abort error if any of the
restart eddies are no longer contained by the new SEM box.

\NOTE
Not using the SEM information in the restart file will result in a transient
pressure wave at the start of the restart simulation due to the discontinuity in
the SEM boundary condition between simulations.
\end{nmldescbody}

%==========================================================================

\FINALFIX{\newpage}
\begin{nmldescitem}
Tent_Function_Continuity = T
\end{nmldescitem}
\inp*[nmlSEM::TentFunctionContinuity]{}[Tent\_Function\_Continuity][]

\begin{nmldescbody}
This input specifies which tent function to use for the SEM.
%
Recall from \sref{SEC:SEM:theory:SumFluctuations} that the tent function
used by GFR is defined as
%
\begin{equation*}
  f_{\wedge}(a) = 
  \begin{cases}
    g(a) & \text{if } \abs{a} < 1 \\
    0    & \text{otherwise}
  \end{cases}
\end{equation*}
%
This input specifies the function $g(a)$ to use with the tent function, and the
input value generally corresponds to the continuity, $C^n$, of the overall tent
function:\\[-5pt]

\begin{tabular}{ @{} r @{ \qquad $g(a)=$} rl @{\quad : } l @{} }
  -1 & $\sqrt{\tfrac{     3}{    2}}$ & $\mspace{-17mu}\left(1-\abs{a}\right)$ & $C^0$ continuous (default) \\[5pt]
   0 & $\sqrt{\tfrac{    15}{   16}}$ & $\mspace{-17mu}\left(1-a^2\right)$     & $C^0$ continuous  \\[5pt]
   1 & $\sqrt{\tfrac{   315}{  256}}$ & $\mspace{-17mu}\left(1-a^2\right)^2$   & $C^1$ continuous  \\[5pt]
   2 & $\sqrt{\tfrac{  6006}{ 4096}}$ & $\mspace{-17mu}\left(1-a^2\right)^3$   & $C^2$ continuous  \\[5pt]
   3 & $\sqrt{\tfrac{109395}{65536}}$ & $\mspace{-17mu}\left(1-a^2\right)^4$   & $C^3$ continuous
\end{tabular}
\\[5pt]
\NOTE
The non-default tent functions have yet to be thoroughly tested.
\end{nmldescbody}

%==========================================================================


\subsection{nml\_chemistry}
\label{SEC:nml_chemistry}
\inp*[nmlChemistry]{}[nml\_chemistry][]

This \hyperref[SEC:User_Input]{namelist block}
is used to specify control of chemically reacting flows.

\begin{nmlblocklst}[
,label={LST:nml_chemistry}
,caption={Namelist nml\_chemistry.}
]
&nml_chemistry
  Disable_Chemical_Reactions       = F
  Semi_Implicit_Chemistry          = 0
  Ignition_Method                  = 0
  Ignition_Start_Timestep          = 0
  Ignition_Duration                = 1
  Ignition_Prod2Reac_L2_Ratio_Goal = -1.0
  Ignition_Max_Temperature         = 1600.0         ! [K]
  Ignition_Spark_Cap               = 1.0e+12        ! [N/m^2]
  Ignition_Source_Magnitude        = 20.0
  Ignition_Source_Radius           = 0.0            ! [m]
  Ignition_Source_Location         = [0.0,0.0,0.0]  ! [m]
/
\end{nmlblocklst}

%==========================================================================

\begin{nmldescitem}
Disable_Chemical_Reactions = F
\end{nmldescitem}
\inp*[nmlChemistry::DisableChemicalReactions]{}[Disable\_Chemical\_Reactions][]

\begin{nmldescbody}
If running a multi-component gas mixture, setting this to true will disable
all chemical reactions, \thatis all reaction source terms are set to zero,
resulting in a simulation of an inert gas mixture.
 
When reactions are disabled, all non-existent chemical species will be
dropped from the simulation to reduce computational cost. For example, if
running the 7-species ($\mathrm{H_2}$, $\mathrm{O_2}$, $\mathrm{H_2 O}$,
OH, H, O, $\mathrm{N_2}$)\slash 7-reaction mechanism for a hydrogen\slash
air mixture with reactions disabled and the flow initialized to a
mixture of $\mathrm{H_2}$, $\mathrm{O_2}$, and $\mathrm{N_2}$,
then the product species $\mathrm{H_2 O}$, OH, H, and O will be dropped
completely from the simulation.
%
If a simulation is run with reactions disabled, any dropped species
will not be written to the restart files. However, these non-reacting
restart files can still be used to restart a simulation with reactions
enabled and the product species will be added to the simulation with an
initial mass fraction of zero.
\end{nmldescbody}

%==========================================================================

\begin{nmldescitem}
Semi_Implicit_Chemistry = 0
\end{nmldescitem}
\inp*[nmlChemistry::SemiImplicitChemistry]{}[Semi\_Implicit\_Chemistry][]

\begin{nmldescbody}
If running a multi-component gas mixture with reactions, this input controls
the numerical treatment of the reaction source terms:

\begin{tabular}{ @{} r @{ = } p{0.85\linewidth} @{} }
   0 & Standard explicit treatment of the reaction source terms, which is
       likely to require the smallest time step.\\
   1 & Point implicit treatment of the reaction source terms for the species
       transport equations, using an analytical evaluation of the Jacobian
       matrix.\\
   2 & Point implicit treatment of the reaction source terms for the species
       transport equations, using a numerical approximation of the Jacobian
       matrix.\\
  -1 & Use special hard-coded source terms if using the 7-species \slash
       7-reaction hydrogen\slash air reaction mechanism.
       Otherwise, the same as option 0.
\end{tabular}

If the reaction terms require a smaller time step than what would be
required without reactions, then using a point implicit treatment of the
reaction source terms may (or may not) allow for using a time step closer
to the non-reacting case.
\vspace{3pt}
\end{nmldescbody}

%==========================================================================

\FINALFIX{\newpage}
\begin{nmldescitem}
Ignition_Method = 0
\end{nmldescitem}
\inp*[nmlChemistry::IgnitionMethod]{}[Ignition\_Method][]

\begin{nmldescbody}
Specifies the ignition method to use for speeding up the startup of chemical
reactions within the solution.

\begin{tabular}{ @{} r @{ = } p{0.85\linewidth} @{} }
   0 & Ignition process is disabled.\\
   1 & A source term of $\rho c_p \Delta T$ is added to the energy equation each
       Runge-Kutta stage, where $\Delta T = 
       \inp[nmlChemistry::IgnitionSourceMagnitude]{Ignition\_Source\_Magnitude}
       / [\#\ \textsf{RK Stages}]$.\\
   2 & A percentage of the reference total energy, $(\rho e_t)_{\infty}$, is
       added to the energy equation as a source term. The input
       \inptext{Ignition\_Source\_Magnitude} specifies the percentage to use, so
       the source term will be
       $\inptext{Ignition\_Source\_Magnitude}\times(\rho e_t)_{\infty}$.\\
   3 & The temperature value specified with
       \inptext{Ignition\_Source\_Magnitude} is used instead of the actual
       local temperature for the solution point when computing the species
       reaction source terms.
\end{tabular}
\end{nmldescbody}

%==========================================================================

\begin{nmldescitem}
Ignition_Start_Timestep = 0
\end{nmldescitem}
\inp*[nmlChemistry::IgnitionStartTimestep]{}[Ignition\_Start\_Timestep][]

\begin{nmldescbody}
Iteration number to begin using the ignition source terms for chemistry
reactions.
\end{nmldescbody}

%==========================================================================

\begin{nmldescitem}
Ignition_Duration = 1
\end{nmldescitem}
\inp*[nmlChemistry::IgnitionDuration]{}[Ignition\_Duration][]

\begin{nmldescbody}
Specifies the number of time steps for the ignition source term to remain
active once either the temperature at any solution point reaches
\inp[nmlChemistry::IgnitionMaxTemperature]{Ignition\_Max\_Temperature} or the
target goal for the overall ratio of L2(products)\slash L2(reactants),
\inp[nmlChemistry::IgnitionProd2ReacL2RatioGoal]{Ignition\_Prod2Reac\_L2\_Ratio\_Goal}, 
within the global solution is achieved. 
%
The ignition source terms turn off after this number of time steps are
completed.
\end{nmldescbody}

%==========================================================================

\begin{nmldescitem}
Ignition_Prod2Reac_L2_Ratio_Goal = -1.0
\end{nmldescitem}
\inp*[nmlChemistry::IgnitionProd2ReacL2RatioGoal]{}[Ignition\_Prod2Reac\_L2\_Ratio\_Goal][]

\begin{nmldescbody}
Target goal for the ratio of L2(products)\slash L2(reactants) of the simulation
to signify that reactions are occurring. Once this goal is met, the ignition
source terms will continue to be active for the number of time steps specified
by \inp[nmlChemistry::IgnitionDuration]{Ignition\_Duration}, after which, the
ignition source terms will be disabled.
\end{nmldescbody}

%==========================================================================

\begin{nmldescitem}
Ignition_Max_Temperature = 1600.0  ! [K]
\end{nmldescitem}
\inp*[nmlChemistry::IgnitionMaxTemperature]{}[Ignition\_Max\_Temperature][]

\begin{nmldescbody}
Specifies a value for temperature (in Kelvin) that marks when to begin turning
off the ignition source terms. Once this temperature has been reached, the
ignition source terms will continue to be active for the number of time steps
specified by \inp[nmlChemistry::IgnitionDuration]{Ignition\_Duration}, after
which, the ignition source terms will be disabled.
%
Any solution point that reaches this temperature while the ignition source terms
are still active will have its ignition source term set to zero.
%
When using $\inp[nmlChemistry::IgnitionMethod]{Ignition\_Method}=3$, this means
the actual local temperature at the solution point will be used to compute the
reaction source terms instead of the reaction temperature specified by
\inp[nmlChemistry::IgnitionSourceMagnitude]{Ignition\_Source\_Magnitude}.
\end{nmldescbody}

%==========================================================================

\begin{nmldescitem}
Ignition_Spark_Cap = 1.0e+12  ! [N/m^2]
\end{nmldescitem}
\inp*[nmlChemistry::IgnitionSparkCap]{}[Ignition\_Spark\_Cap][]

\begin{nmldescbody}
This specifies a maximum limit that can be used for the energy source term
(in $\mathrm{N}/\mathrm{m}^2$).
%
%\FIXME{
%Are these the correct units for an energy source term? YES.
%\inptext{Ignition\_Source\_Magnitude} uses Kelvin for ignition source term.
%When IgnSrcMag is Kelvin, it is multiplied by $\rho c_p$ to create the source
%term which is then compared to this value.
%}
\end{nmldescbody}

%==========================================================================

\begin{nmldescitem}
Ignition_Source_Magnitude = 20.0
\end{nmldescitem}
\inp*[nmlChemistry::IgnitionSourceMagnitude]{}[Ignition\_Source\_Magnitude][]

\begin{nmldescbody}
Value to use for the ignition source term. The value specified depends on the
ignition method specified by
\inp[nmlChemistry::IgnitionMethod]{Ignition\_Method}.

\begin{InputList}
  \item[$\inptext{Ignition\_Method}=1$]
    This input represents a temperature value (units of \unitTemp{}) that is
    used to compute a source term that is added
    to the energy equation each Runge-Kutta stage. The energy source term is
    evaluated as $\rho c_p \Delta T$ where $\Delta T = 
    \inptext{Ignition\_Source\_Magnitude} / [\#\ \textsf{RK Stages}]$.
  \item[$\inptext{Ignition\_Method}=2$]
    This value represents a percentage value that is used to add a percentage
    of the reference total energy, $(\rho e_t)_{\infty}$, as a source term to
    the energy equation. The energy source term is evaluated as
    $\inptext{Ignition\_Source\_Magnitude}\times(\rho e_t)_{\infty}$.
  \item[$\inptext{Ignition\_Method}=3$]
    This value represents a reaction temperature (units of \unitTemp{}) that
    is used to compute the species reaction source terms instead of using the
    actual local temperature for the solution point.
\end{InputList}
\end{nmldescbody}

%==========================================================================

\begin{nmldescitem}
Ignition_Source_Radius = 0.0  ! [m]
\end{nmldescitem}
\inp*[nmlChemistry::IgnitionSourceRadius]{}[Ignition\_Source\_Radius][]

\begin{nmldescbody}
Radius of influence ($R$) for the ignition source terms,
specified in 
\key[key:GridScaledUnits]{\emph{grid scaled units}} (\thatis meters)
as described in \sref{SEC:Grid_Scaling}.
%
The behavior depends upon the value ($R$) specified:
\\
\begin{tabular}{ @{\qquad} r @{ \quad } p{0.85\linewidth} @{} }
  $<0$ & The ignition source term will immediately decay
         as $\exp (-36 r^2 / R^2 )$ from the point specified by
\inp[nmlChemistry::IgnitionSourceLocation]{Ignition\_Source\_Location}.\\
  $=0$ & The ignition source is disabled.\\
  $>0$ & Every point that is within the ignition source radius ($R$)
         from the \inptext{Ignition\_Source\_Location}
         will get the full value of the energy source term
         based on the input value of
\inp[nmlChemistry::IgnitionSourceMagnitude]{Ignition\_Source\_Magnitude}.
         For all points outside this radius (R),
         the ignition source term will decay as
         $\exp (-36 (r-R)^3 / R^3 )$.
         The higher exponent realizes a slower initial decay.\\
\end{tabular}
\\
\Fref{FIG:Ignition_Source_Radius} compares the decay profiles.

\begin{figure}[h]
  \begin{centering}
    \insertimage[width=3.0in]{Ignition/Ignition_Source_Radius}
    \caption{Ignition source decay.}
    \label{FIG:Ignition_Source_Radius}
  \end{centering}
\end{figure}


\end{nmldescbody}

%==========================================================================

\pagebreak[3]
\begin{nmldescitem}
Ignition_Source_Location = [0.0,0.0,0.0]  ! [m]
\end{nmldescitem}
\inp*[nmlChemistry::IgnitionSourceLocation]{}[Ignition\_Source\_Location][]%

\begin{nmldescbody}
\raggedright % Prevent equation from entering right margin.
The center of the ignition source region, specified in 
\key[key:GridScaledUnits]{\emph{grid scaled units}} (\thatis meters)
as described in \sref{SEC:Grid_Scaling}.
\end{nmldescbody}

%==========================================================================


\subsection{nml\_testcase\_channel}
\label{SEC:nml_testcase_channel}
\inp*[nmlTestcaseChannel]{}[nml\_testcase\_channel][]

This \hyperref[SEC:User_Input]{namelist block} is used to specify inputs for the
rectangular channel flow test case
\\% Prevent next word from entering right margin
($\inp[nmlInitialization::Initialization]{Initialization}=2$) or the
annular-pipe channel flow test case ($\inptext{Initialization}=25$).
%
The inputs that start with \quoterm{Grid\_} are specific to running one of these
cases using the internal grid generator.

\begin{nmlblocklst}[
,label={LST:nml_testcase_channel}
,caption={Namelist nml\_testcase\_channel.}
]
&nml_testcase_channel
  Grid_Annulus_Angle  =   5.0    ! [deg]
  Grid_Ave_Xplus      = -30.0
  Grid_Ave_Yplus      = -15.0
  Grid_Ave_Zplus      = -15.0
  Grid_Yplus_Wall     =   2.0
  Grid_Yplus_Of_One   = 8.42e-6  ! [m]
  Body_Force          =   0.0    ! [N]
  Nondim_Output_Coord = F
/
\end{nmlblocklst}

%==========================================================================

%##%\FINALFIX{\newpage}
\begin{nmldescitem}
Grid_Annulus_Angle = 5.0  ! [deg]
\end{nmldescitem}
\inp*[nmlTestcaseChannel::GridAnnulusAngle]{}[Grid\_Annulus\_Angle][]

\begin{nmldescbody}
%
This input is used when running the annular-pipe channel flow case and
specifies the angle (in degrees) of the circumferential extent of the grid
domain.
%
The input value will be limited to an angle between 0.1 and 90 degrees.
%
The radial distance between the inner and outer walls of the annular pipe
will be set to the channel height used for the rectangular channel flow case.
%
The arc length of the annular pipe \emph{outer wall} will be set to the span used for
the regular channel flow case.
%
From the arc length and circumferential angle, the outer radius of the annulus
will be computed.
%
This means the amount of curvature for the annular pipe will vary depending on
the input \inptext{Grid\_Annulus\_Angle}, with a maximum value of 90 resulting in a
highly curved annular pipe and a minimum value of 0.1 have little-to-no
curvature.
%
\end{nmldescbody}

%==========================================================================

\begin{nmldescitem}
Grid_Ave_Xplus = -30.0
\end{nmldescitem}
\inp*[nmlTestcaseChannel::GridAveXplus]{}[Grid\_Ave\_Xplus][]

\begin{nmldescbody}
When running either of the channel flow problems using the internal grid generator,
this input is used to determine the average grid spacing in the
axial ($x$) direction as follows:
\\
\begin{tabular}{ @{\qquad } r @{ \quad } p{0.75\linewidth} @{} }
  $<0$ & The abs(input value) specifies the average \yp[x] degree of freedom
         spacing. This means the average \yp[x] cell spacing is given by
         $\left|\inptext{Grid\_Ave\_Xplus}\right|\times(\p[+1]{S})$.\\
  $>0$ & The input value specifies the average \yp[x] cell spacing.
\end{tabular}
%
\end{nmldescbody}

%==========================================================================

\begin{nmldescitem}
Grid_Ave_Yplus = -15.0
\end{nmldescitem}
\inp*[nmlTestcaseChannel::GridAveYplus]{}[Grid\_Ave\_Yplus][]

\begin{nmldescbody}
When running either of the channel flow problems using the internal grid generator,
this input is used to determine the average grid spacing in the vertical
($y$) direction specifically at the center-height of the channel.
%
Note, for the annular-pipe channel flow case, this is applied to the radial
spacing.
\\
\begin{tabular}{ @{\qquad } r @{ \quad } p{0.75\linewidth} @{} }
  $<0$ & The abs(input value) specifies the average \yp{} degree of freedom
         spacing at the center of the channel. This means the average \yp{}
         cell spacing at the center of the channel is given by
         $\left|\inptext{Grid\_Ave\_Yplus}\right|\times(\p[+1]{S})$.\\
  $>0$ & The input value specifies the average \yp{} cell spacing at the
         center of the channel.
\end{tabular}

\noindent\NOTE
This input will be used to determine the \yp{} spacing for the cells
at the middle of the channel height, and the input
\inp[nmlTestcaseChannel::GridYplusWall]{Grid\_Yplus\_Wall}
will be used to determine the \yp{} spacing for the cells on channel walls.
%
A tanh spacing will then be used to get the \yp{} cell spacings for
all the cells in between.
\end{nmldescbody}

%==========================================================================

\begin{nmldescitem}
Grid_Ave_Zplus = -15.0
\end{nmldescitem}
\inp*[nmlTestcaseChannel::GridAveZplus]{}[Grid\_Ave\_Zplus][]

\begin{nmldescbody}
When running either of the channel flow problems using the internal grid generator,
this input is used to determine the average grid spacing in the
spanwise ($z$) direction.
%
Note, for the annular-pipe channel flow case, this is applied to the
circumferential spacing.
\\
\begin{tabular}{ @{\qquad } r @{ \quad } p{0.80\linewidth} @{} }
  $<0$ & The abs(input value) specifies the average \yp[z] degree of freedom
         spacing. This means the average \yp[z] cell spacing is given by
         $\left|\inptext{Grid\_Ave\_Zplus}\right|\times(\p[+1]{S})$.\\
  $>0$ & The input value specifies the average \yp[z] cell spacing.
\end{tabular}
\end{nmldescbody}

%==========================================================================

\begin{nmldescitem}
Grid_Yplus_Wall = 2.0
\end{nmldescitem}
\inp*[nmlTestcaseChannel::GridYplusWall]{}[Grid\_Yplus\_Wall][]

\begin{nmldescbody}
When running either of the channel flow problems using the internal grid generator,
this input is used to determine the height of the first cell off the wall.
%
Note, for the annular-pipe channel flow case, this is applied to the radial wall
spacing.
\\
\begin{tabular}{ @{\qquad } r @{ \quad } p{0.75\linewidth} @{} }
  $<0$ & The location of \yp{}=1 is located at the
         $\mathtt{nint}(\left|\inptext{Grid\_Yplus\_Wall}\right|)$-th
         \newline
         degree of freedom off the wall.\\
  $>0$ & This specifies the height (in \yp{}) of the first cell off the
         wall.
\end{tabular}
\end{nmldescbody}

%==========================================================================

\begin{nmldescitem}
Grid_Yplus_Of_One = 8.42e-6  ! [m]
\end{nmldescitem}
\inp*[nmlTestcaseChannel::GridYplusOfOne]{}[Grid\_Yplus\_Of\_One][]

\begin{nmldescbody}
This input specifies the dimensional length (in meters) corresponding to
\yp{}=1 to use with the internal grid generator for either of the channel flow
problems.
%
Note, for the annular-pipe channel flow case, this is applied to the radial wall
spacing.
\end{nmldescbody}

%==========================================================================

\begin{nmldescitem}
Body_Force = 0.0  ! [N]
\end{nmldescitem}
\inp*[nmlTestcaseChannel::BodyForce]{}[Body\_Force][]

\begin{nmldescbody}
Body force (in Newtons) used to drive the channel flow problem.
\end{nmldescbody}

%==========================================================================

\begin{nmldescitem}
Nondim_Output_Coord = F
\end{nmldescitem}
\inp*[nmlTestcaseChannel::NondimOutputCoord]{}[Nondim\_Output\_Coord][]

\begin{nmldescbody}
If running the channel flow problem and this input is set to true,
then all the grid coordinates written to the CGNS solution files
will be nondimensionalized by the channel half height.
\end{nmldescbody}

%==========================================================================

%\begin{nmldescitem}
%Rake_Save_Interval = 100
%\end{nmldescitem}
%\inp*[nmlTestcaseChannel::RakeSaveInterval]{}[Rake\_Save\_Interval][]
%
%\begin{nmldescbody}
%\NOTE
%(not working)
%
%This input allows the option to collect a time history of the flow
%along a rake in order to extract energy spectrums from a channel flow.
%%
%It specifies the number of time steps between saving the
%instantaneous flow profile along a rake to the time history data.
%
%\NOTE
%This can only be used with the channel flow problem using an
%internally generated grid.
%
%\vspace*{0.5em}%
%\FIXME{
%\begin{itemize}
%\item
%Where are the rake points located?
%\item
%Should this be removed for now and added again when it works?
%\end{itemize}
%}
%\end{nmldescbody}

%==========================================================================


\subsection{nml\_testcase\_pulse}
\label{SEC:nml_testcase_pulse}
\inp*[nmlTestcasePulse]{}[nml\_testcase\_pulse][]

This \hyperref[SEC:User_Input]{namelist block}
is used to specify inputs for the density pulse test case.

\begin{nmlblocklst}[
,label={LST:nml_testcase_pulse}
,caption={Namelist nml\_testcase\_pulse.}
]
&nml_testcase_pulse
  Form     = 0
  Form_Dim = 1
  Length   = 1.0            ! [m]
  Center   = [0.5,0.0,0.0]  ! [m] (x,y,z)
  Width    = 0.2            ! [m]
  Gap      = 0.2            ! [m]
  Repeat   = 0
/
\end{nmlblocklst}
\newlength{\negPulseVspace}
\setlength{\negPulseVspace}{-3pt}

%==========================================================================

\begin{nmldescitem}
Form = 0
\end{nmldescitem}
\inp*[nmlTestcasePulse::Form]{}[Form][]

\begin{nmldescbody}
Specifies the form of the density pulse test shape.
\\
\begin{tabular}{ @{\qquad} r @{ = } p{0.85\linewidth} @{} }
  0 & Combo (Gauss,Square,Triangle,Ellipse) \\
  1 & Gauss \\ 
  2 & Square \\
  3 & Triangle \\
  4 & Ellipse
\end{tabular}
\end{nmldescbody}

%==========================================================================

\begin{nmldescitem}
Form_Dim = 1
\end{nmldescitem}
\inp*[nmlTestcasePulse::FormDim]{}[Form\_Dim][]

\begin{nmldescbody}
Dimension of the density test shape:
  1=1d,
  2=2d.
\end{nmldescbody}

%==========================================================================

%##%\FINALFIX{\newpage}
\begin{nmldescitem}
Length = 1.0  ! [m]
\end{nmldescitem}
\inp*[nmlTestcasePulse::Length]{}[Length][]

\begin{nmldescbody}
1D domain length (in meters) for density pulse test internal grid.
\end{nmldescbody}

%==========================================================================

\begin{nmldescitem}
Center = [0.5,0.0,0.0]  ! [m] (x,y,z)
\end{nmldescitem}
\inp*[nmlTestcasePulse::Center]{}[Center][]

\begin{nmldescbody}
$[x,y,z]$ (in meters) at center of first density pulse.
\end{nmldescbody}

%==========================================================================

\begin{nmldescitem}
Width = 0.2  ! [m]
\end{nmldescitem}
\inp*[nmlTestcasePulse::Width]{}[Width][]

\begin{nmldescbody}
Width (in meters) of density pulse.
\end{nmldescbody}

%==========================================================================

\begin{nmldescitem}
Gap = 0.2  ! [m]
\end{nmldescitem}
\inp*[nmlTestcasePulse::Gap]{}[Gap][]

\begin{nmldescbody}
Gap (in meters) between density pulses.
\end{nmldescbody}

%==========================================================================

\begin{nmldescitem}
Repeat = 0
\end{nmldescitem}
\inp*[nmlTestcasePulse::Repeat]{}[Repeat][]

\begin{nmldescbody}
Number of repeat density pulses.
\end{nmldescbody}

%==========================================================================


\subsection{nml\_testcase\_sod}
\label{SEC:nml_testcase_sod}
\inp*[nmlTestcaseSod]{}[nml\_testcase\_sod][]

This \hyperref[SEC:User_Input]{namelist block}
is used to specify inputs for the 1D Sod shock tube problem.

\begin{nmlblocklst}[
,label={LST:nml_testcase_sod}
,caption={Namelist nml\_testcase\_sod.}
]
&nml_testcase_sod
  Right_Vector = [1.0,0.0,0.0]  ! [m] (x,y,z)
/
\end{nmlblocklst}

%==========================================================================

\begin{nmldescitem}
Right_Vector = [1.0,0.0,0.0]  ! [m] (x,y,z)
\end{nmldescitem}
\inp*[nmlTestcaseSod::RightVector]{}[Right\_Vector][]

\begin{nmldescbody}
This defines the right side normal vector of the diaphragm for the
Sod shock tube problem.
%
It is intended for advanced 2D and 3D cases on an external grid
where the flow problem is not aligned with the $+x$-direction.
%
This orientation vector is used to initialize the right side of the
Sod shock tube, which is assumed to be centered at $[x,y,z]=[0,0,0]$.
\end{nmldescbody}

%==========================================================================

\SeeAlso{%
\inp[nmlOutputVariables::ExactDensity]{Exact\_Density},
\inp[nmlOutputVariables::ExactPressure]{Exact\_Pressure}
}

\subsection{nml\_testcase\_tcf}
\label{SEC:nml_testcase_tcf}
\inp*[nmlTestcaseTCF]{}[nml\_testcase\_tcf][]

This \hyperref[SEC:User_Input]{namelist block}
is used to specify inputs for the Taylor-Couette flow test case.
%
This version of the Taylor-Couette flow case is based on that found
in \linecite{Wang_2021}, which uses the initial conditions from the Taylor-Green
vortex problem as inspiration for the initial conditions for this case to speed
up the time it takes for the flow to become fully turbulent.
%
The default values in this namelist determining the domain extents are based on
that paper.
%
The default cell counts here are slightly coarser than their work which used
64, 20, and 80 grid cells in the respective $x$, $r$, and $\theta$ directions.

\begin{nmlblocklst}[
,label={LST:nml_testcase_tcf}
,caption={Namelist nml\_testcase\_tcf.}
]
&nml_testcase_tcf
  Inner_Radius       = 0.5     ! [m]
  Inner_Wall_Spacing = 0.002   ! [m]
  Outer_Radius       = 1.0     ! [m]
  Outer_Wall_Spacing = 0.0025  ! [m]
  Aspect_Ratio       = 6.2832  ! 2*pi
  Axial_Cells        = 24
  Radial_Cells       = 20
  Azimuthal_Cells    = 48
  Theta_Shift        = 0.0     ! [deg]
  Periodic_Sectors   = 1
/
\end{nmlblocklst}

%==========================================================================

\begin{nmldescitem}
Inner_Radius = 0.5  ! [m]
\end{nmldescitem}
\inp*[nmlTestcaseTCF::InnerRadius]{}[Inner\_Radius][]

\begin{nmldescbody}
This gives the radius of the inner wall for the Taylor-Couette flow problem.
\end{nmldescbody}

%==========================================================================

\begin{nmldescitem}
Inner_Wall_Spacing = 0.002  ! [m]
\end{nmldescitem}
\inp*[nmlTestcaseTCF::InnerWallSpacing]{}[Inner\_Wall\_Spacing][]

\begin{nmldescbody}
This gives the wall spacing at the inner wall for the Taylor-Couette flow
problem.
\end{nmldescbody}

%==========================================================================

\begin{nmldescitem}
Outer_Radius = 1.0  ! [m]
\end{nmldescitem}
\inp*[nmlTestcaseTCF::OuterRadius]{}[Outer\_Radius][]

\begin{nmldescbody}
This gives the radius of the outer wall for the Taylor-Couette flow problem.
\end{nmldescbody}

%==========================================================================

\begin{nmldescitem}
Outer_Wall_Spacing = 0.0025  ! [m]
\end{nmldescitem}
\inp*[nmlTestcaseTCF::OuterWallSpacing]{}[Outer\_Wall\_Spacing][]

\begin{nmldescbody}
This gives the wall spacing at the outer wall for the Taylor-Couette flow
problem.
\end{nmldescbody}

%==========================================================================

\begin{nmldescitem}
Aspect_Ratio = 6.2832  ! 2*pi
\end{nmldescitem}
\inp*[nmlTestcaseTCF::AspectRatio]{}[Aspect\_Ratio][]

\begin{nmldescbody}
This is used to specify the axial length, $l_x$, of the domain for the
Taylor-Couette flow problem. The length is determined by the equation:

$
  l_x = \inptext{Aspect\_Ratio} \times
        \left| \inptext{Outer\_Radius}
             - \inptext{Inner\_Radius}
        \right|
$
\end{nmldescbody}

%==========================================================================

\begin{nmldescitem}
Axial_Cells = 24
\end{nmldescitem}
\inp*[nmlTestcaseTCF::AxialCells]{}[Axial\_Cells][]

\begin{nmldescbody}
This specifies the number of axial cells to use for the Taylor-Couette flow
problem.
\end{nmldescbody}

%==========================================================================

\FINALFIX{\newpage}
\begin{nmldescitem}
Radial_Cells = 20
\end{nmldescitem}
\inp*[nmlTestcaseTCF::RadialCells]{}[Radial\_Cells][]

\begin{nmldescbody}
This specifies the number of radial cells to use for the Taylor-Couette flow
problem.
\end{nmldescbody}

%==========================================================================

\begin{nmldescitem}
Azimuthal_Cells = 48
\end{nmldescitem}
\inp*[nmlTestcaseTCF::AzimuthalCells]{}[Azimuthal\_Cells][]

\begin{nmldescbody}
This specifies the number of azimuthal cells to use for the Taylor-Couette
flow problem.
\end{nmldescbody}

%==========================================================================

\begin{nmldescitem}
Theta_Shift = 0.0  ! [deg]
\end{nmldescitem}
\inp*[nmlTestcaseTCF::ThetaShift]{}[Theta\_Shift][]

\begin{nmldescbody}
This specifies an angle (in degrees) to shift the grid in theta for the
Taylor-Couette flow problem. The purpose of this input is to rotate the grid
so that a cell boundary is not located at the angle 0\slash 360 degrees.
\end{nmldescbody}

%==========================================================================

\begin{nmldescitem}
Periodic_Sectors = 1
\end{nmldescitem}
\inp*[nmlTestcaseTCF::PeriodicSectors]{}[Periodic\_Sectors][]

\begin{nmldescbody}
This input can be used to run only a sector of the Taylor-Couette flow
problem. This value specifies how many pie slices to divide the grid into,
with the final grid domain consisting of only one pie slice.
\end{nmldescbody}

%==========================================================================


\subsection{nml\_testcase\_vortex}
\label{SEC:nml_testcase_vortex}
\inp*[nmlTestcaseVortex]{}[nml\_testcase\_vortex][]

This \hyperref[SEC:User_Input]{namelist block}
is used to specify inputs for the vortex test cases.
%
These cases include the different versions of
  the isentropic Euler vortex problem
    ($\inp[nmlInitialization::Initialization]{Initialization}=[\pm5,\pm55,\pm6]$),
  the density profile advection case
    ($\inptext{Initialization}=7$),
  the entropy profile advection case
    ($\inptext{Initialization}=8$),
and
  the Taylor-Green vortex problem
    ($\inp{Initialization}=9$).

\begin{nmlblocklst}[
,label={LST:nml_testcase_vortex}
,caption={Namelist nml\_testcase\_vortex.}
]
&nml_testcase_vortex
  Vortex_Bigr          = 1.0   ! [m]
  Vortex_Biggam        = 5.0   ! [unitless]
  Grid_Lref            = 10.0  ! [m]
  Grid_Random_Perturb  = F
  Grid_Perturb_Percent = 30
  Grid_Hex_To_Pyramids = F
  Grid_Is_3d           = F
/
\end{nmlblocklst}

%==========================================================================

\begin{nmldescitem}
Vortex_Bigr = 1.0  ! [m]
\end{nmldescitem}
\inp*[nmlTestcaseVortex::VortexBigr]{}[Vortex\_Bigr][]

\begin{nmldescbody}
Constant used to specify the size or characteristic length (in meters)
of the isentropic Euler vortex.
\end{nmldescbody}

%==========================================================================

\begin{nmldescitem}
Vortex_Biggam = 5.0  ! [unitless]
\end{nmldescitem}
\inp*[nmlTestcaseVortex::VortexBiggam]{}[Vortex\_Biggam][]

\begin{nmldescbody}
Constant used to specify the strength of the isentropic Euler vortex.
\end{nmldescbody}

%==========================================================================

\begin{nmldescitem}
Grid_Lref = 10.0  ! [m]
\end{nmldescitem}
\inp*[nmlTestcaseVortex::GridLref]{}[Grid\_Lref][]

\begin{nmldescbody}
Reference length (in meters) for the grid domain,
\thatis the grid is a square with bounds of
$\left[ -\inptext{Grid\_Lref},
        +\inptext{Grid\_Lref}
 \right]$
in both the $x$- and $y$-directions.
\end{nmldescbody}

%==========================================================================

\begin{nmldescitem}
Grid_Random_Perturb = F
\end{nmldescitem}
\inp*[nmlTestcaseVortex::GridRandomPerturb]{}[Grid\_Random\_Perturb][]

\begin{nmldescbody}
Logical flag for applying a random perturbation to all the interior grid
nodes for the isentropic Euler vortex and Taylor-Green vortex problems.
\end{nmldescbody}

%==========================================================================

%##%\FINALFIX{\newpage}
\begin{nmldescitem}
Grid_Perturb_Percent = 30
\end{nmldescitem}
\inp*[nmlTestcaseVortex::GridPerturbPercent]{}[Grid\_Perturb\_Percent][]

\begin{nmldescbody}
\raggedright % prevent right margin overflow of input variable
Percent of the uniform grid cell length to perturb the interior grid nodes if
\inp[nmlTestcaseVortex::GridRandomPerturb]{Grid\_Random\_Perturb}
is true.
\end{nmldescbody}

%==========================================================================

\begin{nmldescitem}
Grid_Hex_To_Pyramids = F
\end{nmldescitem}
\inp*[nmlTestcaseVortex::GridHexToPyramids]{}[Grid\_Hex\_To\_Pyramids][]

\begin{nmldescbody}
When using the internal grid generator with one of the 3D vortex problems,
the default is to use hex cells. Setting this to true will create a grid
using only pyramid cells, by first creating the hex-only grid, and then
subdividing each hex cell into 6 pyramids by adding a new node point in the
center of the cell.
\end{nmldescbody}

%==========================================================================

\FINALFIX{\newpage}
\begin{nmldescitem}
Grid_Is_3d = F
\end{nmldescitem}
\inp*[nmlTestcaseVortex::GridIs3d]{}[Grid\_Is\_3d][]

\begin{nmldescbody}
When running one of the traditionally 2D vortex cases (\forexample
isentropic vortex), setting this to true will run a 3D version of the case
by extruding the domain and initial solution in $z$ to create a cube domain.
\end{nmldescbody}

%==========================================================================


\subsection{nml\_interpolation}
\label{SEC:nml_interpolation}
\inp*[nmlInterpolation]{}[nml\_interpolation][]

This \hyperref[SEC:User_Input]{namelist block}
is for use with the startup interpolation utility,
which will interpolate a solution from 1) a set of Plot3d grid and solution
files or 2) any restart file and respective grid file onto a new
(and possibly unstructured) grid file.

\begin{nmlblocklst}[
,label={LST:nml_interpolation}
,caption={Namelist nml\_interpolation.}
]
&nml_interpolation
  Interpto_Grid_Scaling_Factor      =   1.0
  Interpto_Grid_Format              = -99   ! -99 = Unknown
  Interpto_Grid_File                =  ""
  Interpto_Plot3d_Cutfile           =  ""
  Interpto_Plot3d_BCSfile           =  ""
  Interpto_Plot3d_Stride_I          =  -1
  Interpto_Plot3d_Stride_J          =  -1
  Interpto_Plot3d_Stride_K          =  -1
  Interpto_Plot3d_Stride_All        =  -1
  Interpto_Plot3d_Agglomerate_Order =  -1
  Interpto_Solution_Order           =   3
/
\end{nmlblocklst}
\inp*[nmlInterpolation::InterptoGridScalingFactor]{}[Interpto\_Grid\_Scaling\_Factor][]
\inp*[nmlInterpolation::InterptoGridFormat]{}[Interpto\_Grid\_Format][]
\inp*[nmlInterpolation::InterptoGridFile]{}[Interpto\_Grid\_File][]
\inp*[nmlInterpolation::InterptoPlot3dCutfile]{}[Interpto\_Plot3d\_Cutfile][]
\inp*[nmlInterpolation::InterptoPlot3dBCSfile]{}[Interpto\_Plot3d\_BCSfile][]
\inp*[nmlInterpolation::InterptoPlot3dStrideAll]{}[Interpto\_Plot3d\_Stride\_All][]
\inp*[nmlInterpolation::InterptoPlot3dStrideI]{}[Interpto\_Plot3d\_Stride\_I][]
\inp*[nmlInterpolation::InterptoPlot3dStrideJ]{}[Interpto\_Plot3d\_Stride\_J][]
\inp*[nmlInterpolation::InterptoPlot3dStrideK]{}[Interpto\_Plot3d\_Stride\_K][]
\inp*[nmlInterpolation::InterptoPlot3dAgglomerateOrder]{}[Interpto\_Plot3d\_Agglomerate\_Order][]
\inp*[nmlInterpolation::InterptoSolutionOrder]{}[Interpto\_Solution\_Order][]

Each input has the same functionality as the corresponding input
(without the \inptext{Interpto\_} prefix) from
the \inptext{nml\_grid} namelist block of \sref{SEC:nml_grid} or
the \inptext{nml\_numerics} namelist block of \sref{SEC:nml_numerics}.
%
However, those inputs refer to the grid that
the old solution is being interpolated \emph{from}, whereas
these inputs refer to the grid that is being interpolated \emph{to}.

\NOTE
The user must provide a valid \inptext{Interpto\_Grid\_Format},
since the default value of -99 will cause the code to terminate.

%==========================================================================


%##%\FINALFIX{\newpage}
\subsection{nml\_utility}
\label{SEC:nml_utility}
\inp*[nmlUtility]{}[nml\_utility][]

This \hyperref[SEC:User_Input]{namelist block}
is used to switch from solver mode to utility mode.
%
\NOTE
The inputs in this block perform post-processing tasks
without running the flow solver.

\begin{nmlblocklst}[
,label={LST:nml_utility}
,caption={Namelist nml\_utility.}
]
&nml_utility
  Convert_Restart_To_CGNS            = F
  Delete_Converted_Restart_Files     = F
  Restart_List_File                  = ""
  Spatial_Ave_Restart_Channel        = F
  Spatial_Ave_Restart_X_And_Z_Planes = F
  Spatial_Ave_Restart_Z_Planes       = F
/
\end{nmlblocklst}

%==========================================================================

\begin{nmldescitem}
Convert_Restart_To_CGNS = F
\end{nmldescitem}
\inp*[nmlUtility::ConvertRestartToCGNS]{}[Convert\_Restart\_To\_CGNS][]

\begin{nmldescbody}
Flag used to read in a restart file and immediately write a CGNS solution
file, bypassing the solver entirely.
%
This first reads the instantaneous restart file from
\inp[nmlInitialization::RestartFile]{Restart\_File},
computes the requested output variables from
\inptext{nml\_output\_variables} in \sref{SEC:nml_output_variables},
and writes them to a CGNS solution file.
%
Next it reads the time-averaged restart file from
\inp[nmlTimeAve::TimeAveFile]{Time\_Ave\_File}, 
and writes the accumulated data to \file{CGNS\_files/solution.\#.ta\_data.cgns}
and the derived statistics to \file{CGNS\_files/solution.\#.time\_ave.cgns}.
%
Lastly, it reads the list of instantaneous restart files from
\inp[nmlUtility::RestartListFile]{Restart\_List\_File}
and processes each one to compute the requested output variables from
\inptext{nml\_output\_variables} in \sref{SEC:nml_output_variables},
and writes them to a CGNS solution file.
%
Finally, it will write the file \file{CGNS\_files/solution.time.cgns} to
link all the instantaneous solution files together into a single time-history
solution file.

The option 
\inp[nmlUtility::DeleteConvertedRestartFiles]{Delete\_Converted\_Restart\_Files}
controls whether the restart files are deleted after the CGNS output
is complete.
\end{nmldescbody}

%==========================================================================

\begin{nmldescitem}
Delete_Converted_Restart_Files = F
\end{nmldescitem}
\inp*[nmlUtility::DeleteConvertedRestartFiles]{}[Delete\_Converted\_Restart\_Files][]

\begin{nmldescbody}
Specifies whether to delete each restart file from
\inp[nmlUtility::RestartListFile]{Restart\_List\_File}
after it has been processed by
\inp[nmlUtility::ConvertRestartToCGNS]{Convert\_Restart\_To\_CGNS}.
%
This can be used to minimized disk space usage.
%
The
\inp[nmlInitialization::RestartFile]{Restart\_File} and
\inp[nmlTimeAve::TimeAveFile]{Time\_Ave\_File}
will \emph{not} be deleted after conversion to CGNS.

\noindent\myTextColor[red]{\textbf{WARNING:} } This will permanently delete
these files from the file system.
\end{nmldescbody}


%==========================================================================

\begin{nmldescitem}
Restart_List_File = ""
\end{nmldescitem}
\inp*[nmlUtility::RestartListFile]{}[Restart\_List\_File][]

\begin{nmldescbody}
File path to a simple ASCII file that contains a list of file paths to
existing instantaneous restart files for use with
\inp[nmlUtility::ConvertRestartToCGNS]{Convert\_Restart\_To\_CGNS}.
%
Each line of the list file should contain the name of a restart file.
%
If the text on a line does not give a valid path to an existing file, it will
ignore the line and continue with the next line.
%
See also
\inp[nmlUtility::DeleteConvertedRestartFiles]{Delete\_Converted\_Restart\_Files}.
\end{nmldescbody}

%==========================================================================

%##%\FINALFIX{\newpage}
\begin{nmldescitem}
Spatial_Ave_Restart_Channel = F
\end{nmldescitem}
\inp*[nmlUtility::SpatialAveRestartChannel]{}[Spatial\_Ave\_Restart\_Channel][]

\begin{nmldescbody}
This is essentially a specialized version of
\inp[nmlUtility::SpatialAveRestartXAndZPlanes]{Spatial\_Ave\_Restart\_X\_And\_Z\_Planes}
for the channel flow problem
where the 1D $y$ profile is normalized by the channel half-height.
%
It is used to post-process the time-averaged restart file
(\inp[nmlTimeAve::TimeAveFile]{Time\_Ave\_File}).
\end{nmldescbody}

%==========================================================================

\begin{nmldescitem}
Spatial_Ave_Restart_X_And_Z_Planes = F
\end{nmldescitem}
\inp*[nmlUtility::SpatialAveRestartXAndZPlanes]{}[Spatial\_Ave\_Restart\_X\_And\_Z\_Planes][]

\begin{nmldescbody}
Utility to read the time-averaged restart file
(\inp[nmlTimeAve::TimeAveFile]{Time\_Ave\_File})
and spatially average in both the $x$ and $z$ coordinates.
%
This option first produces the $z$-averaged output from
\\%  Prevent next word from entering right margin.
\inp[nmlUtility::SpatialAveRestartZPlanes]{Spatial\_Ave\_Restart\_Z\_Planes},
then additionally averages in the $x$-direction
and outputs a Tecplot ASCII data file containing
a 1D line with the $y$ grid coordinates and the
$x$-$z$-averaged flow variables
(\forexample a hex mesh formed from a $y$ grid line by extruding
in the $x$ direction and then the $z$ direction).

\NOTE
This only works for simple cases like channel flows or flat plates
where the $y$-spacing is the same for every $x$-$z$ coordinate pair.

\end{nmldescbody}

%==========================================================================

\begin{nmldescitem}
Spatial_Ave_Restart_Z_Planes = F
\end{nmldescitem}
\inp*[nmlUtility::SpatialAveRestartZPlanes]{}[Spatial\_Ave\_Restart\_Z\_Planes][]

\begin{nmldescbody}
Utility to read the time-averaged restart file
(\inp[nmlTimeAve::TimeAveFile]{Time\_Ave\_File})
and spatially average in the $z$ coordinates.
%
The output will be a Tecplot ASCII data file containing
a 2D plane with the $x$-$y$ grid coordinates and the
$z$-averaged flow variables.

\NOTE
This only works for simple cases like channel flows or flat plates
where the 2D $x$-$y$ grids are identical for all $z$-planes
(\forexample a structured or unstructured 2D planar $x$-$y$ grid
that is extruded in the $z$ direction).

\end{nmldescbody}

%==========================================================================


\subsection{nml\_debug}
\label{SEC:nml_debug}
\inp*[nmlDebug]{}[nml\_debug][]

This \hyperref[SEC:User_Input]{namelist block}
is used by developers for debugging.

\begin{nmlblocklst}[
,label={LST:nml_debug}
,caption={Namelist nml\_debug.}
]
&nml_debug
  Output_Limiter_Statistics          = F
  Output_Filter_Diagonals            = F
  Output_Plot3d_To_Tecplot           = F
  Output_Bface_Array                 = F
  Dump_Memory_Usage                  = F
  Check_Flux_Point_Coordinates       = F
  Force_Nd_Subroutines               = F
  Newton_Convergence_Goal            = 1.0e-12
  Chemistry_Sub_Iterations           = -1
  Traceback_Unconverged_Newton_Temps = F
  Ignore_Thermo_Newton_Warnings      = F
/
\end{nmlblocklst}

%==========================================================================

\begin{nmldescitem}
Output_Limiter_Statistics = F
\end{nmldescitem}
\inp*[nmlDebug::OutputLimiterStatistics]{}[Output\_Limiter\_Statistics][]

\begin{nmldescbody}
This outputs the number of limited\slash projected cells each time step when
\inp[nmlLimiter::LocalLimiterOption]{Local\_Limiter\_Option} is set to use
projection.
\end{nmldescbody}

%==========================================================================

\begin{nmldescitem}
Output_Filter_Diagonals = F
\end{nmldescitem}
\inp*[nmlDebug::OutputFilterDiagonals]{}[Output\_Filter\_Diagonals][]

\begin{nmldescbody}
Specifies whether to output text files containing the filter diagonals and
filter matrices for all geometry types after they are generated in the
preprocessor.
%
This is used to inspect values for possible errors in the filter
definition.
\end{nmldescbody}

%==========================================================================

\begin{nmldescitem}
Output_Plot3d_To_Tecplot = F
\end{nmldescitem}
\inp*[nmlDebug::OutputPlot3dToTecplot]{}[Output\_Plot3d\_To\_Tecplot][]

\begin{nmldescbody}
Specifies whether to output a Tecplot file of the input Plot3D grid file
after reading it.
%
This can be used to inspect whether the grid file was read correctly.
\end{nmldescbody}

%==========================================================================

\begin{nmldescitem}
Output_Bface_Array = F
\end{nmldescitem}
\inp*[nmlDebug::OutputBfaceArray]{}[Output\_Bface\_Array][]

\begin{nmldescbody}
Specifies whether to have each MPI process output the contents of its
own local bface array to the text file ``bface.[MPI process \#].dat''.
%
The bface array contains the boundary conditions and connectivity
information for each boundary face.
\end{nmldescbody}

%==========================================================================

\begin{nmldescitem}
Dump_Memory_Usage = F
\end{nmldescitem}
\inp*[nmlDebug::DumpMemoryUsage]{}[Dump\_Memory\_Usage][]

\begin{nmldescbody}
Specifies whether to profile and output the memory usage of GFR.
\end{nmldescbody}

%==========================================================================

\begin{nmldescitem}
Check_Flux_Point_Coordinates = F
\end{nmldescitem}
\inp*[nmlDebug::CheckFluxPointCoordinates]{}[Check\_Flux\_Point\_Coordinates][]

\begin{nmldescbody}
Flag used to output files to check that co-located face quantities
are consistent in terms of interpolation between two cells sharing a face.
\end{nmldescbody}

%==========================================================================

\FINALFIX{\newpage}
\begin{nmldescitem}
Force_Nd_Subroutines = F
\end{nmldescitem}
\inp*[nmlDebug::ForceNdSubroutines]{}[Force\_Nd\_Subroutines][]

\begin{nmldescbody}
GFR normally uses optimized 1D, 2D, and 3D versions of various
subroutines.
%
This input forces the use of the generalized ``nD'' versions of the
various subroutines that loop over the number of equations instead
of hardcoding in each equation.
\end{nmldescbody}

%==========================================================================

\begin{nmldescitem}
Newton_Convergence_Goal = 1.0e-12
\end{nmldescitem}
\inp*[nmlDebug::NewtonConvergenceGoal]{}[Newton\_Convergence\_Goal][]

\begin{nmldescbody}
This specifies the convergence goal for all of the Newton iterations on
temperature.
\end{nmldescbody}

%==========================================================================

\begin{nmldescitem}
Chemistry_Sub_Iterations = -1
\end{nmldescitem}
\inp*[nmlDebug::ChemistrySubIterations]{}[Chemistry\_Sub\_Iterations][]

\begin{nmldescbody}
Number of chemistry sub-iterations to use for the KIVA scheme source terms.
%
\NOTE
The number of sub-iterations is set to 10 if
$\inptext{Chemistry\_Sub\_Iterations}<1$.
\end{nmldescbody}

%==========================================================================

\begin{nmldescitem}
Traceback_Unconverged_Newton_Temps = F
\end{nmldescitem}
\inp*[nmlDebug::TracebackUnconvergedNewtonTemps]{}[Traceback\_Unconverged\_Newton\_Temps][]

\begin{nmldescbody}
If running a multi-component gas mixture, setting this input to true
will write a traceback to stdout whenever the Newton solve for
temperature fails to meet the convergence goal.

\NOTE
This only works if using the Intel Fortran compiler and
the namelist input variable
\inp[nmlDebug::IgnoreThermoNewtonWarnings]{Ignore\_Thermo\_Newton\_Warnings}
is set to false.
\end{nmldescbody}

%==========================================================================

\begin{nmldescitem}
Ignore_Thermo_Newton_Warnings = F
\end{nmldescitem}
\inp*[nmlDebug::IgnoreThermoNewtonWarnings]{}[Ignore\_Thermo\_Newton\_Warnings][]

\begin{nmldescbody}
If running a multi-component gas mixture, setting this to true will prevent
any warnings from being written to stdout whenever the Newton solve for
temperature fails to converge.
\end{nmldescbody}

%==========================================================================



% These sections need to be reviewd and removed!
%\subsection{nml\_mms}
\label{SEC:nml_mms}
\inp*[nmlMMS]{}[nml\_mms][]

%==========================================================================

%\AddEntry[nmlMMS::Option]{nml\_mms}{Option}

%==========================================================================

%\AddEntry[nmlMMS::Initialization]{nml\_mms}{Initialization}

%==========================================================================

%\AddEntry[nmlMMS::OutputSolution]{nml\_mms}{Output\_Solution}

%==========================================================================


%\subsection{nml\_removed}
\label{SEC:nml_removed}
\inp*[nmlRemoved]{}[nml\_removed][]

This \hyperref[SEC:User_Input]{namelist block}
is used to track which options have been removed from the previous
input file format.

%==========================================================================

\begin{nmldescitem}
Time_Average_Restart_Files = F
\end{nmldescitem}
\inp*[nmlRemoved::TimeAverageRestartFiles]{}[Time\_Average\_Restart\_Files][]

\begin{nmldescbody}
Logical flag to run a utility that creates a time-averaged solution from a
list of existing restart files given in the file specified by the
\inptext{Restart\_List\_File} input parameter. This utility will finish by
writing out a time-averaged restart file and a time-averaged CGNS solution
file. This utility does not execute anything within the solver.
\end{nmldescbody}

%==========================================================================

%\AddEntry[nmlRemoved::LocInitPts]{nml\_removed}{Loc\_Init\_Pts}

%==========================================================================



\subsection{Changes from Prior Version}
\label{SEC:nml_Changes}

Many new features have been added since the prior version, including:
\begin{itemize}
\item
Mixed unstructured elements are now fully functional, \sref{SEC:Element_Types}.
\item
Rotating reference frame, \sref{SEC:Ref_Frame}.
\item
A number of new time stepping methods, \sref{SEC:Time}.
\item
  Unsteady turbulent inflow via SEM, \sref{SEC:SEM}.
\item
The Multidimensional Limiting Process (MLP).
%The Multidimensional Limiting Process (MLP), \sref{SEC:Limiting}.
\item
The \inp[nmlLimiter::Thermofix/Method]{Thermofix\%Method}
to prevent negative density, temperature, or pressure.
\item
Test cases for:
\begin{itemize}
\item
3D Taylor-Couette flow
% (\sref{SEC:Testcases:TaylorCouetteFlow})
\item
3D channel flow and annular pipe channel flow
% (\sref{SEC:Testcases:ChannelFlow})
\item
1D density pulse advection (\sref{SEC:Testcases:DensityPulse})
\item
1D Sod shock tube          (\sref{SEC:Testcases:Sod})
\item
1D Shu-Osher problem       (\sref{SEC:Testcases:ShuOsher})
\end{itemize}
\end{itemize}


One of the most noticeable changes is the renaming and reorganization of the
user input.
%
The code will attempt to convert an older input file to the current format.
%
However, should you need to manually convert an older input file,
the table below provides a mapping between the old and new inputs.

\noindent\NOTE An \inptext{*} means an equivalent input does not exist in the old namelist.


% This table maps version 1 to version 2 namelist format.

% This rule is used to manually determine column widths.
% \colsep    =  6pt  Pads on L+R.
% \parindent = 15pt
{% beg scope for rule
\newcommand{\myrule}{%
  \rule       {   6pt}{1.0pt}% PadL
  \rule       {   5pt}{1.0ex}%
  \rule[1.0ex]{   6pt}{1.0pt}% PadR
  \rule       {   6pt}{1.0pt}% PadL
  \rule       {   5pt}{1.0ex}%
  \rule[1.0ex]{   6pt}{1.0pt}% PadR
  \rule       {   6pt}{1.0pt}% PadL
  \rule       {2.80in}{1.0ex}%
  \rule[1.0ex]{   6pt}{1.0pt}% PadR
  \rule       {   6pt}{1.0pt}% PadL
  \rule       {   5pt}{1.0ex}%
  \rule[1.0ex]{   6pt}{1.0pt}% PadR
  \rule       {   6pt}{1.0pt}% PadL
  \rule       {2.30in}{1.0ex}%
  \rule[1.0ex]{   6pt}{1.0pt}% PadR
}
%\noindent\myrule
}% end scope for rule

%\begin{table}[ht]
  \begin{center}
    {\normalsize % begin font size
    % \tabcolsep =  6pt ~= 0.0833in Pads on L+R.
    % \parindent = 15pt ~= 0.2083in
    \newlength{\lenS}
    \newlength{\lenA}
    \newlength{\lenB}
    \newlength{\lenC}
    \newlength{\lenD}
    \newlength{\lenE}
    \newlength{\lenF}
    \newlength{\lenG}
    \newlength{\lenH}
    \setlength{\lenS}{6pt}
    \setlength{\lenA}{\parindent}
    \setlength{\lenB}{\parindent}
    \setlength{\lenC}{2.80in}
    \setlength{\lenD}{\parindent}
    \setlength{\lenE}{\parindent}
    \setlength{\lenG}{\lenC+\lenD}
    \setlength{\lenH}{\lenG+\lenE}
    \setlength{\lenF}{\textwidth-\lenA-\lenB-\lenC-\lenD-\lenE-5\lenS}
   %\setlength{\lenF}{\textwidth-\lenA-\lenB-\lenC-\lenD-\lenE}
    \begin{longtable}{ p{\lenA} p{\lenB} p{\lenC} p{\lenD} p{\lenE} p{\lenF} }
%   \begin{longtable}{|p{\lenA}|p{\lenB}|p{\lenC}|p{\lenD}|p{\lenE}|p{\lenF}|}
      \hline
       \multicolumn{3}{l}{Old Namelist}
    &  \multicolumn{2}{l}{New Namelist}
    \\
       &  \multicolumn{2}{l}{Old Variable}
       && \multicolumn{2}{l}{New Variable}
       \\
       && \multicolumn{2}{l}{Description}
       &  
       \\
      \hline
      \endhead
%==========================================================================
       \multicolumn{3}{l}{\&input}
    &  \multicolumn{2}{l}{\&\texttt{nml\_version}}
    \\* %%% Note: \\* prevents pagebreak.
       &  \multicolumn{2}{l}{\inptext{*}}
       && \multicolumn{2}{l}{\inp[nmlVersion::NamelistVersion]{Namelist\_Version}}
       \\
%==========================================================================
       \multicolumn{3}{l}{\&input}
    &  \multicolumn{2}{l}{\&\texttt{nml\_grid}}
    \\* %%% Note: \\* prevents pagebreak.
       &  \multicolumn{2}{l}{\inptext{Grid\_Scaling\_Factor}}
       && \multicolumn{2}{l}{\inp[nmlGrid::GridScalingFactor]{Same}[Grid\_Scaling\_Factor]}
       \\
       &  \multicolumn{2}{l}{\inptext{Grid\_Format}}
       && \multicolumn{2}{l}{\inp[nmlGrid::GridFormat]{Same}[Grid\_Format]}
       \\
       &  \multicolumn{2}{l}{\inptext{Gridfile}}
       && \multicolumn{2}{l}{\inp[nmlGrid::GridFile]{Grid\_File}}
       \\
       &  \multicolumn{2}{l}{\inptext{Plot3d\_Cutfile}}
       && \multicolumn{2}{l}{\inp[nmlGrid::Plot3dCutfile]{Same}[Plot3d\_Cutfile]}
       \\
       &  \multicolumn{2}{l}{\inptext{Plot3d\_BCSfile}}
       && \multicolumn{2}{l}{\inp[nmlGrid::Plot3dBCSfile]{Same}[Plot3d\_BCSfile]}
       \\
       &  \multicolumn{2}{l}{\inptext{Plot3d\_I\_Stride}}
       && \multicolumn{2}{l}{\inp[nmlGrid::Plot3dStrideI]{Plot3d\_Stride\_I}}
       \\
       &  \multicolumn{2}{l}{\inptext{Plot3d\_J\_Stride}}
       && \multicolumn{2}{l}{\inp[nmlGrid::Plot3dStrideJ]{Plot3d\_Stride\_J}}
       \\
       &  \multicolumn{2}{l}{\inptext{Plot3d\_K\_Stride}}
       && \multicolumn{2}{l}{\inp[nmlGrid::Plot3dStrideK]{Plot3d\_Stride\_K}}
       \\
       &  \multicolumn{2}{l}{\inptext{Plot3d\_All\_Strides}}
       && \multicolumn{2}{l}{\inp[nmlGrid::Plot3dStrideAll]{Plot3d\_Stride\_All}}
       \\
       &  \multicolumn{2}{l}{\inptext{Plot3d\_Agglomerate\_Order}}
       && \multicolumn{2}{l}{\inp[nmlGrid::Plot3dAgglomerateOrder]{Same}[Plot3d\_Agglomerate\_Order]}
       \\
       &  \multicolumn{2}{l}{\inptext{Correct\_CGNS\_Grid\_Points}}
       && \multicolumn{2}{l}{\inp[nmlGrid::CorrectCGNSGridPoints]{Same}[Correct\_CGNS\_Grid\_Points]}
       \\
       &  \multicolumn{2}{l}{\inptext{Flame\_Speed\_Test\_DoF}}
       && \multicolumn{2}{l}{\inp[nmlGrid::InternalGridDoF]{Internal\_Grid\_DoF}}
       \\
       &  \multicolumn{2}{l}{\inptext{Pulse\_Test\_DoF}}
       && \multicolumn{2}{l}{\inp[nmlGrid::InternalGridDoF]{Internal\_Grid\_DoF}}
       \\
       &  \multicolumn{2}{l}{\inptext{Shu\_Osher\_Test\_DoF}}
       && \multicolumn{2}{l}{\inp[nmlGrid::InternalGridDoF]{Internal\_Grid\_DoF}}
       \\
       &  \multicolumn{2}{l}{\inptext{Sod\_Test\_DoF}}
       && \multicolumn{2}{l}{\inp[nmlGrid::InternalGridDoF]{Internal\_Grid\_DoF}}
       \\
       &  \multicolumn{2}{l}{\inptext{VortexGrid\_DoF}}
       && \multicolumn{2}{l}{\inp[nmlGrid::InternalGridDoF]{Internal\_Grid\_DoF}}
       \\
       &  \multicolumn{2}{l}{\inptext{Internal\_Grid\_Order}}
       && \multicolumn{2}{l}{\inp[nmlGrid::InternalGridOrder]{Same}[Internal\_Grid\_Order]}
       \\
       &  \multicolumn{2}{l}{\inptext{Metis\_Option}}
       && \multicolumn{2}{l}{\inp[nmlGrid::MetisOption]{Same}[Metis\_Option]}
       \\
       &  \multicolumn{2}{l}{\inptext{Geom\_Wts(1)\%Tria}}
       && \multicolumn{2}{l}{\inp[nmlGrid::GeomWts/Tria]{Same}[Geom\_Wts(:)\%Tria]}
       \\
       &  \multicolumn{2}{l}{\inptext{Geom\_Wts(1)\%Quad}}
       && \multicolumn{2}{l}{\inp[nmlGrid::GeomWts/Quad]{Same}[Geom\_Wts(:)\%Quad]}
       \\
       &  \multicolumn{2}{l}{\inptext{Geom\_Wts(1)\%Tetr}}
       && \multicolumn{2}{l}{\inp[nmlGrid::GeomWts/Tetr]{Same}[Geom\_Wts(:)\%Tetr]}
       \\
       &  \multicolumn{2}{l}{\inptext{Geom\_Wts(1)\%Pyra}}
       && \multicolumn{2}{l}{\inp[nmlGrid::GeomWts/Pyra]{Same}[Geom\_Wts(:)\%Pyra]}
       \\
       &  \multicolumn{2}{l}{\inptext{Geom\_Wts(1)\%Pris}}
       && \multicolumn{2}{l}{\inp[nmlGrid::GeomWts/Pris]{Same}[Geom\_Wts(:)\%Pris]}
       \\
       &  \multicolumn{2}{l}{\inptext{Geom\_Wts(1)\%Hexa}}
       && \multicolumn{2}{l}{\inp[nmlGrid::GeomWts/Hexa]{Same}[Geom\_Wts(:)\%Hexa]}
       \\
%==========================================================================
       \multicolumn{3}{l}{\&input}
    &  \multicolumn{2}{l}{\&\texttt{nml\_reference}}
    \\* %%% Note: \\* prevents pagebreak.
       &  \multicolumn{2}{l}{\inptext{Machref}}
       && \multicolumn{2}{l}{\inp[nmlReference::Mach]{Mach}}
       \\
       &  \multicolumn{2}{l}{\inptext{Reyref}}
       && \multicolumn{2}{l}{\inp[nmlReference::Reynolds]{Reynolds}}
       \\
       &  \multicolumn{2}{l}{\inptext{Gam}}
       && \multicolumn{2}{l}{\inp[nmlReference::Gamma]{Gamma}}
       \\
       &  \multicolumn{2}{l}{\inptext{Rhoref}}
       && \multicolumn{2}{l}{\inp[nmlReference::Density]{Density}}
       \\
       &  \multicolumn{2}{l}{\inptext{Tref}}
       && \multicolumn{2}{l}{\inp[nmlReference::Temperature]{Temperature}}
       \\
       &  \multicolumn{2}{l}{\inptext{Pref}}
       && \multicolumn{2}{l}{\inp[nmlReference::Pressure]{Pressure}}
       \\
       &  \multicolumn{2}{l}{\inptext{Rgasref}}
       && \multicolumn{2}{l}{\inp[nmlReference::GasConstant]{Gas\_Constant}}
       \\
       &  \multicolumn{2}{l}{\inptext{Alpha\_AoAref}}
       && \multicolumn{2}{l}{\inp[nmlReference::AoAAlpha]{AoA\_Alpha}}
       \\
       &  \multicolumn{2}{l}{\inptext{Beta\_AoAref}}
       && \multicolumn{2}{l}{\inp[nmlReference::AoABeta]{AoA\_Beta}}
       \\
       &  \multicolumn{2}{l}{\inptext{Pr}}
       && \multicolumn{2}{l}{\inp[nmlReference::Prandtl]{Prandtl}}
       \\
       &  \multicolumn{2}{l}{\inptext{Schmidt\_Lam}}
       && \multicolumn{2}{l}{\inp[nmlReference::Schmidt]{Schmidt}}
       \\
       &  \multicolumn{2}{l}{\inptext{Suth\_Muref}}
       && \multicolumn{2}{l}{\inp[nmlReference::SuthMuref]{Same}[Suth\_Muref]}
       \\
       &  \multicolumn{2}{l}{\inptext{Suth\_Tref}}
       && \multicolumn{2}{l}{\inp[nmlReference::SuthTref]{Same}[Suth\_Tref]}
       \\
       &  \multicolumn{2}{l}{\inptext{Suth\_Sref}}
       && \multicolumn{2}{l}{\inp[nmlReference::SuthSref]{Same}[Suth\_Sref]}
       \\
%==========================================================================
       \multicolumn{3}{l}{\&input}
    &  \multicolumn{2}{l}{\&\texttt{nml\_physics}}
    \\* %%% Note: \\* prevents pagebreak.
       &  \multicolumn{2}{l}{\inptext{Governing\_Equations}}
       && \multicolumn{2}{l}{\inp[nmlPhysics::GoverningEquations]{Same}[Governing\_Equations]}
       \\
       &  \multicolumn{2}{l}{\inptext{Moving\_Ref\_Frame\%Rotation(1:3)}}
       && \multicolumn{2}{l}{\inp[nmlPhysics::MovingRefFrame/Rotation]{Same}[Moving\_Ref\_Frame\%Rotation]}
       \\
       &  \multicolumn{2}{l}{\inptext{Moving\_Ref\_Frame\%Translation(1:3)}}
       && \multicolumn{2}{l}{\inp[nmlPhysics::MovingRefFrame/Translation]{Same}[Moving\_Ref\_Frame\%Translation]}
       \\
       &  \multicolumn{2}{l}{\inptext{Moving\_Ref\_Frame\%Spinup\%Iter}}
       && \multicolumn{2}{l}{\inp[nmlPhysics::MovingRefFrame/Spinup/Iter]{Same}[Moving\_Ref\_Frame\%Spinup\%Iter]}
       \\
       &  \multicolumn{2}{l}{\inptext{Moving\_Ref\_Frame\%Spinup\%Tstar}}
       && \multicolumn{2}{l}{\inp[nmlPhysics::MovingRefFrame/Spinup/Tstar]{Same}[Moving\_Ref\_Frame\%Spinup\%Tstar]}
       \\
%==========================================================================
       \multicolumn{3}{l}{\&input}
    &  \multicolumn{2}{l}{\&\texttt{nml\_initialization}}
    \\* %%% Note: \\* prevents pagebreak.
       &  \multicolumn{2}{l}{\inptext{Itestcase}}
       && \multicolumn{2}{l}{\inp[nmlInitialization::Initialization]{Initialization}}
       \\
       &  \multicolumn{2}{l}{\inptext{Jet\_Init\_Func}}
       && \multicolumn{2}{l}{\inp[nmlInitialization::JetInitFunc]{Same}[Jet\_Init\_Func]}
       \\
       &  \multicolumn{2}{l}{\inptext{Load\_Restart\_File}}
       && \multicolumn{2}{l}{\inp[nmlInitialization::LoadRestartFile]{Same}[Load\_Restart\_File]}
       \\
       &  \multicolumn{2}{l}{\inptext{Restart\_File}}
       && \multicolumn{2}{l}{\inp[nmlInitialization::RestartFile]{Same}[Restart\_File]}
       \\
       &  \multicolumn{2}{l}{\inptext{Initial\_Solution\_File}}
       && \multicolumn{2}{l}{\inp[nmlInitialization::InitialSolutionFile]{Solution\_File}}
       \\
       &  \multicolumn{2}{l}{\inptext{Initial\_Solution\_Format}}
       && \multicolumn{2}{l}{\inp[nmlInitialization::InitialSolutionFormat]{Solution\_Format}}
       \\
%==========================================================================
       \multicolumn{3}{l}{\&input}
    &  \multicolumn{2}{l}{\&\texttt{nml\_timestep}}
    \\* %%% Note: \\* prevents pagebreak.
       &  \multicolumn{2}{l}{\inptext{Runge\_Kutta\_Scheme}}
       && \multicolumn{2}{l}{\inp[nmlTimestep::RungeKuttaScheme]{Same}[Runge\_Kutta\_Scheme]}
       \\
       &  \multicolumn{2}{l}{\inptext{Minimum\_Timesteps}}
       && \multicolumn{2}{l}{\inp[nmlTimestep::NumTimestepsMin]{Num\_Timesteps\_Min}}
       \\
       &  \multicolumn{2}{l}{\inptext{Num\_Timesteps}}
       && \multicolumn{2}{l}{\inp[nmlTimestep::NumTimesteps]{Same}[Num\_Timesteps]}
       \\
       &  \multicolumn{2}{l}{\inptext{Final\_Time}}
       && \multicolumn{2}{l}{\inp[nmlTimestep::FinalTime]{Same}[Final\_Time]}
       \\
       &  \multicolumn{2}{l}{\inptext{Timestep\_Type}}
       && \multicolumn{2}{l}{\inp[nmlTimestep::TimestepType]{Same}[Timestep\_Type]}
       \\
       &  \multicolumn{2}{l}{\inptext{Constant\_Dt}}
       && \multicolumn{2}{l}{\inp[nmlTimestep::Timestep]{Timestep}}
       \\
       &  \multicolumn{2}{l}{\inptext{CFL}}
       && \multicolumn{2}{l}{\inp[nmlTimestep::CFL]{Same}[CFL]}
       \\
       &  \multicolumn{2}{l}{\inptext{CFL\_Beg}}
       && \multicolumn{2}{l}{\inp[nmlTimestep::CFLRampingBeg]{CFL\_Ramping\_Beg}}
       \\
       &  \multicolumn{2}{l}{\inptext{CFL\_End}}
       && \multicolumn{2}{l}{\inp[nmlTimestep::CFLRampingEnd]{CFL\_Ramping\_End}}
       \\
       &  \multicolumn{2}{l}{\inptext{CFL\_Cycles}}
       && \multicolumn{2}{l}{\inp[nmlTimestep::CFLRampingIterations]{CFL\_Ramping\_Iterations}}
       \\
       &  \multicolumn{2}{l}{\inptext{CFL\_Start\_Iter}}
       && \multicolumn{2}{l}{\inp[nmlTimestep::CFLRampingStartIter]{CFL\_Ramping\_Start\_Iter}}
       \\
       &  \multicolumn{2}{l}{\inptext{Convergence\_Order\_Abs\_Res}}
       && \multicolumn{2}{l}{\inp[nmlTimestep::ConvergenceOrderAbsRes]{Same}[Convergence\_Order\_Abs\_Res]}
       \\
       &  \multicolumn{2}{l}{\inptext{Convergence\_Order\_Max\_Res}}
       && \multicolumn{2}{l}{\inp[nmlTimestep::ConvergenceOrderMaxRes]{Same}[Convergence\_Order\_Max\_Res]}
       \\
%==========================================================================
       \multicolumn{3}{l}{\&input}
    &  \multicolumn{2}{l}{\&\texttt{nml\_bc}}
    \\* %%% Note: \\* prevents pagebreak.
       &  \multicolumn{2}{l}{\inptext{Walls\_Are\_Exact}}
       && \multicolumn{2}{l}{\inp[nmlBC::WallsAreExact]{Same}[Walls\_Are\_Exact]}
       \\
       &  \multicolumn{2}{l}{\inptext{Sub\_Inflow\_Method}}
       && \multicolumn{2}{l}{\inp[nmlBC::SubInflowMethod]{Same}[Sub\_Inflow\_Method]}
       \\
       &  \multicolumn{2}{l}{\inptext{Use\_BC\_Conflicts}}
       && \multicolumn{2}{l}{\inp[nmlBC::UseBCConflicts]{Same}[Use\_BC\_Conflicts]}
       \\
       &  \multicolumn{2}{l}{\inptext{BC\_Input(1)}}
       && \multicolumn{2}{l}{\inpBC[nmlBC::BCInput]{}}
       \\
       &  \multicolumn{2}{l}{\inptext{\%Name}}
       && \multicolumn{2}{l}{\inpBC[nmlBC::BCInput/Name]{Same}[Name]}
       \\
       &  \multicolumn{2}{l}{\inptext{\%Set\_BC\_Value}}
       && \multicolumn{2}{l}{\inpBC[nmlBC::BCInput/SetBCValue]{Same}[Set\_BC\_Value]}
       \\
       &  \multicolumn{2}{l}{\inptext{\%T\_Static}}
       && \multicolumn{2}{l}{\inpBC[nmlBC::BCInput/Temperature]{\%Temperature}[Temperature][]}
       \\
       &  \multicolumn{2}{l}{\inptext{\%P\_Static}}
       && \multicolumn{2}{l}{\inpBC[nmlBC::BCInput/Pressure]{\%Pressure}[Pressure][]}
       \\
       &  \multicolumn{2}{l}{\inptext{\%Rho\_Static}}
       && \multicolumn{2}{l}{\inpBC[nmlBC::BCInput/Density]{\%Density}[Density][]}
       \\
       &  \multicolumn{2}{l}{\inptext{\%Mach}}
       && \multicolumn{2}{l}{\inpBC[nmlBC::BCInput/Mach]{Same}[Mach]}
       \\
       &  \multicolumn{2}{l}{\inptext{\%Vx}}
       && \multicolumn{2}{l}{\inpBC[nmlBC::BCInput/VelocityX]{\%VelocityX}[VelocityX][]}
       \\
       &  \multicolumn{2}{l}{\inptext{\%Vy}}
       && \multicolumn{2}{l}{\inpBC[nmlBC::BCInput/VelocityY]{\%VelocityY}[VelocityY][]}
       \\
       &  \multicolumn{2}{l}{\inptext{\%Vz}}
       && \multicolumn{2}{l}{\inpBC[nmlBC::BCInput/VelocityZ]{\%VelocityZ}[VelocityZ][]}
       \\
       &  \multicolumn{2}{l}{\inptext{\%T\_Total}}
       && \multicolumn{2}{l}{\inpBC[nmlBC::BCInput/TotalTemperature]{\%Total\_Temperature}[Total\_Temperature][]}
       \\
       &  \multicolumn{2}{l}{\inptext{\%P\_Total}}
       && \multicolumn{2}{l}{\inpBC[nmlBC::BCInput/TotalPressure]{\%Total\_Pressure}[Total\_Pressure][]}
       \\
       &  \multicolumn{2}{l}{\inptext{\%Rho\_Total}}
       && \multicolumn{2}{l}{\inpBC[nmlBC::BCInput/TotalDensity]{\%Total\_Density}[Total\_Density][]}
       \\
       &  \multicolumn{2}{l}{\inptext{\%Alpha\_AoA}}
       && \multicolumn{2}{l}{\inpBC[nmlBC::BCInput/AoAAlpha]{\%AoA\_Alpha}[AoA\_Alpha][]}
       \\
       &  \multicolumn{2}{l}{\inptext{\%Beta\_AoA}}
       && \multicolumn{2}{l}{\inpBC[nmlBC::BCInput/AoABeta]{\%AoA\_Beta}[AoA\_Beta][]}
       \\
       &  \multicolumn{2}{l}{\inptext{\%Wall\_Temp}}
       && \multicolumn{2}{l}{\inpBC[nmlBC::BCInput/WallTemp]{Same}[Wall\_Temp]}
       \\
       &  \multicolumn{2}{l}{\inptext{\%Rotation\_Rate(1:3)}}
       && \multicolumn{2}{l}{\inpBC[nmlBC::BCInput/RotationRate]{Same}[Rotation\_Rate]}
       \\
       &  \multicolumn{2}{l}{\inptext{\%Is\_Specified\_In\_MRF}}
       && \multicolumn{2}{l}{\inpBC[nmlBC::BCInput/IsSpecifiedInMRF]{Same}[Is\_Specified\_In\_MRF]}
       \\
       &  \multicolumn{2}{l}{\inptext{\%Mass\_Frac(1:7)}}
       && \multicolumn{2}{l}{\inpBC[nmlBC::BCInput/MassFrac]{Same}[Mass\_Frac]}
       \\
       &  \multicolumn{2}{l}{\inptext{\%Species\_Name(1:7)}}
       && \multicolumn{2}{l}{\inpBC[nmlBC::BCInput/SpeciesName]{Same}[Species\_Name]}
       \\
       &  \multicolumn{2}{l}{\inptext{\%Relax\%Value}}
       && \multicolumn{2}{l}{\inpBC[nmlBC::BCInput/Relax/Value]{Same}[Relax\%Value]}
       \\
       &  \multicolumn{2}{l}{\inptext{\%Relax\%Time}}
       && \multicolumn{2}{l}{\inpBC[nmlBC::BCInput/Relax/Time]{Same}[Relax\%Time]}
       \\
       &  \multicolumn{2}{l}{\inptext{\%Relax\%Iter}}
       && \multicolumn{2}{l}{\inpBC[nmlBC::BCInput/Relax/Iter]{Same}[Relax\%Iter]}
       \\
       &  \multicolumn{2}{l}{\inptext{\%SEM\%Profile\_File}}
       && \multicolumn{2}{l}{\inpBC[nmlBC::BCInput/SEM/ProfileFile]{Same}[SEM\%Profile\_File]}
       \\
       &  \multicolumn{2}{l}{\inptext{\%SEM\%Is\_Cylindrical}}
       && \multicolumn{2}{l}{\inpBC[nmlBC::BCInput/SEM/IsCylindrical]{Same}[SEM\%Is\_Cylindrical]}
       \\
       &  \multicolumn{2}{l}{\inptext{\%SEM\%Length\_Scale}}
       && \multicolumn{2}{l}{\inpBC[nmlBC::BCInput/SEM/LengthScale]{Same}[SEM\%Length\_Scale]}
       \\
       &  \multicolumn{2}{l}{\inptext{\%SEM\%Convective\_Velocity}}
       && \multicolumn{2}{l}{\inpBC[nmlBC::BCInput/SEM/ConvectiveVelocity]{Same}[SEM\%Convective\_Velocity]}
       \\
       &  \multicolumn{2}{l}{\inptext{\%SEM\%Profile\_Height}}
       && \multicolumn{2}{l}{\inpBC[nmlBC::BCInput/SEM/ProfileHeight]{Same}[SEM\%Profile\_Height]}
       \\
       &  \multicolumn{2}{l}{\inptext{\%SEM\%Sub\_In\_BC\_Grp}}
       && \multicolumn{2}{l}{\inpBC[nmlBC::BCInput/SEM/SubInBCGrp]{Same}[SEM\%Sub\_In\_BC\_Grp]}
       \\
       &  \multicolumn{2}{l}{\inptext{\%SEM\%Sub\_In\_BC\_Idx}}
       && \multicolumn{2}{l}{\inpBC[nmlBC::BCInput/SEM/SubInBCIdx]{Same}[SEM\%Sub\_In\_BC\_Idx]}
       \\
       &  \multicolumn{2}{l}{\inptext{\%Periodic\%Matching\_BC\_Name}}
       && \multicolumn{2}{l}{\inpBC[nmlBC::BCInput/Periodic/MatchingBCName]{Same}[Periodic\%Matching\_BC\_Name]}
       \\
       &  \multicolumn{2}{l}{\inptext{\%Periodic\%Matching\_BC\_Idx}}
       && \multicolumn{2}{l}{\inpBC[nmlBC::BCInput/Periodic/MatchingBCIdx]{Same}[Periodic\%Matching\_BC\_Idx]}
       \\
       &  \multicolumn{2}{l}{\inptext{\%Periodic\%RotationCenter(1:3)}}
       && \multicolumn{2}{l}{\inpBC[nmlBC::BCInput/Periodic/RotationCenter]{Same}[Periodic\%RotationCenter]}
       \\
       &  \multicolumn{2}{l}{\inptext{\%Periodic\%RotationAngle(1:3)}}
       && \multicolumn{2}{l}{\inpBC[nmlBC::BCInput/Periodic/RotationAngle]{Same}[Periodic\%RotationAngle]}
       \\
       &  \multicolumn{2}{l}{\inptext{\%Periodic\%RotationAxis(1:4)}}
       && \multicolumn{2}{l}{\inpBC[nmlBC::BCInput/Periodic/RotationAxis]{Same}[Periodic\%RotationAxis]}
       \\
       &  \multicolumn{2}{l}{\inptext{\%Periodic\%Translation(1:3)}}
       && \multicolumn{2}{l}{\inpBC[nmlBC::BCInput/Periodic/Translation]{Same}[Periodic\%Translation]}
       \\
%==========================================================================
       \multicolumn{3}{l}{\&input}
    &  \multicolumn{2}{l}{\&\texttt{nml\_numerics}}
    \\* %%% Note: \\* prevents pagebreak.
       &  \multicolumn{2}{l}{\inptext{All\_Orders}}
       && \multicolumn{2}{l}{\inp[nmlNumerics::AllOrders]{Same}[All\_Orders]}
       \\
       &  \multicolumn{2}{l}{\inptext{Solution\_Order}}
       && \multicolumn{2}{l}{\inp[nmlNumerics::SolutionOrder]{Same}[Solution\_Order]}
       \\
       &  \multicolumn{2}{l}{\inptext{Projection\_Order}}
       && \multicolumn{2}{l}{\inp[nmlNumerics::ProjectionOrder]{Same}[Projection\_Order]}
       \\
       &  \multicolumn{2}{l}{\inptext{Quadrature\_Order}}
       && \multicolumn{2}{l}{\inp[nmlNumerics::QuadratureOrder]{Same}[Quadrature\_Order]}
       \\
       &  \multicolumn{2}{l}{\inptext{Error\_Order}}
       && \multicolumn{2}{l}{\inp[nmlNumerics::ErrorOrder]{Same}[Error\_Order]}
       \\
       &  \multicolumn{2}{l}{\inptext{All\_Points}}
       && \multicolumn{2}{l}{\inp[nmlNumerics::AllPoints]{Same}[All\_Points]}
       \\
       &  \multicolumn{2}{l}{\inptext{Loc\_Solution\_Pts}}
       && \multicolumn{2}{l}{\inp[nmlNumerics::LocSolutionPts]{Same}[Loc\_Solution\_Pts]}
       \\
       &  \multicolumn{2}{l}{\inptext{Loc\_Flux\_Pts}}
       && \multicolumn{2}{l}{\inp[nmlNumerics::LocFluxPts]{Same}[Loc\_Flux\_Pts]}
       \\
       &  \multicolumn{2}{l}{\inptext{Loc\_Quadrature\_Pts}}
       && \multicolumn{2}{l}{\inp[nmlNumerics::LocQuadraturePts]{Same}[Loc\_Quadrature\_Pts]}
       \\
       &  \multicolumn{2}{l}{\inptext{Loc\_Error\_Pts}}
       && \multicolumn{2}{l}{\inp[nmlNumerics::LocErrorPts]{Same}[Loc\_Error\_Pts]}
       \\
       &  \multicolumn{2}{l}{\inptext{Loc\_Triangle\_Pts}}
       && \multicolumn{2}{l}{\inp[nmlNumerics::LocTrianglePts]{Same}[Loc\_Triangle\_Pts]}
       \\
       &  \multicolumn{2}{l}{\inptext{Loc\_Tetrahedra\_Pts}}
       && \multicolumn{2}{l}{\inp[nmlNumerics::LocTetrahedraPts]{Same}[Loc\_Tetrahedra\_Pts]}
       \\
       &  \multicolumn{2}{l}{\inptext{Invs\_Flux\_Method}}
       && \multicolumn{2}{l}{\inp[nmlNumerics::InvsFluxMethod]{Same}[Invs\_Flux\_Method]}
       \\
       &  \multicolumn{2}{l}{\inptext{Visc\_Flux\_Method}}
       && \multicolumn{2}{l}{\inp[nmlNumerics::ViscFluxMethod]{Same}[Visc\_Flux\_Method]}
       \\
       &  \multicolumn{2}{l}{\inptext{Visc\_Variables\_To\_Ave}}
       && \multicolumn{2}{l}{\inp[nmlNumerics::ViscVariablesToAve]{Same}[Visc\_Variables\_To\_Ave]}
       \\
       &  \multicolumn{2}{l}{\inptext{Combine\_Common\_Flux\_Projections}}
       && \multicolumn{2}{l}{\inp[nmlNumerics::CombineCommonFluxProjections]{Same}[Combine\_Common\_Flux\_Projections]}
       \\
       &  \multicolumn{2}{l}{\inptext{Correction\_Function}}
       && \multicolumn{2}{l}{\inp[nmlNumerics::CorrectionFunction]{Same}[Correction\_Function]}
       \\
%==========================================================================
       \multicolumn{3}{l}{\&input}
    &  \multicolumn{2}{l}{\&\texttt{nml\_limiter}}
    \\* %%% Note: \\* prevents pagebreak.
       &  \multicolumn{2}{l}{\inptext{Filtering\_Interval}}
       && \multicolumn{2}{l}{\inp[nmlLimiter::GlobalFilterInterval]{Global\_Filter\_Interval}}
       \\
       &  \multicolumn{2}{l}{\inptext{Filter\_Option}}
       && \multicolumn{2}{l}{\inp[nmlLimiter::GlobalFilterOption]{Global\_Filter\_Option}}
       \\
       &  \multicolumn{2}{l}{\inptext{Filter\_Order}}
       && \multicolumn{2}{l}{\inp[nmlLimiter::GlobalFilterOrder]{Global\_Filter\_Order}}
       \\
       &  \multicolumn{2}{l}{\inptext{Filter\_Alpha}}
       && \multicolumn{2}{l}{\inp[nmlLimiter::GlobalFilterAlpha]{Global\_Filter\_Alpha}}
       \\
       &  \multicolumn{2}{l}{\inptext{Filter\_Etac}}
       && \multicolumn{2}{l}{\inp[nmlLimiter::GlobalFilterEtac]{Global\_Filter\_Etac}}
       \\
       &  \multicolumn{2}{l}{\inptext{Limiting\_Interval}}
       && \multicolumn{2}{l}{\inp[nmlLimiter::LocalLimiterInterval]{Local\_Limiter\_Interval}}
       \\
       &  \multicolumn{2}{l}{\inptext{Limiter\_Option}}
       \\* %%% Note: \\* prevents pagebreak.
       && \multicolumn{4}{p{5.0in}}{%
          The old \inptext{Limiter\_Option} input controlled
          both the marker and limiter.
          \inp[nmlLimiter::LocalLimiterOption]{Local\_Limiter\_Option}
          now controls only the limiter, and
          \inp[nmlLimiter::MarkerOption]{Marker\_Option}
          now controls the marker.
          Both are shown below.
          }
       \\
       &  \multicolumn{2}{l}{\inptext{Limiter\_Option}}
       && \multicolumn{2}{l}{\inp[nmlLimiter::LocalLimiterOption]{Local\_Limiter\_Option}}
       \\* %%% Note: \\* prevents pagebreak.
      %&& \multicolumn{2}{p{\lenG}}{This input has been split into
      %   \inp[nmlLimiter::MarkerOption]{Marker\_Option}
      %   and
      %   \inp[nmlLimiter::LocalLimiterOption]{Local\_Limiter\_Option}.
      %   }
      %&&
      %\\
      %&& \multicolumn{3}{p{\lenH}}{This input has been split into
      %   \inp[nmlLimiter::MarkerOption]{Marker\_Option}
      %   and
      %   \inp[nmlLimiter::LocalLimiterOption]{Local\_Limiter\_Option}.
      %   }
      %&
      %\\
      %&& This input has been split into
      %   \inp[nmlLimiter::MarkerOption]{Marker\_Option}
      %   and
      %   \inp[nmlLimiter::LocalLimiterOption]{Local\_Limiter\_Option}.
      %   &&&
      %\\
       && \multicolumn{2}{l}{
          \begin{tabular}[t]{rl}
            -2 & Jump Marker + Scaled Projection \\
            -1 & Res. Marker + Scaled Projection \\
             0 & None \\
             1 & Res. Marker + {\small\inptext{Limiter\_Filter\_Option}} \\
               & if ( {\small\inptext{Limiter\_Filter\_Option}} == 0 ) then \\
               & \quad use Recursive Projection \\
             2 & Jump Marker + {\small\inptext{Limiter\_Filter\_Option}} \\
             3 & MLP  Marker + MLP limiter \\
          \end{tabular}
          }
       &  \multicolumn{2}{l}{
          \begin{tabular}[t]{rl}
            0 & None \\
            1 & Filtering \\
            2 & Scaled Filtering \\
            3 & MLP \\
            9001 & Recursive Projection \\
            9002 & Scaled    Projection
          \end{tabular}
          }
       \\
       &  \multicolumn{2}{l}{\inptext{Limiter\_Option}}
       && \multicolumn{2}{l}{\inp[nmlLimiter::MarkerOption]{Marker\_Option}}
       \\* %%% Note: \\* prevents pagebreak.
       && \multicolumn{2}{l}{See description above.}
       &  \multicolumn{2}{l}{
          \begin{tabular}[t]{rl}
              0 & Default based on \\
                & {\small\inptext{Local\_Limiter\_Option}} \\
              1 & Resolution marker \\
              2 & Jump       marker \\
              3 & MLP        marker
          \end{tabular}
          }
       \\
       &  \multicolumn{2}{l}{\inptext{Limiter\_Use\_Dynamic\_Filter}}
       && \multicolumn{2}{l}{\inp[nmlLimiter::LocalLimiterOption]{Local\_Limiter\_Option}}
       \\
       && \multicolumn{4}{l}{Input value has a different meaning.
                             Dynamic filtering now called scaled filtering.}
       \\
       &  \multicolumn{2}{l}{\inptext{Limiter\_Filter\_Order}}
       && \multicolumn{2}{l}{\inp[nmlLimiter::LocalFilterOrder]{Local\_Filter\_Order}}
       \\
       &  \multicolumn{2}{l}{\inptext{Limiter\_Filter\_Alpha}}
       && \multicolumn{2}{l}{\inp[nmlLimiter::LocalFilterAlpha]{Local\_Filter\_Alpha}}
       \\
       &  \multicolumn{2}{l}{\inptext{Limiter\_Filter\_Etac}}
       && \multicolumn{2}{l}{\inp[nmlLimiter::LocalFilterEtac]{Local\_Filter\_Etac}}
       \\
       &  \multicolumn{2}{l}{\inptext{*}}
       && \multicolumn{2}{l}{\inp[nmlLimiter::SpongeFilterInterval]{Sponge\_Filter\_Interval}}
       \\
       &  \multicolumn{2}{l}{\inptext{Use\_Sponge\_Filter}}
       && \multicolumn{2}{l}{\inp[nmlLimiter::SpongeFilterOption]{Sponge\_Filter\_Option}}
       \\* %%% Note: \\* prevents pagebreak.
%      && \multicolumn{4}{l}{Input value has changed from logical to integer.}
%      \\
       && \begin{tabular}[t]{ll}
            F & Off \\
            T & On
          \end{tabular}
       &&  \multicolumn{2}{l}{
           \begin{tabular}[t]{ll}
             0 & Off \\
             1 & On
           \end{tabular}
           }
       \\
       &  \multicolumn{2}{l}{\inptext{Sponge\_Filter\_Order}}
       && \multicolumn{2}{l}{\inp[nmlLimiter::SpongeFilterOrder]{Same}[Sponge\_Filter\_Order]}
       \\
       &  \multicolumn{2}{l}{\inptext{Sponge\_Filter\_Alpha}}
       && \multicolumn{2}{l}{\inp[nmlLimiter::SpongeFilterAlpha]{Same}[Sponge\_Filter\_Alpha]}
       \\
       &  \multicolumn{2}{l}{\inptext{Sponge\_Filter\_Etac}}
       && \multicolumn{2}{l}{\inp[nmlLimiter::SpongeFilterEtac]{Same}[Sponge\_Filter\_Etac]}
       \\
       &  \multicolumn{2}{l}{\inptext{Sponge\_Beyond\_X}}
       && \multicolumn{2}{l}{\inp[nmlLimiter::SpongeFilterXStart]{Sponge\_Filter\_X\_Start}}
       \\
%%% Moved Marker_Option under Limiter_Option above.
%      &  \multicolumn{2}{l}{\inptext{*}}
%      && \multicolumn{2}{l}{\inp[nmlLimiter::MarkerOption]{Marker\_Option}}
%      \\
%      &&
%      && \multicolumn{2}{l}{
%         \begin{tabular}[t]{rl}
%             0 & Default based on {\small\inptext{Local\_Limiter\_Option}} \\
%             1 & Resolution marker \\
%             2 & Jump       marker \\
%             3 & MLP        marker
%         \end{tabular}
%         }
%      \\
       &  \multicolumn{2}{l}{\inptext{Resolution\_Marker\_Variable}}
       && \multicolumn{2}{l}{\inp[nmlLimiter::MarkerVariable]{Marker\_Variable}}
       \\* %%% Note: \\* prevents pagebreak.
%      && \multicolumn{4}{l}{Input value has a different meaning.}
%      \\
       && \begin{tabular}[t]{rl}
              0 & All conserved variables \\
            1:5 & Specific conserved variable
          \end{tabular}
       && \multicolumn{2}{l}{
          \begin{tabular}[t]{rl}
             -2 & Ducros (velocity divergence) \\
             -1 & Pressure \\
              0 & Default based on {\small\inptext{Marker\_Option}} \\
            1:5 & Specific conserved variable \\
              6 & All conserved variables
          \end{tabular}
          }
       \\
       &  \multicolumn{2}{l}{\inptext{Jump\_Marker\_Variable}}
       && \multicolumn{2}{l}{\inp[nmlLimiter::MarkerVariable]{Marker\_Variable}}
       \\* %%% Note: \\* prevents pagebreak.
%      && \multicolumn{4}{l}{Input value has a different meaning.}
%      \\
       && \begin{tabular}[t]{rl}
              0 & Pressure \\
            1:5 & Specific conserved variable
          \end{tabular}
       && \multicolumn{2}{l}{
          \begin{tabular}[t]{rl}
             -2 & Ducros (velocity divergence) \\
             -1 & Pressure \\
              0 & Default based on {\small\inptext{Marker\_Option}} \\
            1:5 & Specific conserved variable \\
              6 & All conserved variables
          \end{tabular}
          }
       \\
%      \multicolumn{6}{l}{\phantom{blank line}} \\
%      \multicolumn{6}{l}{\phantom{blank line}} \\
%      \multicolumn{6}{l}{\phantom{blank line}} \\
%      \multicolumn{6}{l}{\phantom{blank line}} \\
       &  \multicolumn{2}{l}{\inptext{MLP\_Marker\_Variable}}
       && \multicolumn{2}{l}{\inp[nmlLimiter::MarkerVariable]{Marker\_Variable}}
       \\* %%% Note: \\* prevents pagebreak.
%      && \multicolumn{4}{l}{Input value has a different meaning.}
%      \\
       && \begin{tabular}[t]{rl}
              0 & Density \\
            1:5 & Specific conserved variable
          \end{tabular}
       && \multicolumn{2}{l}{
          \begin{tabular}[t]{rl}
             -2 & Ducros (velocity divergence) \\
             -1 & Pressure \\
              0 & Default based on {\small\inptext{Marker\_Option}} \\
            1:5 & Specific conserved variable \\
              6 & All conserved variables
          \end{tabular}
          }
       \\
       &  \multicolumn{2}{l}{\inptext{Limiter\_Marker\_Min}}
       && \multicolumn{2}{l}{\inp[nmlLimiter::MarkerMin]{Marker\_Min}}
       \\
       &  \multicolumn{2}{l}{\inptext{Limiter\_Marker\_Max}}
       && \multicolumn{2}{l}{\inp[nmlLimiter::MarkerMax]{Marker\_Max}}
       \\
       &  \multicolumn{2}{l}{\inptext{Limiter\_Marker\_Dpsi}}
       && \multicolumn{2}{l}{\inp[nmlLimiter::MarkerDpsi]{Marker\_Dpsi}}
       \\
       &  \multicolumn{2}{l}{\inptext{Limiting\_Interval\_Output}}
       && \multicolumn{2}{l}{\inp[nmlLimiter::MarkerStoredInterval]{Marker\_Stored\_Interval}}
       \\
       &  \multicolumn{2}{l}{\inptext{Thermofix\%Method}}
       && \multicolumn{2}{l}{\inp[nmlLimiter::Thermofix/Method]{Same}[Thermofix\%Method]}
       \\
       &  \multicolumn{2}{l}{\inptext{Thermofix\%Iterations}}
       && \multicolumn{2}{l}{\inp[nmlLimiter::Thermofix/Iterations]{Same}[Thermofix\%Iterations]}
       \\
       &  \multicolumn{2}{l}{\inptext{Thermofix\%Check\_Cells}}
       && \multicolumn{2}{l}{\inp[nmlLimiter::Thermofix/CheckCells]{Same}[Thermofix\%Check\_Cells]}
       \\
       &  \multicolumn{2}{l}{\inptext{Thermofix\%Check\_Faces}}
       && \multicolumn{2}{l}{\inp[nmlLimiter::Thermofix/CheckFaces]{Same}[Thermofix\%Check\_Faces]}
       \\
       &  \multicolumn{2}{l}{\inptext{Thermofix\%Check\_Nodes}}
       && \multicolumn{2}{l}{\inp[nmlLimiter::Thermofix/CheckNodes]{Same}[Thermofix\%Check\_Nodes]}
       \\
       &  \multicolumn{2}{l}{\inptext{Thermofix\%Verbose}}
       && \multicolumn{2}{l}{\inp[nmlLimiter::Thermofix/Verbose]{Same}[Thermofix\%Verbose]}
       \\
       &  \multicolumn{2}{l}{\inptext{Thermofix\%Scale\_Factor}}
       && \multicolumn{2}{l}{\inp[nmlLimiter::Thermofix/ScaleFactor]{Same}[Thermofix\%Scale\_Factor]}
       \\
       &  \multicolumn{2}{l}{\inptext{Thermofix\%Min\_Pressure}}
       && \multicolumn{2}{l}{\inp[nmlLimiter::Thermofix/MinPressure]{Same}[Thermofix\%MinPressure]}
       \\
       &  \multicolumn{2}{l}{\inptext{Thermofix\%Min\_Temperature}}
       && \multicolumn{2}{l}{\inp[nmlLimiter::Thermofix/MinTemperature]{Same}[Thermofix\%Min\_Temperature]}
       \\
       &  \multicolumn{2}{l}{\inptext{Thermofix\%Min\_Density}}
       && \multicolumn{2}{l}{\inp[nmlLimiter::Thermofix/MinDensity]{Same}[Thermofix\%Min\_Density]}
       \\
%==========================================================================
       \multicolumn{3}{l}{\&input}
    &  \multicolumn{2}{l}{\&\texttt{nml\_time\_ave}}
    \\* %%% Note: \\* prevents pagebreak.
       &  \multicolumn{2}{l}{\inptext{Output\_Time\_Averaging}}
       && \multicolumn{2}{l}{\inp[nmlTimeAve::ComputeTimeAveraging]{Compute\_Time\_Averaging}}
       \\
       &  \multicolumn{2}{l}{\inptext{Num\_Saved\_Time\_Ave\_Restarts}}
       && \multicolumn{2}{l}{\inp[nmlTimeAve::NumSavedTimeAveRestarts]{Same}[Num\_Saved\_Time\_Ave\_Restarts]}
       \\
       &  \multicolumn{2}{l}{\inptext{Time\_Ave\_File}}
       && \multicolumn{2}{l}{\inp[nmlTimeAve::InputTimeAveFile]{Input\_Time\_Ave\_File}}
       \\
       &  \multicolumn{2}{l}{\inptext{Time\_Ave\_Vel\_Is\_Axisymm}}
       && \multicolumn{2}{l}{\inp[nmlTimeAve::TimeAveVelIsAxisymm]{Same}[Time\_Ave\_Vel\_Is\_Axisymm]}
       \\
       &  \multicolumn{2}{l}{\inptext{Start\_Averaging\_Timestep}}
       && \multicolumn{2}{l}{\inp[nmlTimeAve::StartAveragingTimestep]{Same}[Start\_Averaging\_Timestep]}
       \\
       &  \multicolumn{2}{l}{\inptext{Start\_Averaging\_Tstar}}
       && \multicolumn{2}{l}{\inp[nmlTimeAve::StartAveragingTstar]{Same}[Start\_Averaging\_Tstar]}
       \\
%==========================================================================
       \multicolumn{3}{l}{\&input}
    &  \multicolumn{2}{l}{\&\texttt{nml\_output}}
    \\* %%% Note: \\* prevents pagebreak.
       &  \multicolumn{2}{l}{\inptext{Time\_Scaling\_Factor}}
       && \multicolumn{2}{l}{\inp[nmlOutput::TimeScalingFactor]{Same}[Time\_Scaling\_Factor]}
       \\
       &  \multicolumn{2}{l}{\inptext{Restart\_Interval}}
       && \multicolumn{2}{l}{\inp[nmlOutput::RestartInterval]{Same}[Restart\_Interval]}
       \\
       &  \multicolumn{2}{l}{\inptext{Restart\_Frequency}}
       && \multicolumn{2}{l}{\inp[nmlOutput::RestartFrequency]{Same}[Restart\_Frequency]}
       \\
       &  \multicolumn{2}{l}{\inptext{Restart\_Output\_Uses\_Host\_Roots}}
       && \multicolumn{2}{l}{\inp[nmlOutput::RestartOutputUsesHostRoots]{Same}[Restart\_Output\_Uses\_Host\_Roots]}
       \\
       &  \multicolumn{2}{l}{\inptext{Completely\_Disable\_Restart\_Output}}
       && \multicolumn{2}{l}{\inp[nmlOutput::RestartOutputDisabled]{Restart\_Output\_Disabled}}
       \\
       &  \multicolumn{2}{l}{\inptext{Iter\_Out\_Interval}}
       && \multicolumn{2}{l}{\inp[nmlOutput::DisplayInterval]{Display\_Interval}}
       \\
       &  \multicolumn{2}{l}{\inptext{Results\_Interval}}
       && \multicolumn{2}{l}{\inp[nmlOutput::ResidualInterval]{Residual\_Interval}}
       \\
       &  \multicolumn{2}{l}{\inptext{Output\_Summary\_File}}
       && \multicolumn{2}{l}{\inp[nmlOutput::OutputSummaryFile]{Same}[Output\_Summary\_File]}
       \\
       &  \multicolumn{2}{l}{\inptext{Output\_Dir}}
       && \multicolumn{2}{l}{\inp[nmlOutput::OutputDir]{Same}[Output\_Dir]}
       \\
       &  \multicolumn{2}{l}{\inptext{Output\_Interval}}
       && \multicolumn{2}{l}{\inp[nmlOutput::OutputInterval]{Same}[Output\_Interval]}
       \\
       &  \multicolumn{2}{l}{\inptext{Output\_Frequency}}
       && \multicolumn{2}{l}{\inp[nmlOutput::OutputFrequency]{Same}[Output\_Frequency]}
       \\
       &  \multicolumn{2}{l}{\inptext{Output\_Order}}
       && \multicolumn{2}{l}{\inp[nmlOutput::OutputOrder]{Same}[Output\_Order]}
       \\
       &  \multicolumn{2}{l}{\inptext{Dmat\_Grad\_Temp\_Order}}
       && \multicolumn{2}{l}{\inp[nmlOutput::DmatGradTempOrder]{Same}[Dmat\_Grad\_Temp\_Order]}
       \\
       &  \multicolumn{2}{l}{\inptext{Loc\_Output\_Pts}}
       && \multicolumn{2}{l}{\inp[nmlOutput::LocOutputPts]{Same}[Loc\_Output\_Pts]}
       \\
       &  \multicolumn{2}{l}{\inptext{CGNS\_High\_Order\_Output}}
       && \multicolumn{2}{l}{\inp[nmlOutput::CGNSHighOrderOutput]{Same}[CGNS\_High\_Order\_Output]}
       \\
       &  \multicolumn{2}{l}{\inptext{CGNS\_Mixed\_Output}}
       && \multicolumn{2}{l}{\inp[nmlOutput::CGNSMixedOutput]{Same}[CGNS\_Mixed\_Output]}
       \\
       &  \multicolumn{2}{l}{\inptext{Completely\_Disable\_CGNS}}
       && \multicolumn{2}{l}{\inp[nmlOutput::CGNSOutputDisabled]{CGNS\_Output\_Disabled}}
       \\
       &  \multicolumn{2}{l}{\inptext{Probes\_Collection\_Interval}}
       && \multicolumn{2}{l}{\inp[nmlOutput::ProbesCollectionInterval]{same}[Probes\_Collection\_Interval]}
       \\
       &  \multicolumn{2}{l}{\inptext{Probes\_Write\_Interval}}
       && \multicolumn{2}{l}{\inp[nmlOutput::ProbesWriteInterval]{same}[Probes\_Write\_Interval]}
       \\
       &  \multicolumn{2}{l}{\inptext{Point\_Probes}}
       && \multicolumn{2}{l}{\inp[nmlOutput::PointProbes]{same}[Point\_Probes]}
       \\
       &  \multicolumn{2}{l}{\inptext{Rake\_Probes}}
       && \multicolumn{2}{l}{\inp[nmlOutput::RakeProbes]{same}[Rake\_Probes]}
       \\
       &  \multicolumn{2}{l}{\inptext{Profile\_IO\_All}}
       && \multicolumn{2}{l}{\inp[nmlOutput::ProfileIOAll]{Same}[Profile\_IO\_All]}
       \\
       &  \multicolumn{2}{l}{\inptext{Profile\_IO\_Restart}}
       && \multicolumn{2}{l}{\inp[nmlOutput::ProfileIORestart]{Same}[Profile\_IO\_Restart]}
       \\
       &  \multicolumn{2}{l}{\inptext{Profile\_IO\_CGNS}}
       && \multicolumn{2}{l}{\inp[nmlOutput::ProfileIOCGNS]{Same}[Profile\_IO\_CGNS]}
       \\
       &  \multicolumn{2}{l}{\inptext{Profile\_IO\_CGNS\_Subs}}
       && \multicolumn{2}{l}{\inp[nmlOutput::ProfileIOCGNSSubs]{Same}[Profile\_IO\_CGNS\_Subs]}
       \\
       &  \multicolumn{2}{l}{\inptext{Lustre\_Stripe\_Count}}
       && \multicolumn{2}{l}{\inp[nmlOutput::LustreStripeCount]{Same}[Lustre\_Stripe\_Count]}
       \\
       &  \multicolumn{2}{l}{\inptext{Lustre\_Stripe\_Size}}
       && \multicolumn{2}{l}{\inp[nmlOutput::LustreStripeSize]{Same}[Lustre\_Stripe\_Size]}
       \\
%==========================================================================
       \multicolumn{3}{l}{\&input}
    &  \multicolumn{2}{l}{\&\texttt{nml\_output\_variables}}
    \\* %%% Note: \\* prevents pagebreak.
       &  \multicolumn{2}{l}{\inptext{Output\_Conserved\_Variables}}
       && \multicolumn{2}{l}{\inp[nmlOutputVariables::ConservedVariables]{Conserved\_Variables}}
       \\
       &  \multicolumn{2}{l}{\inptext{Output\_Momentum}}
       && \multicolumn{2}{l}{\inp[nmlOutputVariables::Momentum]{Momentum}}
       \\
       &  \multicolumn{2}{l}{\inptext{Output\_Vorticity}}
       && \multicolumn{2}{l}{\inp[nmlOutputVariables::Vorticity]{Vorticity}}
       \\
       &  \multicolumn{2}{l}{\inptext{Output\_Cylindrical\_Velocity}}
       && \multicolumn{2}{l}{\inp[nmlOutputVariables::CylindricalVelocity]{Cylindrical\_Velocity}}
       \\
       &  \multicolumn{2}{l}{\inptext{Output\_Cylindrical\_Coordinates}}
       && \multicolumn{2}{l}{\inp[nmlOutputVariables::CylindricalCoordinates]{Cylindrical\_Coordinates}}
       \\
       &  \multicolumn{2}{l}{\inptext{Output\_Relative\_Velocity}}
       && \multicolumn{2}{l}{\inp[nmlOutputVariables::RelativeVelocity]{Relative\_Velocity}}
       \\
       &  \multicolumn{2}{l}{\inptext{Output\_Relative\_Conditions}}
       && \multicolumn{2}{l}{\inp[nmlOutputVariables::RelativeConditions]{Relative\_Conditions}}
       \\
       &  \multicolumn{2}{l}{\inptext{Output\_Thermo\_Variables}}
       && \multicolumn{2}{l}{\inp[nmlOutputVariables::ThermoVariables]{Thermo\_Variables}}
       \\
       &  \multicolumn{2}{l}{\inptext{Output\_Temperature}}
       && \multicolumn{2}{l}{\inp[nmlOutputVariables::Temperature]{Temperature}}
       \\
       &  \multicolumn{2}{l}{\inptext{Output\_Gradients}}
       && \multicolumn{2}{l}{\inp[nmlOutputVariables::Gradients]{Gradients}}
       \\
       &  \multicolumn{2}{l}{\inptext{Output\_Transport\_Variables}}
       && \multicolumn{2}{l}{\inp[nmlOutputVariables::TransportVariables]{Transport\_Variables}}
       \\
       &  \multicolumn{2}{l}{\inptext{Output\_Stagnation\_Variables}}
       && \multicolumn{2}{l}{\inp[nmlOutputVariables::StagnationVariables]{Stagnation\_Variables}}
       \\
       &  \multicolumn{2}{l}{\inptext{Output\_Stagnation\_Pressure}}
       && \multicolumn{2}{l}{\inp[nmlOutputVariables::StagnationPressure]{Stagnation\_Pressure}}
       \\
       &  \multicolumn{2}{l}{\inptext{Output\_Stagnation\_Temperature}}
       && \multicolumn{2}{l}{\inp[nmlOutputVariables::StagnationTemperature]{Stagnation\_Temperature}}
       \\
       &  \multicolumn{2}{l}{\inptext{Output\_Marker\_Variables}}
       && \multicolumn{2}{l}{\inp[nmlOutputVariables::MarkerVariables]{Marker\_Variables}}
       \\
       &  \multicolumn{2}{l}{\inptext{Output\_Qcriterion}}
       && \multicolumn{2}{l}{\inp[nmlOutputVariables::Qcriterion]{Qcriterion}}
       \\
       &  \multicolumn{2}{l}{\inptext{Output\_Species\_Mass\_Fractions}}
       && \multicolumn{2}{l}{\inp[nmlOutputVariables::SpeciesMassFraction]{Species\_Mass\_Fraction}}
       \\
       &  \multicolumn{2}{l}{\inptext{Output\_Species\_Density}}
       && \multicolumn{2}{l}{\inp[nmlOutputVariables::SpeciesDensity]{Species\_Density}}
       \\
       &  \multicolumn{2}{l}{\inptext{Output\_Species\_Mole\_Fractions}}
       && \multicolumn{2}{l}{\inp[nmlOutputVariables::SpeciesMoleFraction]{Species\_Mole\_Fraction}}
       \\
       &  \multicolumn{2}{l}{\inptext{Output\_Extra\_Species\_Variables}}
       && \multicolumn{2}{l}{\inp[nmlOutputVariables::ExtraSpeciesVariables]{Extra\_Species\_Variables}}
       \\
       &  \multicolumn{2}{l}{\inptext{Output\_Debug\_Species\_Variables}}
       && \multicolumn{2}{l}{\inp[nmlOutputVariables::DebugSpeciesVariables]{Debug\_Species\_Variables}}
       \\
       &  \multicolumn{2}{l}{\inptext{Output\_Grdnt\_Species\_Variables}}
       && \multicolumn{2}{l}{\inp[nmlOutputVariables::GrdntSpeciesVariables]{Grdnt\_Species\_Variables}}
       \\
%==========================================================================
       \multicolumn{3}{l}{\&input}
    &  \multicolumn{2}{l}{\&\texttt{nml\_sem}}
    \\* %%% Note: \\* prevents pagebreak.
       &  \multicolumn{2}{l}{\inptext{Use\_SEM\_Profile\_For\_Initialization}}
       && \multicolumn{2}{l}{\inp[nmlSEM::UseProfileForInitialization]{Use\_Profile\_For\_Initialization}}
       \\
       &  \multicolumn{2}{l}{\inptext{SEM\_Use\_Random\_Seed\_From\_Restart}}
       && \multicolumn{2}{l}{\inp[nmlSEM::UseRandomSeedFromRestart]{Use\_Random\_Seed\_From\_Restart}}
       \\
       &  \multicolumn{2}{l}{\inptext{SEM\_Eddy\_Volume\_Is\_Spherical}}
       && \multicolumn{2}{l}{\inp[nmlSEM::UseSphericalEddyVolume]{Use\_Spherical\_Eddy\_Volume}}
       \\
       &  \multicolumn{2}{l}{\inptext{SEM\_Include\_Thermo\_Fluctuations}}
       && \multicolumn{2}{l}{\inp[nmlSEM::IncludeThermoFluctuations]{Include\_Thermo\_Fluctuations}}
       \\
       &  \multicolumn{2}{l}{\inptext{SEM\_Box\_Num\_Eddies\_Multiplier}}
       && \multicolumn{2}{l}{\inp[nmlSEM::BoxNumEddiesMultiplier]{Box\_Num\_Eddies\_Multiplier}}
       \\
%==========================================================================
       \multicolumn{3}{l}{\&input}
    &  \multicolumn{2}{l}{\&\texttt{nml\_chemistry}}
    \\* %%% Note: \\* prevents pagebreak.
       &  \multicolumn{2}{l}{\inptext{Disable\_Chemical\_Reactions}}
       && \multicolumn{2}{l}{\inp[nmlChemistry::DisableChemicalReactions]{Same}[Disable\_Chemical\_Reactions]}
       \\
       &  \multicolumn{2}{l}{\inptext{Semi\_Implicit\_Chemistry}}
       && \multicolumn{2}{l}{\inp[nmlChemistry::SemiImplicitChemistry]{Same}[Semi\_Implicit\_Chemistry]}
       \\
       &  \multicolumn{2}{l}{\inptext{Ignition\_Method}}
       && \multicolumn{2}{l}{\inp[nmlChemistry::IgnitionMethod]{Same}[Ignition\_Method]}
       \\
       &  \multicolumn{2}{l}{\inptext{Ignition\_Start\_Timestep}}
       && \multicolumn{2}{l}{\inp[nmlChemistry::IgnitionStartTimestep]{Same}[Ignition\_Start\_Timestep]}
       \\
       &  \multicolumn{2}{l}{\inptext{Ignition\_Duration}}
       && \multicolumn{2}{l}{\inp[nmlChemistry::IgnitionDuration]{Same}[Ignition\_Duration]}
       \\
       &  \multicolumn{2}{l}{\inptext{Ignition\_Prod2Reac\_L2\_Ratio\_Goal}}
       && \multicolumn{2}{l}{\inp[nmlChemistry::IgnitionProd2ReacL2RatioGoal]{Same}[Ignition\_Prod2Reac\_L2\_Ratio\_Goal]}
       \\
       &  \multicolumn{2}{l}{\inptext{Ignition\_Max\_Temperature}}
       && \multicolumn{2}{l}{\inp[nmlChemistry::IgnitionMaxTemperature]{Same}[Ignition\_Max\_Temperature]}
       \\
       &  \multicolumn{2}{l}{\inptext{Ignition\_Spark\_Cap}}
       && \multicolumn{2}{l}{\inp[nmlChemistry::IgnitionSparkCap]{Same}[Ignition\_Spark\_Cap]}
       \\
       &  \multicolumn{2}{l}{\inptext{Ignition\_Source\_Magnitude}}
       && \multicolumn{2}{l}{\inp[nmlChemistry::IgnitionSourceMagnitude]{Same}[Ignition\_Source\_Magnitude]}
       \\
       &  \multicolumn{2}{l}{\inptext{Ignition\_Source\_Radius}}
       && \multicolumn{2}{l}{\inp[nmlChemistry::IgnitionSourceRadius]{Same}[Ignition\_Source\_Radius]}
       \\
       &  \multicolumn{2}{l}{\inptext{Ignition\_Source\_Location}}
       && \multicolumn{2}{l}{\inp[nmlChemistry::IgnitionSourceLocation]{Same}[Ignition\_Source\_Location]}
       \\
%==========================================================================
       \multicolumn{3}{l}{\&input}
    &  \multicolumn{2}{l}{\&\texttt{nml\_testcase\_channel}}
    \\
       &  \multicolumn{2}{l}{\inptext{Pipe\_Channel\_Angle}}
       && \multicolumn{2}{l}{\inp[nmlTestcaseChannel::GridAnnulusAngle]{Grid\_Annulus\_Angle}}
       \\
       &  \multicolumn{2}{l}{\inptext{Channel\_Grid\_Ave\_Xplus}}
       && \multicolumn{2}{l}{\inp[nmlTestcaseChannel::GridAveXplus]{Grid\_Ave\_Xplus}}
       \\
       &  \multicolumn{2}{l}{\inptext{Channel\_Grid\_Ave\_Zplus}}
       && \multicolumn{2}{l}{\inp[nmlTestcaseChannel::GridAveZplus]{Grid\_Ave\_Zplus}}
       \\
       &  \multicolumn{2}{l}{\inptext{Channel\_Grid\_Ave\_Yplus}}
       && \multicolumn{2}{l}{\inp[nmlTestcaseChannel::GridAveYplus]{Grid\_Ave\_Yplus}}
       \\
       &  \multicolumn{2}{l}{\inptext{Channel\_Grid\_Yplus\_Wall}}
       && \multicolumn{2}{l}{\inp[nmlTestcaseChannel::GridYplusWall]{Grid\_Yplus\_Wall}}
       \\
       &  \multicolumn{2}{l}{\inptext{Channel\_Yplus\_Of\_One}}
       && \multicolumn{2}{l}{\inp[nmlTestcaseChannel::GridYplusOfOne]{Grid\_Yplus\_Of\_One}}
       \\
       &  \multicolumn{2}{l}{\inptext{Channel\_Body\_Force}}
       && \multicolumn{2}{l}{\inp[nmlTestcaseChannel::BodyForce]{Body\_Force}}
       \\
       &  \multicolumn{2}{l}{\inptext{Channel\_Nondim\_Output\_Coord}}
       && \multicolumn{2}{l}{\inp[nmlTestcaseChannel::NondimOutputCoord]{Nondim\_Output\_Coord}}
       \\
%==========================================================================
       \multicolumn{3}{l}{\&input}
    &  \multicolumn{2}{l}{\&\texttt{nml\_testcase\_pulse}}
    \\* %%% Note: \\* prevents pagebreak.
       &  \multicolumn{2}{l}{\inptext{Pulse\_Test\_Form}}
       && \multicolumn{2}{l}{\inp[nmlTestcasePulse::Form]{Form}}
       \\
       &  \multicolumn{2}{l}{\inptext{Pulse\_Test\_Form\_Dim}}
       && \multicolumn{2}{l}{\inp[nmlTestcasePulse::FormDim]{Form\_Dim}}
       \\
       &  \multicolumn{2}{l}{\inptext{Pulse\_Test\_Length}}
       && \multicolumn{2}{l}{\inp[nmlTestcasePulse::Length]{Length}}
       \\
       &  \multicolumn{2}{l}{\inptext{Pulse\_Test\_Center}}
       && \multicolumn{2}{l}{\inp[nmlTestcasePulse::Center]{Center}}
       \\
       &  \multicolumn{2}{l}{\inptext{Pulse\_Test\_Width}}
       && \multicolumn{2}{l}{\inp[nmlTestcasePulse::Width]{Width}}
       \\
       &  \multicolumn{2}{l}{\inptext{Pulse\_Test\_Gap}}
       && \multicolumn{2}{l}{\inp[nmlTestcasePulse::Gap]{Gap}}
       \\
       &  \multicolumn{2}{l}{\inptext{Pulse\_Test\_Repeat}}
       && \multicolumn{2}{l}{\inp[nmlTestcasePulse::Repeat]{Repeat}}
       \\
%==========================================================================
       \multicolumn{3}{l}{\&input}
    &  \multicolumn{2}{l}{\&\texttt{nml\_testcase\_sod}}
    \\* %%% Note: \\* prevents pagebreak.
       &  \multicolumn{2}{l}{\inptext{Sod\_Test\_Right\_Vector}}
       && \multicolumn{2}{l}{\inp[nmlTestcaseSod::RightVector]{Right\_Vector}}
       \\
%==========================================================================
       \multicolumn{3}{l}{\&input}
    &  \multicolumn{2}{l}{\&\texttt{nml\_testcase\_tcf}}
    \\* %%% Note: \\* prevents pagebreak.
       &  \multicolumn{2}{l}{\inptext{Tcf\_Inner\_Radius}}
       && \multicolumn{2}{l}{\inp[nmlTestcaseTCF::InnerRadius]{Inner\_Radius}}
       \\
       &  \multicolumn{2}{l}{\inptext{Tcf\_Inner\_Wall\_Spacing}}
       && \multicolumn{2}{l}{\inp[nmlTestcaseTCF::InnerWallSpacing]{Inner\_Wall\_Spacing}}
       \\
       &  \multicolumn{2}{l}{\inptext{Tcf\_Outer\_Radius}}
       && \multicolumn{2}{l}{\inp[nmlTestcaseTCF::OuterRadius]{Outer\_Radius}}
       \\
       &  \multicolumn{2}{l}{\inptext{Tcf\_Outer\_Wall\_Spacing}}
       && \multicolumn{2}{l}{\inp[nmlTestcaseTCF::OuterWallSpacing]{Outer\_Wall\_Spacing}}
       \\
       &  \multicolumn{2}{l}{\inptext{Tcf\_Aspect\_Ratio}}
       && \multicolumn{2}{l}{\inp[nmlTestcaseTCF::AspectRatio]{Aspect\_Ratio}}
       \\
       &  \multicolumn{2}{l}{\inptext{Tcf\_Axial\_Cells}}
       && \multicolumn{2}{l}{\inp[nmlTestcaseTCF::AxialCells]{Axial\_Cells}}
       \\
       &  \multicolumn{2}{l}{\inptext{Tcf\_Radial\_Cells}}
       && \multicolumn{2}{l}{\inp[nmlTestcaseTCF::RadialCells]{Radial\_Cells}}
       \\
       &  \multicolumn{2}{l}{\inptext{Tcf\_Azimuthal\_Cells}}
       && \multicolumn{2}{l}{\inp[nmlTestcaseTCF::AzimuthalCells]{Azimuthal\_Cells}}
       \\
       &  \multicolumn{2}{l}{\inptext{Tcf\_Theta\_Shift}}
       && \multicolumn{2}{l}{\inp[nmlTestcaseTCF::ThetaShift]{Theta\_Shift}}
       \\
       &  \multicolumn{2}{l}{\inptext{Tcf\_Periodic\_Sectors}}
       && \multicolumn{2}{l}{\inp[nmlTestcaseTCF::PeriodicSectors]{Periodic\_Sectors}}
       \\
%==========================================================================
       \multicolumn{3}{l}{\&input}
    &  \multicolumn{2}{l}{\&\texttt{nml\_testcase\_vortex}}
    \\* %%% Note: \\* prevents pagebreak.
       &  \multicolumn{2}{l}{\inptext{Vortex\_Bigr}}
       && \multicolumn{2}{l}{\inp[nmlTestcaseVortex::VortexBigr]{Same}[Vortex\_Bigr]}
       \\
       &  \multicolumn{2}{l}{\inptext{Vortex\_Biggam}}
       && \multicolumn{2}{l}{\inp[nmlTestcaseVortex::VortexBiggam]{Same}[Vortex\_Biggam]}
       \\
       &  \multicolumn{2}{l}{\inptext{VortexGrid\_Lref}}
       && \multicolumn{2}{l}{\inp[nmlTestcaseVortex::GridLref]{Grid\_Lref}}
       \\
       &  \multicolumn{2}{l}{\inptext{VortexGrid\_Random\_Perturb}}
       && \multicolumn{2}{l}{\inp[nmlTestcaseVortex::GridRandomPerturb]{Grid\_Random\_Perturb}}
       \\
       &  \multicolumn{2}{l}{\inptext{VortexGrid\_Perturb\_Percent}}
       && \multicolumn{2}{l}{\inp[nmlTestcaseVortex::GridPerturbPercent]{Grid\_Perturb\_Percent}}
       \\
       &  \multicolumn{2}{l}{\inptext{VortexGrid\_Hex\_To\_Pyramids}}
       && \multicolumn{2}{l}{\inp[nmlTestcaseVortex::GridHexToPyramids]{Grid\_Hex\_To\_Pyramids}}
       \\
       &  \multicolumn{2}{l}{\inptext{VortexGrid\_Is\_3d}}
       && \multicolumn{2}{l}{\inp[nmlTestcaseVortex::GridIs3d]{Grid\_Is\_3d}}
       \\
%==========================================================================
       \multicolumn{3}{l}{\&input}
    &  \multicolumn{2}{l}{\&\texttt{nml\_interpolation}}
    \\* %%% Note: \\* prevents pagebreak.
       &  \multicolumn{2}{l}{\inptext{Interpto\_Grid\_Scaling\_Factor}}
       && \multicolumn{2}{l}{\inp[nmlInterpolation::InterptoGridScalingFactor]{Same}[Interpto\_Grid\_Scaling\_Factor]}
       \\
       &  \multicolumn{2}{l}{\inptext{Interpto\_Grid\_Format}}
       && \multicolumn{2}{l}{\inp[nmlInterpolation::InterptoGridFormat]{Same}[Interpto\_Grid\_Format]}
       \\
       &  \multicolumn{2}{l}{\inptext{Interpto\_Gridfile}}
       && \multicolumn{2}{l}{\inp[nmlInterpolation::InterptoGridFile]{Interpto\_Grid\_File}}
       \\
       &  \multicolumn{2}{l}{\inptext{Interpto\_Plot3d\_Cutfile}}
       && \multicolumn{2}{l}{\inp[nmlInterpolation::InterptoPlot3dCutfile]{Same}[Interpto\_Plot3d\_Cutfile]}
       \\
       &  \multicolumn{2}{l}{\inptext{Interpto\_Plot3d\_BCSfile}}
       && \multicolumn{2}{l}{\inp[nmlInterpolation::InterptoPlot3dBCSfile]{Same}[Interpto\_Plot3d\_BCSfile]}
       \\
       &  \multicolumn{2}{l}{\inptext{Interpto\_Plot3d\_I\_Stride}}
       && \multicolumn{2}{l}{\inp[nmlInterpolation::InterptoPlot3dStrideI]{Interpto\_Plot3d\_Stride\_I}}
       \\
       &  \multicolumn{2}{l}{\inptext{Interpto\_Plot3d\_J\_Stride}}
       && \multicolumn{2}{l}{\inp[nmlInterpolation::InterptoPlot3dStrideJ]{Interpto\_Plot3d\_Stride\_J}}
       \\
       &  \multicolumn{2}{l}{\inptext{Interpto\_Plot3d\_K\_Stride}}
       && \multicolumn{2}{l}{\inp[nmlInterpolation::InterptoPlot3dStrideK]{Interpto\_Plot3d\_Stride\_K}}
       \\
       &  \multicolumn{2}{l}{\inptext{Interpto\_Plot3d\_All\_Strides}}
       && \multicolumn{2}{l}{\inp[nmlInterpolation::InterptoPlot3dStrideAll]{Interpto\_Plot3d\_Stride\_All}}
       \\
       &  \multicolumn{2}{l}{\inptext{Interpto\_Plot3d\_Agglomerate\_Order}}
       && \multicolumn{2}{l}{\inp[nmlInterpolation::InterptoPlot3dAgglomerateOrder]{Same}[Interpto\_Plot3d\_Agglomerate\_Order]}
       \\
       &  \multicolumn{2}{l}{\inptext{Interpto\_Solution\_Order}}
       && \multicolumn{2}{l}{\inp[nmlInterpolation::InterptoSolutionOrder]{Same}[Interpto\_Solution\_Order]}
       \\
%==========================================================================
       \multicolumn{3}{l}{\&input}
    &  \multicolumn{2}{l}{\&\texttt{nml\_utility}}
    \\* %%% Note: \\* prevents pagebreak.
       &  \multicolumn{2}{l}{\inptext{Convert\_Restart\_To\_CGNS}}
       && \multicolumn{2}{l}{\inp[nmlUtility::ConvertRestartToCGNS]{Same}[Convert\_Restart\_To\_CGNS]}
       \\
       &  \multicolumn{2}{l}{\inptext{Delete\_Converted\_Restart\_Files}}
       && \multicolumn{2}{l}{\inp[nmlUtility::DeleteConvertedRestartFiles]{Same}[Delete\_Converted\_Restart\_Files]}
       \\
       &  \multicolumn{2}{l}{\inptext{Restart\_List\_File}}
       && \multicolumn{2}{l}{\inp[nmlUtility::RestartListFile]{Same}[Restart\_List\_File]}
       \\
       &  \multicolumn{2}{l}{\inptext{Channel\_Spatial\_Ave\_Restart}}
       && \multicolumn{2}{l}{\inp[nmlUtility::SpatialAveRestartChannel]{Spatial\_Ave\_Restart\_Channel}}
       \\
       &  \multicolumn{2}{l}{\inptext{Spatial\_Ave\_Restart\_X\_And\_Z\_Planes}}
       && \multicolumn{2}{l}{\inp[nmlUtility::SpatialAveRestartXAndZPlanes]{Same}[Spatial\_Ave\_Restart\_X\_And\_Z\_Planes]}
       \\
       &  \multicolumn{2}{l}{\inptext{Spatial\_Ave\_Restart\_Z\_Planes}}
       && \multicolumn{2}{l}{\inp[nmlUtility::SpatialAveRestartZPlanes]{Same}[Spatial\_Ave\_Restart\_Z\_Planes]}
       \\
%==========================================================================
       \multicolumn{3}{l}{\&input}
    &  \multicolumn{2}{l}{\&\texttt{nml\_debug}}
    \\* %%% Note: \\* prevents pagebreak.
       &  \multicolumn{2}{l}{\inptext{Output\_Limiter\_Statistics}}
       && \multicolumn{2}{l}{\inp[nmlDebug::OutputLimiterStatistics]{Same}[Output\_Limiter\_Statistics]}
       \\
       &  \multicolumn{2}{l}{\inptext{Output\_Filter\_Diagonals}}
       && \multicolumn{2}{l}{\inp[nmlDebug::OutputFilterDiagonals]{Same}[Output\_Filter\_Diagonals]}
       \\
       &  \multicolumn{2}{l}{\inptext{Output\_Plot3d\_To\_Tecplot}}
       && \multicolumn{2}{l}{\inp[nmlDebug::OutputPlot3dToTecplot]{Same}[Output\_Plot3d\_To\_Tecplot]}
       \\
       &  \multicolumn{2}{l}{\inptext{Output\_Bface\_Array}}
       && \multicolumn{2}{l}{\inp[nmlDebug::OutputBfaceArray]{Same}[Output\_Bface\_Array]}
       \\
       &  \multicolumn{2}{l}{\inptext{Dump\_Memory\_Usage}}
       && \multicolumn{2}{l}{\inp[nmlDebug::DumpMemoryUsage]{Same}[Dump\_Memory\_Usage]}
       \\
       &  \multicolumn{2}{l}{\inptext{Check\_Flux\_Point\_Coordinates}}
       && \multicolumn{2}{l}{\inp[nmlDebug::CheckFluxPointCoordinates]{Same}[Check\_Flux\_Point\_Coordinates]}
       \\
       &  \multicolumn{2}{l}{\inptext{Force\_Nd\_Subroutines}}
       && \multicolumn{2}{l}{\inp[nmlDebug::ForceNdSubroutines]{Same}[Force\_Nd\_Subroutines]}
       \\
       &  \multicolumn{2}{l}{\inptext{Conv\_Temp}}
       && \multicolumn{2}{l}{\inp[nmlDebug::NewtonConvergenceGoal]{Newton\_Convergence\_Goal}}
       \\
       &  \multicolumn{2}{l}{\inptext{Chemistry\_Sub\_Iterations}}
       && \multicolumn{2}{l}{\inp[nmlDebug::ChemistrySubIterations]{Same}[Chemistry\_Sub\_Iterations]}
       \\
       &  \multicolumn{2}{l}{\inptext{Traceback\_Unconverged\_Newton\_Temps}}
       && \multicolumn{2}{l}{\inp[nmlDebug::TracebackUnconvergedNewtonTemps]{Same}[Traceback\_Unconverged\_Newton\_Temps]}
       \\
       &  \multicolumn{2}{l}{\inptext{Ignore\_Thermo\_Newton\_Warnings}}
       && \multicolumn{2}{l}{\inp[nmlDebug::IgnoreThermoNewtonWarnings]{Same}[Ignore\_Thermo\_Newton\_Warnings]}
       \\[0.5em]
       \multicolumn{5}{l}{
       The variables below have been removed from the input namelist.
       }
       \\
%==========================================================================
       \multicolumn{3}{l}{\&input}
    &  \multicolumn{2}{l}{\&\texttt{nml\_removed}}
    \\* %%% Note: \\* prevents pagebreak.
       &  \multicolumn{2}{l}{\inptext{Parametric\_Agglomeration}}
       && \multicolumn{2}{l}{Same}
       \\
       &  \multicolumn{2}{l}{\inptext{Interpto\_All\_Orders}}
       && \multicolumn{2}{l}{Same}
       \\
       &  \multicolumn{2}{l}{\inptext{Interpto\_Projection\_Order}}
       && \multicolumn{2}{l}{Same}
       \\
       &  \multicolumn{2}{l}{\inptext{Interpto\_Quadrature\_Order}}
       && \multicolumn{2}{l}{Same}
       \\
       &  \multicolumn{2}{l}{\inptext{Interpto\_Output\_Order}}
       && \multicolumn{2}{l}{Same}
       \\
       &  \multicolumn{2}{l}{\inptext{Interpto\_Error\_Order}}
       && \multicolumn{2}{l}{Same}
       \\
       &  \multicolumn{2}{l}{\inptext{Prt}}
       && \multicolumn{2}{l}{\texttt{Prandtl\_Turb}}
       \\
       &  \multicolumn{2}{l}{\inptext{Schmidt\_Turb}}
       && \multicolumn{2}{l}{Same}
       \\
       &  \multicolumn{2}{l}{\inptext{Tinref}}
       && \multicolumn{2}{l}{\texttt{Turb\_Intensity}}
       \\
       &  \multicolumn{2}{l}{\inptext{Tvrref}}
       && \multicolumn{2}{l}{\texttt{Turb\_Viscosity}}
       \\
       &  \multicolumn{2}{l}{\inptext{Turbulence\_Model}}
       && \multicolumn{2}{l}{Same}
       \\
       &  \multicolumn{2}{l}{\inptext{Newton\_Rlx}}
       && \multicolumn{2}{l}{\texttt{Newton\_Relax}}
       \\
       &  \multicolumn{2}{l}{\inptext{Use\_Improved\_SEM\_Normalization}}
       && \multicolumn{2}{l}{\texttt{Use\_Improved\_Initialization}}
       \\
       &  \multicolumn{2}{l}{\inptext{Average\_SEM\_Inflow\_Plane}}
       && \multicolumn{2}{l}{\texttt{Average\_Inflow\_Plane}}
       \\
       &  \multicolumn{2}{l}{\inptext{MMS\_Opt}}
       && \multicolumn{2}{l}{\texttt{Option}}
       \\
       &  \multicolumn{2}{l}{\inptext{MMS\_Init}}
       && \multicolumn{2}{l}{\texttt{Initialization}}
       \\
       &  \multicolumn{2}{l}{\inptext{MMS\_Output\_Solution}}
       && \multicolumn{2}{l}{\texttt{Output\_Solution}}
       \\
       &  \multicolumn{2}{l}{\inptext{Loc\_Init\_Pts}}
       && \multicolumn{2}{l}{Same}
       \\
       &  \multicolumn{2}{l}{\inptext{Use\_Old\_Restart\_Format}}
       && \multicolumn{2}{l}{Same}
       \\
       &  \multicolumn{2}{l}{\inptext{CFL\_Reduction\_For\_Nan\_Restart}}
       && \multicolumn{2}{l}{Same}
       \\
       &  \multicolumn{2}{l}{\inptext{Use\_Unrolled\_Dot\_Products}}
       && \multicolumn{2}{l}{Same}
       \\
       &  \multicolumn{2}{l}{\inptext{Flux\_Divergence}}
       && \multicolumn{2}{l}{Same}
       \\
       &  \multicolumn{2}{l}{\inptext{Pyramid\_Filter\_Option}}
       && \multicolumn{2}{l}{Same}
       \\
       &  \multicolumn{2}{l}{\inptext{Ideal\_Total\_Temperature}}
       && \multicolumn{2}{l}{Same}
       \\
       &  \multicolumn{2}{l}{\inptext{Limit\_Total\_Temperature}}
       && \multicolumn{2}{l}{Same}
       \\
       &  \multicolumn{2}{l}{\inptext{Filter\_Eps}}
       && \multicolumn{2}{l}{\texttt{Global\_Filter\_Eps}}
       \\
       &  \multicolumn{2}{l}{\inptext{Limiter\_Filter\_Eps}}
       && \multicolumn{2}{l}{\texttt{Local\_Filter\_Eps}}
       \\
       &  \multicolumn{2}{l}{\inptext{Limiter\_Marker\_C1}}
       && \multicolumn{2}{l}{\texttt{Marker\_C1}}
       \\
       &  \multicolumn{2}{l}{\inptext{Limiter\_Marker\_C2}}
       && \multicolumn{2}{l}{\texttt{Marker\_C2}}
       \\
       &  \multicolumn{2}{l}{\inptext{Output\_Eperk\_CFL}}
       && \multicolumn{2}{l}{Same}
       \\
       &  \multicolumn{2}{l}{\inptext{Disable\_Header\_Rewrites}}
       && \multicolumn{2}{l}{Same}
       \\
       &  \multicolumn{2}{l}{\inptext{CGNS\_Use\_Queue}}
       && \multicolumn{2}{l}{Same}
       \\
       &  \multicolumn{2}{l}{\inptext{Use\_Async\_For\_Error}}
       && \multicolumn{2}{l}{Same}
       \\
       &  \multicolumn{2}{l}{\inptext{Output\_Bnd\_Profiles}}
       && \multicolumn{2}{l}{Same}
       \\
       &  \multicolumn{2}{l}{\inptext{Channel\_Init\_File}}
       && \multicolumn{2}{l}{\texttt{Init\_File}}
       \\
       &  \multicolumn{2}{l}{\inptext{Time\_Average\_Restart\_Files}}
       && \multicolumn{2}{l}{Same}
       \\
       &  \multicolumn{2}{l}{\inptext{Postproc\_Debug\_Grid}}
       && \multicolumn{2}{l}{Same}
       \\
       &  \multicolumn{2}{l}{\inptext{Dump\_Fluxes}}
       && \multicolumn{2}{l}{Same}
       \\
       &  \multicolumn{2}{l}{\inptext{Channel\_Rake\_Save\_Interval}}
       && \multicolumn{2}{l}{\texttt{Rake\_Save\_Interval}}
       \\
       &  \multicolumn{2}{l}{\inptext{Track\_Limiter\_Filter\_Dissipation}}
       && \multicolumn{2}{l}{\texttt{Local\_Filter\_Dissipation}}
       \\
       &  \multicolumn{2}{l}{\inptext{Continuous\_Output}}
       && \multicolumn{2}{l}{Same}
       \\
       &  \multicolumn{2}{l}{\inptext{Interpolate\_Before\_Output\_Variable}}
       && \multicolumn{2}{l}{Same}
       \\
       &  \multicolumn{2}{l}{\inptext{Num\_Rk\_Stages}}
       && \multicolumn{2}{l}{\texttt{Runge\_Kutta\_Stages}}
       \\
       &  \multicolumn{2}{l}{\inptext{Geom\_Wts(1)\%Edge}}
       && \multicolumn{2}{l}{Same}
       \\
       &  \multicolumn{2}{l}{\inptext{Geom\_Wts(1)\%Face\_Tria}}
       && \multicolumn{2}{l}{Same}
       \\
       &  \multicolumn{2}{l}{\inptext{Geom\_Wts(1)\%Face\_Quad}}
       && \multicolumn{2}{l}{Same}
       \\
      \hline
    \end{longtable}
    }%end font size
  \end{center}
%\end{table}

\RemoveFromHook{cmd/subsection/before}
